% appendix_body.tex --- IEEEtran appendix body, converted from paper_draft/APPENDICES.md.
% The wrapper (main.tex) issues \appendices before \input, so each \section below is
% lettered A, B, ... automatically. Do NOT add \appendices here. Order A--F is preserved.

Appendix~A supports Sec.~4 (Theorem~A and Corollaries 1--2); Appendix~B supports Sec.~4 (Theorem~A$'$, parametric-transversality proof); Appendix~C supports Sec.~5 (Propositions 1--2, condition ($\bigstar$), Theorem~B); Appendix~D supports Sec.~6 (Theorems C--D, Fisher geometry, low-photon); Appendix~E supports Sec.~7 (SRHT Walsh condition, permutation bounds); Appendix~F supports Sec.~8 (Theorem~1, specializations, Prop~3). Heuristic or genericity steps are labeled as such inside each proof; every appendix ends its theorems with a ``What this does not claim'' paragraph.

\section{Square-design non-identifiability}

This appendix proves Theorem~A of Sec.~4 together with its two corollaries, in the notation of the main text: noiseless bucket record $R_n = a_n B_n$, carrier $B_n = (MT)_n$, log-gain $\ell = \log a$ (entrywise), drift class $S$ a linear subspace of $\mathbb{R}^N$ with $\operatorname{span}\{\mathbf 1\}\subseteq S$ and $\dim S = p$. This appendix is the archival form of these results under the manuscript labels Theorem~A, Corollary~1, Corollary~2.

\subsection{A.0 Standing assumptions}

\begin{itemize}
\item \textbf{(A1) Model.} Noiseless record $R_n = a_n(MT)_n$, $n = 1,\dots,N$, with $N = K$ and $M \in \mathbb{R}^{K\times K}$ \textbf{invertible} --- Hadamard, Fourier, random square, or otherwise. Write $c = MT$ for the carrier coefficient vector, so $R = e^{\ell}\odot c$ ($\odot$ = entrywise product).
\item \textbf{(A2) Gain class.} $a_n > 0$ for all $n$; the admissible log-gains form (or contain a neighborhood inside) the linear subspace $S \ni \mathbf 1$, $\dim S = p \ge 2$; write $p - 1 = \dim(S/\operatorname{span}\{\mathbf 1\})$ for the nonconstant drift dimension.
\item \textbf{(A3) Object class.} $T \in \mathbb{R}^K$ unconstrained (Theorem~A, Corollary~1); constrained classes are treated in Corollary~2.
\item \textbf{(A4) Carrier support.} For the clean dimension statements we assume full carrier support, $c_n = (MT)_n \neq 0$ for all $n$; this holds for generic $T$ (each $\{c_n = 0\}$ is a hyperplane). Zero buckets are discussed in Remark~A.2.
\item \textbf{(A5) Gauge.} The global-scale gauge orbit of a pair $(\ell, T)$ is $\mathcal G(\ell,T) = \{(\ell + t\mathbf 1,\, e^{-t}T) : t \in \mathbb{R}\}$; identifiability is always meant modulo $\mathcal G$.
\end{itemize}

\subsection{A.1 Theorem A (square non-identifiability, any invertible design)}

\begin{namedthm}{Theorem A}
Assume (A1)--(A5). For every $s \in S$ the pair
\[
\ell' = \ell + s, \qquad c' = e^{-s}\odot c, \qquad T'(s) = M^{-1}\!\left(e^{-s}\odot MT\right)
\]
reproduces the bucket record exactly: $e^{\ell'_n} c'_n = R_n$ for all $n$. If $s$ is nonconstant on $\operatorname{supp}(MT)$ --- under (A4), simply nonconstant --- then $(\ell', T'(s)) \notin \mathcal G(\ell, T)$ and $T'(s) \neq \gamma T$ for every $\gamma \in \mathbb{R}$. The set of admissible pairs consistent with the record therefore contains a real-analytic family of dimension $\dim S = p$, of which the gauge orbit accounts for only $\dim(S \cap \operatorname{span}\{\mathbf 1\}) = 1$; the non-gauge ambiguity dimension is at least $p - 1 \ge 1$. In particular, gain--object separation from a square design is not identifiable modulo gauge whenever $S$ contains a nonconstant element.
\end{namedthm}

\noindent\emph{Proof.} \textbf{(Data equality.)} $e^{\ell'_n}c'_n = e^{\ell_n + s_n}\, e^{-s_n} c_n = e^{\ell_n} c_n = R_n$, and $T'(s) = M^{-1}c'$ is well defined because $M$ is invertible. Both steps are exact; no smallness of $s$ is used.

\textbf{(Distinctness of the object.)} $T'(s) = \gamma T$ for some $\gamma$ iff $MT'(s) = \gamma MT$, i.e. $e^{-s_n} c_n = \gamma c_n$ for every $n$, i.e. $e^{-s_n} = \gamma$ for every $n \in \operatorname{supp}(c)$. Since $x \mapsto e^{-x}$ is injective, this forces $s$ constant on $\operatorname{supp}(c)$. Contrapositively, $s$ nonconstant on $\operatorname{supp}(MT)$ implies $T'(s)$ is not proportional to $T$.

\textbf{(Distinctness modulo gauge.)} $(\ell + s, T'(s)) \in \mathcal G(\ell,T)$ requires $\ell + s = \ell + t\mathbf 1$ in the gain coordinate, i.e. $s = t\mathbf 1$ (log is injective on positive gains). A nonconstant $s$ therefore leaves the orbit; note that when $s = t\mathbf 1$ the construction reduces exactly to the gauge, $T'(t\mathbf 1) = e^{-t}T$, so the family contains the orbit and nothing of the orbit is double-counted.

\textbf{(Dimension.)} The map $s \mapsto (\ell + s, T'(s))$ is injective (read off the first coordinate) and real-analytic in $s$, hence its image is a $p$-dimensional analytically parametrized family of exact solutions. Intersecting with the gauge orbit uses up $\dim(S \cap \operatorname{span}\{\mathbf 1\}) = 1$ dimension, leaving non-gauge ambiguity dimension $\ge p - 1$. \hfill$\blacksquare$

The display in the main text (Sec.~4) is the instance $s = \bar\ell\,\mathbf 1 - \ell$: the drift is absorbed \emph{in one step} into the relabelled object $T' = M^{-1}\operatorname{diag}(a)MT$ (up to the gauge constant), leaving a constant gain. The hypothesis ``$a$ nonconstant on the support of $MT$'' in the main-text statement is exactly the distinctness condition established above.

\paragraph{What this does not claim.} Theorem~A is a statement about the square case $N = K$ and an unconstrained object; for $N > K$ the construction fails unless $e^{-s}\odot MT$ remains in $\operatorname{range}(M)$, and the correct analysis is the tall-design identifiability theory of Theorem~A$'$ (see Appendix~B). It does not say the record is uninformative about the gain in a statistical model: under random illumination, the \emph{relative} gain is consistently estimable (Theorem~B; see Appendix~C), a different --- statistical, asymptotic --- notion of identifiability which Theorem~A neither contradicts nor implies. Finally, Theorem~A is about uniqueness, not conditioning: it makes no claim about how well- or ill-posed any estimator is.

\subsection{A.2 Corollary 1 (slowness does not help)}

\begin{namedthm}{Corollary 1}
Assume (A1)--(A5).

(i) If the admissible log-gain class is the linear space $S$ itself, the entire ambiguity family of Theorem~A lies inside the class: $\ell \in S,\ s \in S \Rightarrow \ell + s \in S$. In particular, choosing $s = \bar\ell\,\mathbf 1 - \ell \in S$ produces a \textbf{constant-gain} exact explanation of the record with the static object $T' = M^{-1}(e^{\ell - \bar\ell}\odot MT)$: a slowly drifting gain over the true object is indistinguishable from no drift at all over a relabelled object.

(ii) If the admissible class is a slowness ball $\mathcal B_L = \{\ell \in S : \lVert \ell \rVert_* \le L\}$ for any seminorm $\lVert\cdot\rVert_*$ vanishing on constants (e.g. a H\"older or low-pass-energy seminorm), and $\ell$ is interior ($\lVert\ell\rVert_* < L$), then every nonconstant $s \in S$ with $\lVert s\rVert_*$ small enough keeps $\ell + s \in \mathcal B_L$. The local non-gauge ambiguity dimension remains $p - 1$: slowness bounds the \textbf{size} of the ambiguity, never its \textbf{dimension}.
\end{namedthm}

\noindent\emph{Proof.} (i) is linearity of $S$ plus Theorem~A applied with the stated $s$; nonconstancy of $\ell$ on $\operatorname{supp}(c)$ makes the constant-gain alternative distinct modulo gauge. (ii) Subadditivity of the seminorm gives $\lVert \ell + s\rVert_* \le \lVert\ell\rVert_* + \lVert s\rVert_* \le L$ once $\lVert s\rVert_* \le L - \lVert\ell\rVert_*$; the set of such $s$ is a neighborhood of $0$ in $S$ and Theorem~A applies to each nonconstant member. \hfill$\blacksquare$

\paragraph{What this does not claim.} Slowness is not useless in general --- in the tall-design regime the whole point is that slow drift means small $p$, making the threshold $N \ge K + p$ cheap to satisfy (Sec.~4). The corollary only asserts that at $N = K$ a slowness prior \emph{by itself} removes nothing from the ambiguity, because the ambiguity directions live in the slow class itself.

\subsection{A.3 Corollary 2 (nonnegativity is insufficient; known support zeros can restore local identifiability)}

Throughout this corollary write $H$ for the square invertible design ($H = M$; the letter matches the fixed-basis usage, e.g. Hadamard), and define the gain-absorption operator
\[
M_s := H^{-1}\operatorname{diag}(e^{-s})H, \qquad T'(s) = M_s T,
\]
so that the ambiguity set at $T$ under an object prior class $\mathcal T$ is $V_{\mathcal T}(T) = \{s \in S : M_sT \in \mathcal T\}$, and identifiability modulo gauge is the assertion $V_{\mathcal T}(T) = S \cap \operatorname{span}\{\mathbf 1\}$.

\begin{namedthm}{Corollary 2}
Assume (A1)--(A5).

(i) (Positivity insufficient.) Let $\mathcal T_+ = \{T : T_j \ge 0\ \forall j\}$ and let $T$ have full support, $T_j > 0$ for all $j$. Then there is $r > 0$ such that every $s \in S$ with $\lVert s\rVert_\infty < r$ satisfies $M_sT > 0$; explicitly, it suffices that
\[
\lVert H^{-1}\rVert_{2}\,\lVert HT\rVert_2\,(e^{r} - 1) < \min_j T_j .
\]
Hence $V_{\mathcal T_+}(T)$ contains a full neighborhood of $0$ in $S$ and the local non-gauge ambiguity dimension is still $p - 1$.

(ii) (Known support zeros; local criterion.) Let $\Omega \subset \{1,\dots,K\}$, $|\Omega| = q$, and $\mathcal T_\Omega = \{T : T_{\Omega^c} = 0,\ T \ge 0\}$, with the true $T \in \mathcal T_\Omega$, $T_\Omega > 0$. Then local identifiability modulo gauge holds at $T$ if and only if (at the first-order level)
\[
\ker\!\big(P_{\Omega^c} H^{-1}\operatorname{diag}(HT)\,\big|_S\big) \;=\; S \cap \operatorname{span}\{\mathbf 1\},
\]
where $P_{\Omega^c}$ is coordinate projection onto $\Omega^c$. The gauge direction always lies in the kernel; when the kernel is exactly the gauge, no ambiguity exists in a quantifiable ball: with $\gamma$ the smallest singular value of the linearized map on a complement $S_\perp$ of $\operatorname{span}\{\mathbf 1\}$ in $S$ and $C_H$ a second-order remainder constant (below), every $s$ with $0 < \lVert s_\perp\rVert_2 < \gamma / C_H$ violates the support constraint.

(iii) (Generic counts --- dimension-counting arguments, labeled as such.) The linearized map in (ii) sends the $(p-1)$-dimensional space $S/\operatorname{span}\{\mathbf 1\}$ into $\mathbb{R}^{K-q}$; injectivity is possible only if $K_0 := K - q \ge p - 1$, i.e. \textbf{at least one known zero per nonconstant drift dimension}. For generic $(H, T_\Omega)$ this necessary count is expected to be sufficient. For unknown support of size $\le q$, dimension counting predicts the threshold $K \ge 2q + (p-1)$, up to endpoint conventions.
\end{namedthm}

\noindent\emph{Proof.} \textbf{(i)} $M_0T = T$ and $\lVert M_sT - T\rVert_2 = \lVert H^{-1}\big((e^{-s} - \mathbf 1)\odot HT\big)\rVert_2 \le \lVert H^{-1}\rVert_2\,\lVert HT\rVert_2\,\max_n|e^{-s_n} - 1| \le \lVert H^{-1}\rVert_2\,\lVert HT\rVert_2\,(e^{\lVert s\rVert_\infty} - 1)$. If this is below $\min_j T_j$ then $M_sT > 0$ entrywise. Since the positive orthant is open and $s \mapsto M_sT$ is continuous, positivity excludes nothing locally; Theorem~A supplies the distinctness. \hfill$\blacksquare$\ \textnormal{(i)}

\textbf{(ii)} Define the support-constraint map $F : S \to \mathbb{R}^{K-q}$, $F(s) = P_{\Omega^c} M_s T$. It is real-analytic, $F(0) = P_{\Omega^c}T = 0$, and its differential at $0$ is
\[
DF_0\,u = -\,P_{\Omega^c} H^{-1}\operatorname{diag}(HT)\,u,
\]
by differentiating $e^{-tu} = \mathbf 1 - tu + O(t^2)$ inside $M_{tu}$. The gauge is always in the kernel: $H^{-1}\operatorname{diag}(HT)\mathbf 1 = H^{-1}(HT) = T$ and $P_{\Omega^c}T = 0$; consistently, $F(t\mathbf 1) = e^{-t}F(0) = 0$ exactly, for all $t$.

\emph{Necessity of the criterion (first-order sense).} If the kernel strictly exceeds the gauge, some nonconstant $u \in S$ satisfies the linearized support constraint. This is a first-order ambiguity direction: it makes $T$ a critical point of the constrained model map and hence a Fisher-information null direction for the corresponding functional (cf. Remark~A.1), so local \emph{differential} identifiability fails. Whether the direction integrates to an exact ambiguity curve requires a further constant-rank or analyticity argument and can fail at second order; we claim only the first-order (differential) statement, which is what the criterion tests.

\emph{Sufficiency in a ball.} Suppose the kernel equals the gauge and let $\gamma = \sigma_{\min}(DF_0|_{S_\perp}) > 0$. Entrywise Taylor with remainder, $|e^{-x} - 1 + x| \le \tfrac{x^2}{2}e^{|x|}$, gives for $\lVert s\rVert_\infty \le 1$
\[
\begin{gathered}
\lVert F(s) - DF_0\,s\rVert_2 \;\le\; C_H\,\lVert s\rVert_2^2,\\
C_H := \tfrac{e}{2}\,\lVert H^{-1}\rVert_2\,\lVert HT\rVert_2 .
\end{gathered}
\]
For gauge-fixed $s = s_\perp \in S_\perp$ (any $s$ may be reduced to this form using $F(s + t\mathbf 1) = e^{-t}F(s)$, which rescales but never annihilates $F$),
\[
\begin{gathered}
\lVert F(s_\perp)\rVert_2 \;\ge\; \gamma\,\lVert s_\perp\rVert_2 - C_H\,\lVert s_\perp\rVert_2^2 \;>\; 0\\
\text{for } 0 < \lVert s_\perp\rVert_2 < \gamma / C_H .
\end{gathered}
\]
A nonzero $F(s)$ means $M_sT \notin \mathcal T_\Omega$ (support constraint violated), so no non-gauge ambiguity exists in that ball. This is the quantitative inverse-function-theorem argument in explicit constants. \hfill$\blacksquare$\ \textnormal{(ii)}

\textbf{(iii)} \emph{These are genericity/dimension-count arguments, not proofs, and we state them as such.} The necessity direction is rigorous: a linear map from a $(p-1)$-dimensional space cannot be injective into $\mathbb{R}^{K-q}$ if $K - q < p - 1$, so at least $p - 1$ known zeros are required. The sufficiency direction is generic: the entries of $P_{\Omega^c}H^{-1}\operatorname{diag}(HT)|_{S_\perp}$ are polynomial in $(H, T)$, so full column rank holds off a proper algebraic subvariety \emph{provided a single witness exists}; the dimension count $K - q \ge p - 1$ makes a witness expected but, for a \emph{fixed} named basis (e.g. Hadamard with a specific low-pass $S$), the rank must be verified directly --- the correct fixed-basis statement is rank-based, not dimensional. The unknown-support count $K \ge 2q + (p-1)$ arises from demanding that the $q$-sparse model avoid all of its gain-transformed copies $M_s\Sigma_q$, in analogy with the generic sparsity thresholds of Kech--Krahmer~\cite{KechKrahmer1603.07316} for generic bilinear maps; it is a heuristic prediction of the generic threshold, unverified for structured designs, and the uniform no-support-switching condition it summarizes is strictly stronger than (ii). \hfill$\blacksquare$\ \textnormal{(heuristic scope)}

\paragraph{What this does not claim.} Part (i) says nonnegativity fails to restore \emph{identifiability}; it may still improve estimator conditioning and is routinely useful as regularization --- no claim is made either way. Part (ii) is local: it excludes ambiguities in a ball around $T$ and is silent about isolated global aliases, which can exist even with a nonsingular linearization; global uniqueness needs a separate compactness/non-vanishing argument outside the ball. The criterion is also stated at fixed, \emph{known} support; for unknown support only the heuristic count of part (iii) is offered, and part (iii) makes no claim for any specific fixed basis without an explicit rank check.

\subsection{A.4 Remarks}

\paragraph{Remark A.1 (Fisher-theoretic reading).} Because the ambiguity of Theorem~A is an explicit real-analytic curve $t \mapsto (\ell + ts, T'(ts))$ along which the noiseless data are constant, every nonconstant $s \in S/\operatorname{span}\{\mathbf 1\}$ yields an exact null direction of the (gauge-quotient) Fisher information in any regular dominated noise model built on this record: the Cram\'er--Rao bound is infinite for any functional that varies along the curve. Here the implication ``explicit ambiguity curve $\Rightarrow$ Fisher null direction'' is exact, not merely first-order.

\paragraph{Remark A.2 (zero buckets).} If some $c_n = (MT)_n = 0$, two separate pathologies arise, which the statements above deliberately quarantine via (A4): the gain values $a_n$ at zero-carrier frames are entirely unobservable (a trivial extra ambiguity, of dimension $\dim(S \cap \{s : s = 0 \text{ on } \operatorname{supp}(c)\})$ beyond the count of Theorem~A), and log-domain processing of $R_n = 0$ is undefined, an estimation-side issue treated by the calibrated soft-log machinery of the low-photon analysis (Sec.~6; see Appendix~D), not an identifiability issue.

\section{Tall-Design Identifiability: the Parametric-Transversality Proof}

This appendix proves the tall-design identifiability results quoted as \textbf{Theorem~A$'$} in Sec.~4, restated in full as Theorem~A$'$(i)--(iv) in Sec.~B.8: the generic local threshold with a genuine \emph{iff} at $N=K+p-1$ for $K\ge2$, genuine (exact, positive-dimensional) failure below that wall, generic exact identifiability at $N\ge K+p$, uniform identifiability over the nonvanishing-carrier set at $N\ge 2K+p-1$, and the degenerate $K=1$ stratum. The proof route is \textbf{parametric transversality} (a rowwise design-perturbation witness plus Fubini), verified adversarially in review round R11; it holds for \textbf{every fixed drift subspace $S\ni\mathbf 1$} --- no genericity of $S$ is assumed anywhere, so the physical Fourier low-pass space $S_{\mathrm{LP}}$ is covered directly, with no transfer lemmas. Sharpness (necessity) of the two global thresholds is \textbf{not} claimed as a theorem; its status is stated precisely after Theorem~A$'$. This appendix \textbf{supersedes} the earlier incidence-variety/stratum-counting route (see Sec.~B.9); the superseded ``sharp iff'' statements at $N=K+p$ and $N=2K+p-1$ are withdrawn.

\subsection{B.0 Setting, notation, and the reduction}

The bucket record is $R_n=a_nB_n$ with $B=MT$: in vector form the data map is
\[
\begin{gathered}
\Phi_M(T,s)\;=\;D_{e^{s}}\,M\,T,\qquad M\in\mathbb R^{N\times K},\ N\ge K,\\
s\in S,\ T\in\mathbb R^{K},
\end{gathered}
\]
where $D_x=\operatorname{diag}(x)$, $e^{s}$ is componentwise, and the log-gain trajectory is confined to the fixed, known drift class $S\subset\mathbb R^{N}$ with $\dim S=p\ge2$ and $\mathbf 1\in S$. Write $y=MT$ for the carrier and define
\[
\begin{gathered}
\mathcal U=\{(M,T):(MT)_n\neq0\ \forall n\},\\
\mathcal U_M=\{T\in\mathbb R^K:(MT)_n\neq0\ \forall n\},\\
\mathcal R=\{M\in\mathbb R^{N\times K}:\operatorname{rank}M<K\}.
\end{gathered}
\]
$\mathcal U$ is open with null complement (each $\{m_n^\top T=0\}$ is a proper algebraic subset); $\mathcal R$ is a proper algebraic subset (vanishing of all $K\times K$ minors), hence Lebesgue-null. \textbf{Every exceptional set below explicitly includes $\mathcal R$ (or $\mathcal R\times\mathbb R^K$).} The gauge is $(T,s)\sim(cT,\,s-\log c\,\mathbf 1)$, $c>0$; identifiability is always modulo this gauge. ``For a.e.''\ always refers to Lebesgue measure on the stated parameter space.

\begin{namedthm}{Lemma B.0 (reduction)}
Assume $\operatorname{rank}M=K$ and $(M,T)\in\mathcal U$, and fix any $s\in S$. Then $(T',s')\in\mathbb R^K\times S$ is an alias of $(T,s)$ --- i.e.\ $\Phi_M(T',s')=\Phi_M(T,s)$ --- iff, with $d:=s-s'\in S$,
\begin{equation}
M\,T'\;=\;D_{e^{d}}\,M\,T.
\tag{B.1}
\end{equation}
If $d=a\mathbf 1$ is constant, then (B.1) reads $MT'=e^{a}MT$ and full column rank gives $T'=e^{a}T$ --- exactly the gauge orbit with $c=e^{a}$; conversely every gauge point solves (B.1) with constant $d$. Hence $(M,T)$ is identifiable iff (B.1) has no solution $(T',d)$ with $d\in S\setminus\mathbb R\mathbf 1$. Moreover, fixing a complement $S_0$ of $\mathbb R\mathbf 1$ in $S$ ($\dim S_0=p-1$) and writing $d=d_0+a\mathbf 1$, the normalization $(T',d)\mapsto(e^{-a}T',d_0)$ maps solutions to solutions; so $(M,T)$ is identifiable iff (B.1) has no solution with $d\in S_0\setminus\{0\}$.
\end{namedthm}

\noindent\textbf{Proof.} $D_{e^{s}}$ is invertible, so record equality is equivalent to (B.1) with $d=s-s'$. Constant $d=a\mathbf 1$: row $n$ reads $m_n^\top T'=e^{a}y_n$, i.e.\ $MT'=e^{a}MT$, and $\operatorname{rank}M=K$ forces $T'=e^{a}T$, which is the gauge point $(cT,\,s-\log c\,\mathbf 1)$ with $c=e^{a}>0$. The equivariance $(T',d)\mapsto(e^{a}T',d+a\mathbf 1)$ is a direct substitution in (B.1). \hfill$\blacksquare$

Two structural remarks. \emph{(a) The reduction eliminates $s$:} identifiability is a property of $(M,T,S)$ alone, uniform over the true drift --- this is the source of the ``for every $s\in S$ simultaneously'' quantifier in every theorem below. \emph{(b) Full column rank is necessary:} for rank-deficient $M$ the equivalence ``constant $d$ $\Leftrightarrow$ gauge-trivial'' is \textbf{false} even on $\mathcal U$. Counterexample ($\operatorname{rank}M=1$): $N=K=2$, both rows of $M$ equal to $(0,1)$, $T=(0,1)^\top$, $T'=(1,1)^\top$: $MT=MT'=(1,1)^\top$, so $(T',d=0)$ solves (B.1) with constant $d$, yet $T'\notin\mathbb R T$ --- an object-kernel ambiguity invisible to the normalized $d\neq0$ incidence set. This is why $\mathcal R$ is carried explicitly in every exceptional set.

\paragraph{What this does not claim.} The reduction is exact but presupposes $\operatorname{rank}M=K$ and a nonvanishing carrier; neither is a restriction for the a.e.\ statements (both failure sets are null), but both must be --- and are --- carried explicitly.

\subsection{B.1 Lemma B.1 (row dichotomy)}

\begin{namedthm}{Lemma B.1}
Let $(M,T)\in\mathcal U$ and let $(T',d)$ solve (B.1) with $d$ nonconstant ($d\notin\mathbb R\mathbf 1$). Then $T'$ is not a scalar multiple of $T$; consequently $T'-e^{d_n}T\neq0$ for every $n=1,\dots,N$.
\end{namedthm}

\noindent\textbf{Proof.} Suppose $T'=\beta T$. Row $n$ of (B.1): $\beta y_n=e^{d_n}y_n$; since $y_n\neq0$ on $\mathcal U$, $e^{d_n}=\beta$ for all $n$, so $d$ is constant --- contradiction. For the display: if $T'=e^{d_n}T$ for a single $n$, then $T'\parallel T$ ($T\neq0$ because $y=MT\neq0$). \hfill$\blacksquare$

Only \emph{global} nonconstancy of $d$ is used --- partial constancy (some equal components) is harmless --- and the only division is by $y_n$, legitimate on $\mathcal U$.

\subsection{B.2 Lemma B.2 (joint submersivity)}

Define $F:\mathbb R^{N\times K}\times\mathbb R^K\times\mathbb R^K\times S_0\to\mathbb R^N$,
\begin{equation}
\begin{gathered}
F(M,T,T',d)\;=\;MT'-D_{e^{d}}MT,\\
F_n=m_n^\top\bigl(T'-e^{d_n}T\bigr).
\end{gathered}
\tag{B.2}
\end{equation}

\begin{namedthm}{Lemma B.2}
At every zero of $F$ with $(M,T)\in\mathcal U$ and $d\in S_0\setminus\{0\}$, the differential of $F$ with respect to $M$ alone is surjective onto $\mathbb R^N$.
\end{namedthm}

\noindent\textbf{Proof.} Perturb row $n$ only: $\delta M=e_n\,\delta m_n^\top$ gives $\delta F=e_n\cdot\delta m_n^\top(T'-e^{d_n}T)$, which spans $\mathbb R e_n$ as $\delta m_n$ ranges over $\mathbb R^K$, provided $T'-e^{d_n}T\neq0$ --- which holds at every such zero for every $n$ by Lemma B.1. Rows are decoupled, so the image contains every $e_n$. \hfill$\blacksquare$

Only $\delta M$ is used: no genericity of $S$, no structure of $d$ beyond nonconstancy, and the object is untouched. This rowwise witness is uniform over the whole zero set --- the ingredient the superseded route lacked.

\subsection{B.3 A measure lemma, and generic exact identifiability at \texorpdfstring{$N\ge K+p$}{N >= K+p}}

\begin{namedthm}{Lemma B.3 ($C^1$ image of a lower-dimensional manifold is null)}
Let $X$ be a second-countable $C^1$ manifold with $\dim X=m<q$, and $f:X\to\mathbb R^q$ a $C^1$ map. Then $f(X)$ has $q$-dimensional Lebesgue measure zero.
\end{namedthm}

\noindent\textbf{Proof.} Second countability gives a countable chart cover $\varphi_i:U_i\to\mathbb R^m$, and each chart domain is exhausted by countably many compact cubes; by countable subadditivity it suffices to treat one compact cube $Q$ of side $\ell$. On $Q$, $g=f\circ\varphi_i^{-1}$ is $C^1$, hence Lipschitz with some constant $L$. Subdividing $Q$ into $k^m$ subcubes of side $\ell/k$, each image lies in a ball of radius $L\sqrt m\,\ell/k$ in $\mathbb R^q$, so the outer measure of $g(Q)$ is at most $k^m\,c_q(2L\sqrt m\,\ell/k)^q=C\,k^{m-q}\to0$ since $m<q$. \hfill$\blacksquare$

No properness or compactness of $f$ is needed; the zero manifolds below are embedded in Euclidean space, hence second countable and $\sigma$-compact automatically. The statement is \textbf{false} for merely continuous maps (space-filling curves) --- the $C^1$ regularity supplied by Lemma B.2 is load-bearing.

\begin{namedthm}{Theorem B.4 (generic exact identifiability; Theorem A$'$(ii))}
Fix $S\ni\mathbf 1$, $\dim S=p\ge2$. If $N\ge K+p$, there is a Lebesgue-null set $E_{\mathrm{gen}}\subset\mathbb R^{N\times K}\times\mathbb R^K$ such that every $(M,T)\notin E_{\mathrm{gen}}$ is identifiable: for every $s\in S$, the only $(T',s')\in\mathbb R^K\times S$ with $\Phi_M(T',s')=\Phi_M(T,s)$ are the gauge orbit of $(T,s)$.
\end{namedthm}

\noindent\textbf{Proof.} Work on the open set $\Omega=\{(M,T,T',d):(M,T)\in\mathcal U,\ T'\in\mathbb R^K,\ d\in S_0\setminus\{0\}\}$. By Lemma B.2, $0$ is a regular value of $F$ on $\Omega$, so $Z=F^{-1}(0)\cap\Omega$ is an embedded real-analytic submanifold of codimension $N$:
\[
\dim Z=(NK+K)+K+(p-1)-N.
\]
The projection $\pi:Z\to\mathbb R^{N\times K}\times\mathbb R^K$ onto $(M,T)$ is $C^1$ and
\[
\dim Z-\dim(\mathbb R^{N\times K}\times\mathbb R^K)=K+p-1-N\le-1
\]
when $N\ge K+p$, so $\pi(Z)$ is Lebesgue-null by Lemma B.3. By the reduction (Lemma B.0, which requires $\operatorname{rank}M=K$),
\[
\{\text{non-identifiable pairs in }\mathcal U\}\;\subset\;(\mathcal R\times\mathbb R^K)\,\cup\,\pi(Z),
\]
and the full non-identifiable set is contained in $(\mathcal R\times\mathbb R^K)\cup\pi(Z)\cup\mathcal U^{c}$ --- three null sets. Set $E_{\mathrm{gen}}$ to their union. \hfill$\blacksquare$

\paragraph{Remark B.4.1 (quantifier order, proved).} For a.e.\ $(M,T)$: for \textbf{every} $s\in S$, no non-gauge alias --- the null set $E_{\mathrm{gen}}$ is independent of $s$, because the reduction eliminated $s$ before any measure-theoretic step. This does \textbf{not} imply that one design works for every object; that is Theorem B.5 at its own larger threshold. ($T'=0$ needs no special handling: $T'=0$ forces $D_{e^{d}}y=0$, impossible on $\mathcal U$.)

\paragraph{Remark B.4.2 (fixed-object slice).} For each \textbf{fixed} $T\neq0$: a.e.\ $M$ identifies that particular $T$ when $N\ge K+p$ (the exceptional design set may depend on $T$). Proof: freeze $T$, run the same argument with $F_T(M,T',d)=MT'-D_{e^{d}}MT$ on $\{M:(M,T)\in\mathcal U\}\times\mathbb R^K\times(S_0\setminus\{0\})$; Lemma B.2's witness uses only $\delta M$, the zero set has dimension $NK+K+(p-1)-N\le NK-1$ when $N\ge K+p$, and Lemma B.3 applies; add $\mathcal R$ and the proper algebraic set $\{M:(M,T)\notin\mathcal U\}$. This slice form is what covers the fixed structured objects of the Fig.~6 protocol.

\paragraph{What this does not claim.} No claim of uniformity over objects for a fixed design (see Theorem B.5 and its counterexample); no claim at the boundary $N=K+p-1$ (necessity is open --- Sec.~B.8, sharpness paragraph); nothing about conditioning, noise, or estimators; deterministic structured designs are not covered by any a.e.\ statement.

\subsection{B.4 Uniform identifiability on the nonvanishing-carrier set at \texorpdfstring{$N\ge 2K+p-1$}{N >= 2K+p-1}}

\begin{namedthm}{Theorem B.5 (uniform sufficiency on $\mathcal U_M$; Theorem A$'$(iii))}
Fix $S\ni\mathbf 1$, $\dim S=p\ge2$. If $N\ge 2K+p-1$, there is a Lebesgue-null set $E_{\mathrm{unif}}\subset\mathbb R^{N\times K}$ such that for every $M\notin E_{\mathrm{unif}}$: every $T\in\mathcal U_M$ is identifiable, for every $s\in S$ simultaneously.
\end{namedthm}

\noindent\textbf{Proof.} The solution set of (B.1) is scaling-equivariant --- $(T,T',d)$ solves iff $(cT,cT',d)$ does, $c\neq0$ --- and nonvanishing carrier and nonconstancy are scale-invariant, so it suffices to rule out solutions with $|T|=1$. Let
\[
\begin{gathered}
\Omega'=\{(M,T,T',d):|T|=1,\ (M,T)\in\mathcal U,\\
\hspace{5.2em} T'\in\mathbb R^K,\ d\in S_0\setminus\{0\}\},
\end{gathered}
\]
an open subset of $\mathbb R^{N\times K}\times S^{K-1}\times\mathbb R^K\times S_0$. Lemma B.2 holds verbatim on $\Omega'$ (the sphere constrains $T$, not $M$), so $Z'=F^{-1}(0)\cap\Omega'$ is an embedded real-analytic manifold of codimension $N$ with
\[
\begin{aligned}
\dim Z'&=NK+(K-1)+K+(p-1)-N\\
&=NK+(2K+p-2-N),
\end{aligned}
\]
and the projection $\pi':Z'\to\mathbb R^{N\times K}$ onto $M$ has deficit $2K+p-2-N\le-1$ when $N\ge2K+p-1$; $\pi'(Z')$ is null by Lemma B.3. Set $E_{\mathrm{unif}}=\mathcal R\cup\pi'(Z')$ (rank-deficient designs added explicitly: the reduction needs $\operatorname{rank}M=K$). If $M\notin E_{\mathrm{unif}}$ and $(T',d)$ solved (B.1) for some $T\in\mathcal U_M$, $d\in S_0\setminus\{0\}$, then rescaling by $c=1/|T|$ would put $(M,cT,cT',d)\in Z'$, i.e.\ $M\in\pi'(Z')$ --- contradiction. Lemma B.0 concludes. \hfill$\blacksquare$

\paragraph{Remark B.5.1 (why $\mathcal U_M$ cannot be dropped: a zero-carrier counterexample at the threshold).} ``Every nonzero object'' is \textbf{false} even at $N=2K+p-1$. Take $K=2$, $p=2$, $N=5=2K+p-1$, $S=\operatorname{span}\{\mathbf 1,e_1\}$, $T=(1,1)^\top$, and $M$ with rows $(1,-1),(1,0),(0,1),(1,1),(1,2)$, so $MT=(0,1,1,2,3)^\top$ --- full column rank, first carrier coordinate zero. For every $t\neq0$, $d=t\,e_1\in S$ is nonconstant and $(T',d)=(T,\,t\,e_1)$ solves (B.1): row 1 reads $m_1^\top T'=e^{t}\cdot0=0$, rows 2--5 are unchanged. So $(T,s)$ and $(T,s-t\,e_1)$ produce identical records for every $t$ --- a one-parameter family of non-gauge aliases. This mechanism is available for \textbf{any} design: every object on the hyperplane $\{m_1^\top T=0\}$ acquires such aliases whenever $S$ contains a vector supported on the zero-carrier row. The correct conclusion is the stated one: for a.e.\ $M$, every $T\in\mathcal U_M$, every $s\in S$.

\paragraph{What this does not claim.} Uniformity is over nonvanishing-carrier objects and drifts, not over designs: an adversarially chosen $M$ can defeat the bound. Necessity of $2K+p-1$ is not claimed (Sec.~B.8). No claim about recovery algorithms or noise.

\subsection{B.5 Lemma W: position of the twisted column space}

For $(M,T)\in\mathcal U$ with $\operatorname{rank}M=K$ define
\[
W(M,T)\;:=\;D_y^{-1}\operatorname{col}(M)\subset\mathbb R^N,\qquad y=MT,
\]
a $K$-dimensional subspace with $\mathbf 1\in W$ (since $D_y\mathbf 1=y\in\operatorname{col}(M)$). Aliases relate to $W$ by: \emph{$d\in S$ solves (B.1) for some $T'$ iff $e^{d}\in W(M,T)$} --- indeed $D_{e^{d}}y=D_y e^{d}$, so (B.1) reads $MT'=D_ye^{d}$, i.e.\ $e^{d}=D_y^{-1}MT'\in W$; conversely $e^{d}=D_y^{-1}Mu$ gives the (unique, by full column rank) solution $T'=u$.

\begin{namedthm}{Lemma W}
Fix any subspace $S\ni\mathbf 1$ with $\dim S=p\le N$, and $K\ge1$. For a.e.\ $(M,T)\in\mathcal U$, with $W=W(M,T)$:
\begin{equation}
\begin{aligned}
\dim(S\cap W)&=1+\max\{0,\ (p-1)+(K-1)-(N-1)\}\\
&=1+\max\{0,\,K+p-N-1\}.
\end{aligned}
\tag{B.3}
\end{equation}
No genericity of $S$ is assumed or used --- only $\mathbf 1\in S$.
\end{namedthm}

\noindent\textbf{Proof.} \emph{(Step 1: basis completion; fixed $T$.)} Fix $T\neq0$ and complete it to a basis $B_T=[T,b_2,\dots,b_K]\in GL_K$. The map $M\mapsto MB_T=(y,v_2,\dots,v_K)$ is a linear isomorphism of $\mathbb R^{N\times K}$, hence maps null sets to null sets both ways: a.e.\ $M$ $\Leftrightarrow$ a.e.\ $(y,v_2,\dots,v_K)$.

\emph{(Step 2: fixed-$y$ isomorphism.)} For fixed $y\in(\mathbb R^*)^N$ (a co-null condition), set $w_j:=D_y^{-1}v_j$. For fixed $y$ this is a linear isomorphism of $\mathbb R^{N(K-1)}$ in $(v_2,\dots,v_K)$ --- it preserves ``a.e.''\ and introduces no correlation between the $w_j$ and the fixed $S$. Then $W=D_y^{-1}\operatorname{span}(y,v_2,\dots,v_K)=\operatorname{span}(\mathbf 1,w_2,\dots,w_K)$.

\emph{(Step 3: quotient.)} Let $q:\mathbb R^N\to\mathbb R^N/\mathbb R\mathbf 1$ (dimension $N-1$). $A:=q(S)$ has dimension $p-1$ (the kernel $\mathbb R\mathbf 1$ lies in $S$); $q(W)=\operatorname{span}(qw_2,\dots,qw_K)$.

\emph{(Step 4: general position against the fixed $A$.)} For a.e.\ $(w_2,\dots,w_K)$: the $q(w_j)$ are independent and $\dim(A\cap q(W))=\max\{0,(p-1)+(K-1)-(N-1)\}$. Failure is the vanishing of suitable maximal minors of the matrix $[A\text{-basis}\mid qw_2\cdots qw_K]$ --- a polynomial condition in the $w_j$, not identically zero (a witness exists: against one fixed subspace, vectors can be chosen stepwise outside a finite union of proper subspaces), hence a proper algebraic --- null --- subset. This is general position with respect to the fixed $A$ only; no genericity of $S$.

\emph{(Step 5: $q(S\cap W)=q(S)\cap q(W)$, via $\mathbf 1\in W$.)} ``$\subseteq$'' is trivial. ``$\supseteq$'': if $x\in S$, $x'\in W$ and $q(x)=q(x')$, then $x-x'\in\mathbb R\mathbf 1\subseteq W$, so $x\in W$, i.e.\ $x\in S\cap W$. The kernel of $q$ on $S\cap W$ is exactly $\mathbb R\mathbf 1$ (as $\mathbf 1\in S\cap W$), whence $\dim(S\cap W)=1+\dim(q(S)\cap q(W))$.

\emph{(Step 6: Fubini twice, with Borel measurability.)} The quantifier order so far: for every fixed $T\neq0$, for a.e.\ $y$, for a.e.\ $(v_j)$ --- hence, by Steps 1--2 and Fubini on $\mathbb R^N\times\mathbb R^{N(K-1)}$, for every fixed $T\neq0$, for a.e.\ $M$, formula (B.3) holds. The failure set $B=\{(M,T)\in\mathcal U:\text{(B.3) fails}\}$ is Borel: the dimension condition is a Boolean combination of rank conditions (minor vanishing/nonvanishing) on matrices whose entries are rational in $(M,T)$ on $\mathcal U$ (the $w_j$ are entries of $M$ divided by coordinates of $MT$). Measurability licenses a second Fubini, over $T$: every $T$-slice of $B$ is $M$-null, so $B$ is null in $\mathbb R^{N\times K}\times\mathbb R^K$. \hfill$\blacksquare$

\paragraph{Remark W.1 ($K=1$).} For $K=1$ the lemma \textbf{holds}: $W=\mathbb R\mathbf 1$ identically, $\dim(S\cap W)=1$, and (B.3) also gives $1+\max\{0,p-N\}=1$ since $p\le N$. What degenerates at $K=1$ is the necessity of the global thresholds (Sec.~B.7), not Lemma W.

\paragraph{What this does not claim.} (B.3) is an a.e.\ statement over $(M,T)$; it says nothing about any \emph{particular} design (e.g.\ structured $M$), and nothing about the conditioning of the intersection, only its dimension.

\subsection{B.6 The generic local threshold (\texorpdfstring{$K\ge2$}{K >= 2}) and genuine failure below the wall}

\paragraph{Kernel derivation (independent; no import from the superseded route).} Fix $(M,T)\in\mathcal U$ with $\operatorname{rank}M=K$ and any $s\in S$. The differential of $\Phi_M$ at $(T,s)$ is
\begin{equation}
D\Phi_M(T,s)[\delta T,\delta s]\;=\;D_{e^{s}}\bigl(M\,\delta T+D_y\,\delta s\bigr),\qquad y=MT.
\tag{B.4}
\end{equation}
The diagonal factor is invertible, so $(\delta T,\delta s)\in\ker\Leftrightarrow M\delta T=-D_y\delta s$; the right side lies in $\operatorname{col}(M)$ iff $\delta s\in D_y^{-1}\operatorname{col}(M)=W(M,T)$, and then $\delta T=-(M^\top M)^{-1}M^\top D_y\delta s$ is unique. Hence $\ker D\Phi_M(T,s)\cong S\cap W$ via $(\delta T,\delta s)\mapsto\delta s$, and
\[
\begin{gathered}
\dim\ker D\Phi_M(T,s)=\dim(S\cap W),\\
\operatorname{rank}D\Phi_M(T,s)=K+p-\dim(S\cap W),
\end{gathered}
\]
independently of $s$. The gauge tangent is the line spanned by $(-T,\mathbf 1)$ (from $c\mapsto(cT,s-\log c\,\mathbf 1)$), and it lies in the kernel since $MT=D_y\mathbf 1$; as $\mathbf 1\in S\cap W$ always, the maximal gauge-compatible rank is $K+p-1$ and the defect from it is $\dim(S\cap W)-1$. This derivation is self-contained: nothing is imported from the superseded route, and no genericity of $S$ is used. By Lemma W, for a.e.\ $(M,T)\in\mathcal U$ and every $s$:
\[
\operatorname{rank}D\Phi_M(T,s)=\min\{N,\ K+p-1\}.
\]

\begin{namedthm}{Theorem B.6 (generic local iff; Theorem A$'$(i), rank part)}
Fix $S\ni\mathbf 1$, $\dim S=p\ge2$, and $K\ge2$. There is a Lebesgue-null set $E_{\mathrm{loc}}\subset\mathbb R^{N\times K}\times\mathbb R^K$ --- take $E_{\mathrm{loc}}=(\mathcal R\times\mathbb R^K)\cup\mathcal U^c\cup B$ with $B$ the failure set of Lemma W --- such that for every $(M,T)\notin E_{\mathrm{loc}}$ with $T\in\mathcal U_M$ and every $s\in S$: $\operatorname{rank}D\Phi_M(T,s)=\min\{N,K+p-1\}$, and $(T,s)$ is locally identifiable modulo gauge \textbf{iff} $N\ge K+p-1$.
\end{namedthm}

\noindent\textbf{Proof.} The rank formula is the kernel derivation plus Lemma W. \emph{Sufficiency at $N\ge K+p-1$ (ordinary local injectivity modulo gauge):} write $S=\mathbb R\mathbf 1\oplus S_0$; the bookkeeping
\[
S\cap W=\mathbb R\mathbf 1\oplus(S_0\cap W)
\]
holds because for $x\in S\cap W$, $x=a\mathbf 1+x_0$ with $x_0\in S_0$ gives $x_0=x-a\mathbf 1\in W$. At maximal rank $\dim(S\cap W)=1$, so $S_0\cap W=0$. Restrict $\Phi_M$ to the gauge slice $\mathbb R^K\times(s+S_0)$, of dimension $K+p-1$: by (B.4) its differential at $(T,s)$ has kernel $\cong S_0\cap W=0$, so the restricted map is an immersion there, hence (constant-rank/immersion theorem) a local embedding --- locally injective on the slice. Any nearby alias is gauge-normalized into the slice by a factor $c$ near $1$, so ordinary local identifiability modulo gauge holds. \emph{Necessity at $N\le K+p-2$:} Theorem B.7 below supplies a positive-dimensional manifold of exact aliases through $(T,s)$ --- genuine failure, not merely a rank diagnostic. \hfill$\blacksquare$

\begin{namedthm}{Theorem B.7 (genuine failure below the wall; Theorem A$'$(i), below-wall part)}
Fix $S\ni\mathbf 1$, $K\ge2$, $N\le K+p-2$. For every $(M,T)\notin E_{\mathrm{loc}}$ with $T\in\mathcal U_M$ and every $s\in S$: through $(T,s)$ passes a real-analytic manifold of exact aliases of dimension $K+p-N-1\ge1$; every point with $d\neq0$ is gauge-inequivalent to $(T,s)$, and the aliased objects satisfy $T'(d)\to T$ as $d\to0$.
\end{namedthm}

\noindent\textbf{Proof.} Let $P=\Pi_{W^\perp}$ be the orthogonal projection onto $W^\perp$ ($\dim W^\perp=N-K$ by full column rank) and define the real-analytic map $\varphi:S_0\to W^\perp$, $\varphi(d)=P(e^{d})$. Then $e^{d}\in W\Leftrightarrow\varphi(d)=0$ and $\varphi(0)=P\mathbf 1=0$. The differential is $D\varphi_0=P|_{S_0}$ ($\tfrac{d}{dt}\big|_0e^{td}=d$ componentwise; $P$ linear). By the bookkeeping $S\cap W=\mathbb R\mathbf 1\oplus(S_0\cap W)$ and Lemma W with $N\le K+p-2$ (the max in (B.3) is attained), $\dim(S_0\cap W)=K+p-N-1$, so
\[
\operatorname{rank}D\varphi_0=(p-1)-(K+p-N-1)=N-K=\dim W^\perp:
\]
$D\varphi_0$ is surjective. By the (real-analytic) submersion/implicit-function theorem, $\varphi^{-1}(0)$ is, near $0$, a real-analytic manifold of dimension $(p-1)-(N-K)=K+p-N-1\ge1$. For each $d$ in it, $e^{d}\in W$; with the left inverse $L=(M^\top M)^{-1}M^\top$ set
\[
T'(d)\;:=\;L\,D_y\,e^{d},
\]
so that $MT'(d)=D_ye^{d}$ ($D_ye^{d}=D_{e^{d}}y\in\operatorname{col}(M)$ precisely when $e^{d}\in W$): $(T'(d),d)$ solves (B.1), $T'(0)=Ly=T$, and $T'(d)\to T$ by continuity. Every $d\neq0$ lies in $S_0\setminus\{0\}$, hence $d\notin\mathbb R\mathbf 1$: a genuine non-gauge alias --- $(T,s)$ and $(T'(d),s-d)$ produce identical records but are gauge-inequivalent. \hfill$\blacksquare$

\paragraph{Consistency check at $N=K$.} At $N=K$ (square, $\operatorname{rank}M=K$): $W=\mathbb R^N$, $W^\perp=0$, $\varphi\equiv0$, and the alias manifold is a full neighborhood of $0$ in $S_0$ --- ambiguity dimension $p-1$, exactly Theorem~A. There is no hidden $N\ge K+1$ assumption, and no contradiction ($p\ge2$ gives $K<K+p\le2K+p-1$, so the square case lies strictly below both global thresholds).

\paragraph{What this does not claim.} The local iff is generic in $(M,T)$ --- a particular structured design must be checked separately. Local identifiability at $N\ge K+p-1$ does not imply global uniqueness (far aliases at the boundary are exactly the open sharpness question, Sec.~B.8). Nothing about the singular values of the Jacobian, i.e.\ conditioning.

\subsection{B.7 The \texorpdfstring{$K=1$}{K=1} stratum}

\begin{namedthm}{Theorem B.8 ($K=1$: global identifiability; Theorem A$'$(iv))}
Let $K=1$. For every $M$, every $T\in\mathcal U_M$, and every $s\in S$, the pair $(T,s)$ is globally identifiable modulo gauge, for every admissible $N$.
\end{namedthm}

\noindent\textbf{Proof.} $T\in\mathcal U_M$ means $T\neq0$ and $m_n\neq0$ for all $n$ (scalars $m_nT\neq0$). Any solution of (B.1) satisfies $e^{d_n}=m_nT'/(m_nT)=T'/T$ for every $n$, so $d$ is constant and $c:=T'/T=e^{d_1}>0$ --- the gauge. \hfill$\blacksquare$

The local threshold at $K=1$ reads $N\ge K+p-1=p$, which is \textbf{automatically satisfied} ($S\subseteq\mathbb R^N$ forces $p\le N$) --- the threshold is vacuous, not ``inapplicable''. Consequently neither $N\ge K+p$ nor $N\ge2K+p-1$ is necessary at $K=1$: corner strata break naive universal-necessity claims, which is why all necessity statements in this appendix are restricted to $K\ge2$.

\subsection{B.8 Theorem A\texorpdfstring{$'$}{'} --- the manuscript statement, and sharpness status}

\begin{namedthm}{Theorem A$'$ (proved tall-design identifiability results)}
Let $K\ge1$, $N\ge K$, and let $S\le\mathbb R^N$ be a fixed, known linear subspace of dimension $p\ge2$ with $\mathbf 1\in S$. For $M\in\mathbb R^{N\times K}$ define $\Phi_M(T,s)=D_{e^{s}}MT$ and $\mathcal U_M=\{T\in\mathbb R^K:(MT)_n\neq0$ for every $n\}$. Parameters are identified modulo the positive global-scale gauge $(T,s)\sim(cT,\,s-\log c\,\mathbf 1)$, $c>0$.

(i) \textbf{(Generic local threshold and genuine failure below it, $K\ge2$.)} There is a Lebesgue-null set $E_{\mathrm{loc}}\subset\mathbb R^{N\times K}\times\mathbb R^K$ such that for every $(M,T)\notin E_{\mathrm{loc}}$ with $T\in\mathcal U_M$ and every $s\in S$: $\operatorname{rank}D\Phi_M(T,s)=\min\{N,K+p-1\}$. Consequently $(T,s)$ is locally identifiable modulo gauge iff $N\ge K+p-1$. If $N\le K+p-2$ then a gauge-normalized neighborhood of $(T,s)$ contains a real-analytic manifold of exact aliases of dimension $K+p-N-1\ge1$, and every other sufficiently nearby point of it is gauge-inequivalent.

(ii) \textbf{(Generic exact-identifiability sufficiency.)} If $N\ge K+p$, there is a Lebesgue-null set $E_{\mathrm{gen}}\subset\mathbb R^{N\times K}\times\mathbb R^K$ such that for every $(M,T)\notin E_{\mathrm{gen}}$ and every $s\in S$: $D_{e^{s'}}MT'=D_{e^{s}}MT$ with $(T',s')\in\mathbb R^K\times S$ implies $T'=cT$, $s'=s-\log c\,\mathbf 1$ for a unique $c>0$. Exact identifiability holds simultaneously for every possible true drift $s\in S$. This statement is \textbf{not} uniform over all objects for one fixed $M$.

(iii) \textbf{(Uniform exact-identifiability sufficiency on the nonvanishing-carrier set.)} If $N\ge2K+p-1$, there is a Lebesgue-null set $E_{\mathrm{unif}}\subset\mathbb R^{N\times K}$ such that for every $M\notin E_{\mathrm{unif}}$, every $T\in\mathcal U_M$, every $s\in S$, and every $(T',s')\in\mathbb R^K\times S$: $D_{e^{s'}}MT'=D_{e^{s}}MT$ implies $T'=cT$, $s'=s-\log c\,\mathbf 1$ for a unique $c>0$. The same generic design works simultaneously for all nonvanishing-carrier objects and all true drifts in $S$.

(iv) \textbf{(Degenerate one-dimensional object stratum.)} If $K=1$ then for every $M$ and every $T\in\mathcal U_M$ and every $s\in S$, the pair $(T,s)$ is globally identifiable modulo gauge, for every admissible $N$. Neither $N\ge K+p$ nor $N\ge2K+p-1$ is necessary when $K=1$.
\end{namedthm}

\noindent\textbf{Proof.} (i) = Theorems B.6--B.7 ($E_{\mathrm{loc}}$ from Lemma W); (ii) = Theorem B.4; (iii) = Theorem B.5; (iv) = Theorem B.8 (which holds for every $M$, not just a.e.). \hfill$\blacksquare$

\begin{statusblock}\textbf{Sharpness status.} Parts (ii) and (iii) are sufficient conditions, not proved if-and-only-if thresholds. For $K\ge2$, the below-wall part of (i) proves generic exact failure whenever $N\le K+p-2$. The remaining necessity assertion for the $K+p$ bound is the boundary case $N=K+p-1$: dimension counting predicts nonlocal aliases there, but dominance of the bad-pair projection has not been proved. Likewise, dimension counting predicts that a generic design fails uniform identifiability for some object when $N\le2K+p-2$, but this necessity direction has not been proved in the intermediate regime. Thus $K+p$ and $2K+p-1$ are conjecturally sharp for $K\ge2$ under generic/nondegenerate drift geometry, not universal sharp thresholds for every fixed $S$. The $K=1$ stratum explicitly rules out any universal necessity claim.\end{statusblock}

\paragraph{What this does not claim.} Necessity of the two global thresholds (see the sharpness paragraph); anything about deterministic structured designs (the ordered-Hadamard obstruction of Theorem~A is unaffected, and no a.e.\ statement covers a fixed named design); anything about conditioning, noise robustness, or estimators --- uniqueness at a threshold is compatible with arbitrarily bad stability; and no coverage of objects constrained to lower-dimensional sets (e.g.\ exactly sparse) --- open constraint sets (e.g.\ positivity with nonvanishing carrier) inherit all a.e.\ statements by restriction. The Fig.~6 experiment (Sec.~9) is evidence for the \emph{local rank wall} of part (i) --- with the fixed-object slice covered by Remark B.4.2 --- and probes algorithmic recovery; it neither tests nor proves the sharpness predictions.

\subsection{B.9 Relation to the superseded incidence-variety route}

Earlier versions of this appendix proved Theorem~A$'$ through a collision incidence variety: a stratum-counting lemma over the gain-ratio manifold, a dominant-stratum rank hypothesis $\operatorname{rank}[M,-D_bM]=\min(N,2K)$, a determinantal-transversality argument for the rank-drop locus, and two low-pass lemmas (a level-multiplicity bound for $S_{\mathrm{LP}}$ and a generic-rank lemma for $[M,-DM]$) feeding a three-level transfer proposition. \textbf{That route is superseded (R11), and its claims are withdrawn or replaced as follows.}

\begin{itemize}
\item The former ``sharp iff'' statements at $N\ge K+p$ (generic) and $N\ge2K+p-1$ (uniform) are \textbf{withdrawn}: their necessity directions rested on unproved dominant-projection and determinantal-transversality assertions. The proved statements are the sufficiency directions, now established by Theorems B.4--B.5 above, together with the genuinely sharp \emph{local} iff at $N=K+p-1$ (Theorem B.6) and genuine below-wall failure (Theorem B.7). Sharpness of the two global thresholds is a labeled prediction (Sec.~B.8).
\item The former uniform claim quantified over ``every nonzero object''; that is \textbf{false} at the stated threshold (Remark B.5.1) and is replaced by uniformity over the nonvanishing-carrier set $\mathcal U_M$.
\item The former route needed genericity of the drift family (or the low-pass transfer lemmas to emulate it). The parametric-transversality proof holds for \textbf{every fixed $S\ni\mathbf 1$}, so the low-pass transfer machinery (level-multiplicity bound, generic-rank lemma, three-level transfer hierarchy, and the $L_S$ multiplicity check) is no longer needed for Theorem~A$'$ and is not restated here; the withdrawn material remains available in the archived drafts for provenance.
\item The $K=1$, $N=p=2$ counterexample that broke the old claimed generic threshold is absorbed as the degenerate stratum, Theorem~A$'$(iv).
\end{itemize}

\paragraph{What this section does not claim.} Nothing in the superseded route is relied upon anywhere in Secs.~B.0--B.8; conversely, no assertion of the superseded route beyond the ones re-proved above should be quoted from earlier drafts.

\section{Carrier statistics and the stationarity anchor}

This appendix proves Proposition 1, formalizes condition ($\bigstar$) quantitatively, proves Proposition 2 (with counterexamples in both directions), and proves Theorem~B. Throughout, $R_n = a_n B_n$ is the bucket record, $B_n = \langle I_n, T\rangle = \sum_{j=1}^K I_{n,j} T_j$ the carrier, $\ell_n = \log a_n$ the log-gain, $S$ the drift class with $p=\dim S$, $S_1=\sum_j T_j$, $S_2=\|T\|_2^2$, $K_{\mathrm{eff}}=S_1^2/S_2$, $K_\infty=S_2/\|T\|_\infty^2$, $K_4=S_2^2/\sum_j T_j^4$, and $W$ denotes the window length of the blind estimator $\hat g_n = W^{-1}\sum_{k\in W_n} Y_k$ applied to $Y_n=\log R_n$. Asymptotic statements are for a triangular array indexed by $(K,N)$: the object $T=T^{(K)}$ and pattern law may vary with $K$, and limits are taken along $K,N\to\infty$.

\subsection{C.1 Proposition 1 --- random-basis carrier statistics}

\begin{namedthm}{Proposition 1 (carrier moments, stationarity, CLT, Berry--Esseen)}
Assume:
\begin{itemize}
\item \textbf{(P1)} $T_j\ge 0$ for all $j$, $S_1>0$, $S_2>0$.
\item \textbf{(P2)} The pattern entries $\{I_{n,j}\}$ are i.i.d.\ across $n$ and $j$ with mean $\mu_I>0$ and variance $\sigma_I^2>0$; the entries are nonnegative (so $B_n\ge 0$ a.s., and $B_n>0$ a.s.\ whenever $P(I=0)^{|\mathrm{supp}\,T|}=0$).
\item \textbf{(P3)} (for the CLT) Spikiness $K_\infty\to\infty$ along the array, with the family $\{(I-\mu_I)^2\}$ uniformly integrable (automatic for a fixed pattern law).
\item \textbf{(P4)} (for Berry--Esseen) $\nu_3 := \mathbb{E}|I-\mu_I|^3/\sigma_I^3<\infty$.
\end{itemize}
Then:

\textbf{(i) Exact moments and stationarity.} $\mathbb{E}B_n=\mu_I S_1$, $\operatorname{Var}B_n=\sigma_I^2 S_2$, hence $\mathrm{CV}_B=(\sigma_I/\mu_I)/\sqrt{K_{\mathrm{eff}}}$. The sequence $\{B_n\}_{n\ge1}$ is i.i.d., hence strictly stationary and white.

\textbf{(ii) CLT.} Under (P1)--(P3), $(B_n-\mu_I S_1)/(\sigma_I\sqrt{S_2}) \Rightarrow \mathcal N(0,1)$.

\textbf{(iii) Berry--Esseen.} Under (P4), for a universal constant $C$,
\[
\begin{aligned}
\sup_x\Big|\,P\Big(\tfrac{B_n-\mu_I S_1}{\sigma_I\sqrt{S_2}}\le x\Big)-\Phi(x)\Big|
&\;\le\; C\,\nu_3\,\frac{\sum_j |T_j|^3}{S_2^{3/2}}\\
&\;\le\; C\,\nu_3\,K_\infty^{-1/2}.
\end{aligned}
\]
The correct CLT small parameter is $K_\infty^{-1/2}$, \textbf{not} $K_{\mathrm{eff}}^{-1}$.

\textbf{(iv) Higher cumulants and the log carrier.} For $r\ge 2$ the cumulants are $\kappa_r(B_n)=\kappa_r(I)\sum_j T_j^r$, so the standardized law of $B_n$ depends on $T$ through \emph{all} ratios $\sum_j T_j^r/S_2^{r/2}$, $r\ge3$; these vanish as $K_\infty\to\infty$ but are not controlled by $K_{\mathrm{eff}}$ alone. Writing $B_n=\mu_I S_1(1+u_n)$ with $\mathbb{E}u_n=0$, $\operatorname{Var}u_n=(\sigma_I/\mu_I)^2/K_{\mathrm{eff}}$, a third-order Taylor expansion gives
\[
\begin{aligned}
\mathbb{E}\log B_n&=\log(\mu_I S_1)-\tfrac12\operatorname{Var}(u_n)+O(\mathbb{E}|u_n|^3),\\
\operatorname{Var}(\log B_n)&=\frac{(\sigma_I/\mu_I)^2}{K_{\mathrm{eff}}}+O(\mathbb{E}|u_n|^3).
\end{aligned}
\]
\end{namedthm}

\noindent\emph{Proof.} (i) Linearity of expectation and independence across $j$ give the mean and (Bienaym\'e) the variance; independence across $n$ gives the i.i.d.\ claim. (ii) Apply the Lindeberg--Feller CLT \cite[Thm.~27.2]{Billingsley1995} to the triangular array $X_j=T_j(I_{n,j}-\mu_I)$ with $s^2=\sigma_I^2 S_2$. Since $|X_j|\le \|T\|_\infty|I-\mu_I|$, the event $\{|X_j|>\varepsilon s\}$ is contained in $\{|I-\mu_I|>\varepsilon\sigma_I\sqrt{K_\infty}\}$, so the Lindeberg fraction is
\[
\begin{aligned}
&\frac{1}{\sigma_I^2 S_2}\sum_j T_j^2\,\mathbb{E}\big[(I-\mu_I)^2\mathbf 1\{|X_j|>\varepsilon s\}\big]\\
&\qquad\le\;\frac{1}{\sigma_I^2}\,\mathbb{E}\big[(I-\mu_I)^2\mathbf 1\{|I-\mu_I|>\varepsilon\sigma_I\sqrt{K_\infty}\}\big]\;\longrightarrow\;0
\end{aligned}
\]
by uniform integrability as $K_\infty\to\infty$. (iii) Esseen's inequality for sums of independent, non-identically distributed variables \cite[Ch.~V]{Petrov1975} gives the bound with third-moment ratio $\nu_3\sum_j|T_j|^3/S_2^{3/2}$, and $\sum_j|T_j|^3\le\|T\|_\infty S_2$ yields the $K_\infty^{-1/2}$ form. (iv) Cumulant homogeneity ($\kappa_r(cX)=c^r\kappa_r(X)$) and additivity over independent summands give the cumulant formula. The log expansion is Taylor's theorem with Lagrange remainder applied to $\log(1+u)$ on the event $\{|u_n|\le\tfrac12\}$. \hfill$\blacksquare$

\begin{statusblock}\textbf{Status of the log expansion (flagged).} Step (iv) is a \textbf{controlled moment expansion, not a distribution-free theorem}: making the $O(\mathbb{E}|u_n|^3)$ remainder rigorous, and ensuring integrability of $\log B_n$ near $B_n=0$, requires a high-probability bound $P(|u_n|>\tfrac12)$ decaying fast enough --- e.g.\ sub-exponential pattern entries plus $K_\infty$-driven concentration, or a strictly positive pattern floor $I\ge I_{\min}>0$. This step is a moment calculation; it enters only through the variance scale $\sigma_{\mathrm{abs}}^2=\sigma_{\mathrm{LR}}^2\asymp 1/K_{\mathrm{eff}}$ (per-frame, summable-absolute, and long-run variances all coincide for the i.i.d.\ carrier) entering rate constants, never in identifiability statements. Zero and near-zero buckets are excluded here; the calibrated soft-log $\psi_\alpha(c)=\log(c+\alpha)$ of Appendix~D removes this restriction with a photon-dependent sensitivity factor.\end{statusblock}

\paragraph{What this does not claim.} The proposition does not say the carrier law is determined by $(\mu_B,\sigma_B)$ --- that reduction holds only at the Gaussian-approximation level; the exact law feels every weighted cumulant of $T$. For gain recovery this residual object dependence enters $m_T=\mathbb{E}\log B_n$ only as a \emph{time-constant} scalar, absorbed by the global-scale gauge. It also does not assert cross-row decorrelation or window-energy whitening for structured (SRHT) carriers --- those require the Walsh-spectrum/$K_4$ conditions of Appendix~E, where $K_{\mathrm{eff}}$ alone is again insufficient.

\subsection{C.2 The quantitative stationarity condition (\texorpdfstring{$\bigstar$}{*})}

\begin{nameddef}{Definition ($\bigstar$)$(W,\epsilon,\delta)$}
Let $\{B_n\}_{n\le N}$ be the carrier and $W_n$ length-$W$ windows covering $\{1,\dots,N\}$. The carrier satisfies the \emph{quantitative stationarity anchor} with tolerance $\epsilon$ and confidence $1-\delta$ if there exists a scalar $m_T\in\mathbb R$ (allowed to depend on $T$, not on $n$) such that
\[
P\Big(\max_{1\le n\le N}\Big|\,W^{-1}\!\!\sum_{k\in W_n}\log B_k-m_T\Big|\le \epsilon\Big)\ \ge\ 1-\delta .
\]
We say \textbf{($\bigstar$) holds asymptotically} if for every $\delta>0$ there are windows $W=W_N\to\infty$, $W/N\to0$, with $\epsilon=\epsilon_{W,N,\delta}\to0$. The scalar $m_T$ is exactly the quantity absorbed by the global-scale gauge $a\mapsto ca$, $T\mapsto T/c$ (which shifts $m_T\mapsto m_T-\log c$, $\ell\mapsto\ell+\log c$, leaving $Y$ invariant).
\end{nameddef}

\subsection{C.3 Proposition 2 --- (\texorpdfstring{$\bigstar$}{*}) is a window-estimator criterion, not a universal one}

\begin{namedthm}{Proposition 2}
Assume the drift is slow in the sense of hypothesis (B3) below, with $C_\beta L\,(W/N)^\beta\to0$. Then:

\textbf{(a) Equivalence for the windowed corrector.} The blind windowed estimator $\hat g_n=W^{-1}\sum_{k\in W_n}Y_k$ is uniformly consistent for $\ell_n+m_T$ (in probability, uniformly over $n$) \textbf{if and only if} ($\bigstar$) holds asymptotically.

\textbf{(b) ($\bigstar$) is not necessary for identifiability.} There are designs violating ($\bigstar$) that are nonetheless identifiable: a tall design with $N\ge K+p$ is exactly algebraically identifiable modulo scale for a.e.\ design--object pair $(M,T)$, simultaneously for every drift $s\in S$ (Theorem~A$'$(ii); see Appendix~B.4 and Remark~B.4.1), with \emph{no} stationarity whatsoever --- the a.e.-$(M,T)$ quantifier does not cover deterministic structured (e.g.\ ordered) designs, which must be checked individually (Sec.~B.8); likewise, $K_0\ge p-1$ known support zeros restore local identifiability under ordered Hadamard (see Appendix~A).

\textbf{(c) ($\bigstar$) is not sufficient for algebraic identifiability.} There are settings where ($\bigstar$) holds yet exact separation of $(a,T)$ is impossible: (c1) any \emph{square} random design $N=K$ --- Proposition 1 gives ($\bigstar$), yet Theorem~A applies to every invertible realization, so a nonconstant $s\in S$ still maps $(\ell,T)$ to an exact alias; (c2) a uniform random re-ordering $\pi$ of the \emph{deterministic} ordered Hadamard record makes ($\bigstar$) hold without changing the coefficient multiset --- hence without adding any algebraic information. Quantitatively, Serfling's inequality for sampling without replacement~\cite{Serfling1974} gives, for the finite population $g_k=\log B_k$ with variance $\tau_g^2$ and range $R_g$,
\[
P\Big(\Big|W^{-1}\!\!\sum_{i=n}^{n+W-1} g_{\pi(i)}-\overline g\Big|\ge t\Big)\ \le\ 2\exp\!\Big[-\frac{W t^2}{2\tau_g^2+\tfrac23 R_g t}\Big],
\]
so a union bound over the $\le N$ windows yields ($\bigstar$)$(W,\epsilon,\delta)$ whenever $\log N\ll W$ (and $W$ remains below the drift coherence time). This device fails when some $g_k$ are undefined (zero buckets) or the population log-variance/range is uncontrolled, and it does \textbf{not} deliver the small carrier variance $\sigma_{\mathrm{LR}}^2\asymp1/K_{\mathrm{eff}}$ of genuinely random nonnegative patterns.
\end{namedthm}

\noindent\emph{Proof.} (a) Decompose exactly
\begin{align*}
\hat g_n-(\ell_n+m_T)\;=\;&\underbrace{\Big(W^{-1}\!\!\sum_{k\in W_n}\ell_k-\ell_n\Big)}_{\text{drift bias}}\\
&+\;\underbrace{\Big(W^{-1}\!\!\sum_{k\in W_n}\log B_k-m_T\Big)}_{\text{carrier deviation}} .
\end{align*}
The drift bias is $\le C_\beta L\,(W/N)^\beta\to0$ deterministically by (B3). Hence $\max_n|\hat g_n-(\ell_n+m_T)|\to0$ in probability iff the maximal carrier deviation does --- which is verbatim the asymptotic form of ($\bigstar$). Both implications are immediate from this identity. (b) Theorem~A$'$ (see Appendix~B) is proved without any distributional assumption on the carrier order; the support-zero criterion is Corollary~2 (see Appendix~A). (c1) Theorem~A (see Appendix~A) holds for \emph{any} invertible $M$, in particular a random draw, conditionally on the realization; ($\bigstar$), a statement about the temporal law, is unaffected. (c2) Serfling's bound plus the union over windows; the multiset $\{g_k\}$, hence every order-independent (algebraic) functional of the data, is invariant under $\pi$. \hfill$\blacksquare$

\paragraph{What this does not claim.} Part (a) is an equivalence for \emph{one estimator family} (windowed local averages of $\log R_n$; the same proof covers kernel weights of matching order). It is not a characterization of identifiability by all mechanisms --- (b) and (c) are precisely the two failure modes of that over-reading. Nor does ($\bigstar$) by itself bound the \emph{rate}: rates require the variance model of Theorem~B, and reconstruction quality after gain correction additionally requires the corrected linear inverse problem to be well posed (see Appendix~F).

\subsection{C.4 Theorem B --- statistical relative-gain identifiability}

\begin{namedthm}{Theorem B}
Assume:
\begin{itemize}
\item \textbf{(B1) Signal model.} $Y_n=\log R_n=\ell_n+m_T+z_n$, $n=1,\dots,N$, where $m_T=\mathbb{E}\log B_n$ is a time-constant scalar (finite by $B_n>0$ a.s.\ and $\mathbb{E}|\log B_n|<\infty$), and $z_n=\log B_n-m_T$.
\item \textbf{(B2) Carrier noise.} $\{z_n\}$ is centered and strictly stationary, and \emph{either} (B2a) $\alpha$-mixing with $\mathbb{E}|z_0|^{2+\eta}<\infty$ and $\sum_{h\ge1}\alpha(h)^{\eta/(2+\eta)}<\infty$, which by Davydov's covariance inequality~\cite{Davydov1968} gives absolutely summable autocovariances --- write $\gamma(h)=\operatorname{Cov}(z_0,z_h)$ and $\sigma_{\mathrm{abs}}^2:=\sum_{h\in\mathbb Z}|\gamma(h)|<\infty$ for the \textbf{summable-autocovariance constant}, the quantity entering every non-asymptotic variance bound below --- and hence also a finite \textbf{signed long-run variance} $\sigma_{\mathrm{LR}}^2:=\gamma(0)+2\sum_{h\ge1}\gamma(h)\ (\le\sigma_{\mathrm{abs}}^2)$, which is the \emph{asymptotic} variance of window means, $W\operatorname{Var}(W^{-1}\sum_{k\in W_n}z_k)\to\sigma_{\mathrm{LR}}^2$; the two constants play distinct roles and must not be conflated; \emph{or} (B2b), for high-probability bounds, $\beta$-mixing with $\beta(k)\le\beta_0\exp[-(k/b)^\kappa]$ and $\|z_n\|_{\psi_1}\le M$.
\item \textbf{(B3) Slow drift (normalized time).} $\ell\in S$ embeds in a normalized-time H\"older ball: $|\ell_j-\ell_k|\le L\,(|j-k|/N)^{\beta}$, $\beta\in(0,2]$ --- equivalently the per-frame constant $L_a=L\,N^{-\beta}$ of Appendix~D --- matched to the window order so the windowed bias obeys $|W^{-1}\sum_{k\in W_n}\ell_k-\ell_n|\le C_\beta L\,(W/N)^\beta$ (for $\beta\in(1,2]$ this requires centered symmetric windows so the first-order bias cancels).
\item \textbf{(B4) Window asymptotics.} Both $W=W_N\to\infty$ \emph{and} $W/N\to0$; for fixed $L$ the windowed drift bias $C_\beta L\,(W/N)^\beta\to0$ then holds automatically (for a triangular array with drift constant $L=L_N$ growing with $N$, require additionally $L\,(W/N)^\beta\to0$).
\end{itemize}
Then:

\textbf{(i) Pointwise consistency with rate.} Under (B2a), the \emph{non-asymptotic} bound $\mathbb{E}\big[\hat g_n-(\ell_n+m_T)\big]^2\le C\big(L^2(W/N)^{2\beta}+\sigma_{\mathrm{abs}}^2/W\big)\to0$ holds, the windowed bias being $O(L\,(W/N)^\beta)$; the \emph{asymptotic} variance of the window mean is $\sigma_{\mathrm{LR}}^2/W$ (signed long-run variance, (B2a)). Under (B2b), with block length $q=\lceil b\,[\log(4W/\delta)]^{1/\kappa}\rceil$, with probability $\ge1-\delta$,
\[
\begin{aligned}
\big|\hat g_n-(\ell_n+m_T)\big|\ \le\ &C_\beta L\,(W/N)^\beta\\
&+C\Big[\sqrt{\tfrac{\nu_q^2\,q\log(4/\delta)}{W}}+\tfrac{Mq\log(4/\delta)}{W}\Big],
\end{aligned}
\]
where $\nu_q^2$ is a block-variance proxy bounded by a multiple of $\sigma_{\mathrm{abs}}^2$; uniformity over $n\le N$ follows with $\delta\mapsto\delta/N$.

\textbf{(ii) Centered-gain recovery.} $\hat\ell_n^{\,c}:=\hat g_n-N^{-1}\sum_{m}\hat g_m$ satisfies $\mathbb{E}\big[\hat\ell_n^{\,c}-(\ell_n-\bar\ell)\big]^2\le 4C\big(L^2(W/N)^{2\beta}+\sigma_{\mathrm{abs}}^2/W\big)$; hence the relative gain $\ell_n-\bar\ell$ (equivalently $a_n$ up to one multiplicative constant) is identifiable and consistently estimable.

\textbf{(iii) Gauge.} The scalar $m_T+\bar\ell$ is \emph{not} identifiable: for every $c>0$, the substitution $a\mapsto ca$, $T\mapsto T/c$ leaves the law of $\{Y_n\}$ exactly invariant while shifting $\ell\mapsto\ell+\log c$, $m_T\mapsto m_T-\log c$. This is precisely the constant direction $\mathrm{span}\{1\}\subset S$ of Theorem~A, and it is the only ambiguity remaining.
\end{namedthm}

\noindent\emph{Proof.} (i) Bias--variance split as in Proposition 2(a). Bias: (B3). Variance: for a stationary sequence, $\operatorname{Var}\big(W^{-1}\sum_{k\in W_n}z_k\big)=W^{-1}\sum_{|h|<W}(1-|h|/W)\gamma(h)\le W^{-1}\sum_h|\gamma(h)|=\sigma_{\mathrm{abs}}^2/W$ --- the non-asymptotic bound, with the summable-autocovariance constant; as $W\to\infty$ the same expression converges to $\sigma_{\mathrm{LR}}^2/W\,(1+o(1))$ with the \emph{signed} long-run variance, which is the correct asymptotic constant and can be strictly smaller. Finiteness of both is (B2a)/Davydov~\cite{Davydov1968}; $(x+y)^2\le2x^2+2y^2$ combines bias and variance. For (B2b), partition the window into alternating blocks of length $q$; the exponential $\beta$-mixing coupling replaces blocks by independent copies up to total-variation cost $W\beta(q)/q\le\delta/2$ by the choice of $q$, and Bernstein's inequality for independent sub-exponential block sums (the Merlev\`ede--Peligrad--Rio inequality~\cite{MerlevedePeligradRio2011} with the standard blocking scheme) gives the display. (ii) $\hat\ell_n^{\,c}-(\ell_n-\bar\ell)=[\hat g_n-(\ell_n+m_T)]-N^{-1}\sum_m[\hat g_m-(\ell_m+m_T)]$; Jensen bounds the second term's second moment by the maximal pointwise one, and $(x+y)^2\le2x^2+2y^2$ gives the factor 4. (iii) Direct computation: $B_n(T/c)=B_n(T)/c$, so $Y_n$ is unchanged term by term; conversely, consistency in (ii) implies identifiability of the estimand $\ell_n-\bar\ell$ (two parameter configurations inducing the same law of $\{Y_n\}$ would force the same limit), so the ambiguity is exactly one dimensional. \hfill$\blacksquare$

\paragraph{What this does not claim.} Theorem~B is a \textbf{statistical, asymptotic} statement about \emph{relative}-gain recovery in a triangular array; it is not exact finite-sample algebraic identifiability of $(a,T)$ --- Theorem~A still holds verbatim for a square random design, and the two results are complementary, not in tension. It presupposes $B_n>0$ and integrable logs: in photon-counting records with empty bins the theorem must be replaced by its calibrated soft-log variant (see Appendix~D), where $\psi_\alpha(c)=\log(c+\alpha)$, the calibration curve $m_\alpha$ is inverted, and the same bias--variance structure holds with sensitivity $\kappa_\alpha$ and photon-starved variance $\sim1/(W\bar\lambda)$. The constants $\sigma_{\mathrm{abs}}^2$ and $\sigma_{\mathrm{LR}}^2$ are model inputs (they coincide for the i.i.d.\ carrier, where $\gamma(h)=0$ for $h\neq0$); identifying them with $(\sigma_I/\mu_I)^2/K_{\mathrm{eff}}$ invokes the flagged expansion of Proposition 1(iv). Finally, consistent gain correction does not by itself guarantee good reconstruction --- the corrected inverse problem must be independently well posed, and the propagation of the residual gain variance $v$ through a basis with leverage $B_L$, floor $C_0/N$, and the $\rho_{\mathrm{pair}}$ flip boundary is the subject of Appendix~F.

\section{Estimation rates, minimax lower bounds, and Fisher geometry}

This appendix proves the rate and optimality statements of Sec.~6: the high-probability windowed-estimator bound (Theorem~C's engine), its minimax optimality at the level of rates (Theorem~D), the quotient-Fisher geometry underlying the identifiability dichotomy, and the low-photon soft-log rate. Notation is locked to the main text: $R_n=a_nB_n$, $\ell=\log a\in S$ with $\dim S=p$ and $\mathrm{span}\{1\}\subset S$, $Y_n=\log R_n=\ell_n+m_T+z_n$ with $m_T=\mathbb{E}\log B_n$ a time-independent scalar, windows $W$, and the spread functionals $K_{\mathrm{eff}}$, $K_\infty$ as in Appendix~C. The H\"older smoothness order is written $\beta\in(0,2]$ with seminorm $L_a$ (the drift constant of Theorem~C); the symbol $\alpha$ is reserved for the soft-log offset $\psi_\alpha(c)=\log(c+\alpha)$.

Throughout, the \emph{discrete H\"older ball} $H_\beta(L_a)$ is defined operationally: $\ell$ belongs to $H_\beta(L_a)$ relative to a window family $\{W_n\}$ of order $\beta$ if the deterministic smoothing bias obeys
\begin{equation}
\Big|\,W^{-1}\!\!\sum_{k\in W_n}\ell_k-\ell_n\Big|\;\le\;C_\beta\,L_a\,W^{\beta}\qquad\text{for all }n. \tag{D.1}
\end{equation}
For $\beta\in(0,1]$ a one-sided moving average of $\ell$ with modulus $|\ell_k-\ell_n|\le L_a|k-n|^{\beta}$ satisfies (D.1). For $\beta\in(1,2]$, (D.1) requires a \emph{centered symmetric} window (or a kernel of order $\lfloor\beta\rfloor$), so that the first-order bias cancels and only curvature $L_2$ survives; a one-sided window does not attain $\beta>1$.

\subsection{D.1 Lemma D.1 (windowed gain-estimation rate under mixing)}

\paragraph{Statement.} Assume:
\begin{enumerate}
\item \emph{(Model)} $Y_n=\ell_n+m_T+z_n$, with $\{z_n\}$ strictly stationary and $\mathbb{E}z_n=0$.
\item \emph{(Mixing)} $\{z_n\}$ is $\beta$-mixing with $\beta(k)\le\beta_0\exp[-(k/b)^{\kappa}]$ for some $\beta_0,b>0$, $\kappa>0$.
\item \emph{(Tails)} $\|z_n\|_{\psi_1}\le M$ (sub-exponential).
\item \emph{(Smoothness)} $\ell\in H_\beta(L_a)$ in the sense (D.1), with the window order matched to $\beta$ as above.
\end{enumerate}
Let $\hat g_n=W^{-1}\sum_{k\in W_n}Y_k$ and set $q=\lceil b\,[\log(4W/\delta)]^{1/\kappa}\rceil$. Then for each fixed $n$ and $\delta\in(0,1)$, with probability at least $1-\delta$,
\begin{align*}
\big|\hat g_n-(\ell_n+m_T)\big|
&\;\le\;C_\beta L_a W^{\beta}\\
&\quad+\;C\Big[\sqrt{\tfrac{\nu_q^2\,q\,\log(4/\delta)}{W}}\;+\;\tfrac{M\,q\,\log(4/\delta)}{W}\Big], \tag{D.2}
\end{align*}
where $\nu_q^2=\operatorname{Var}\!\big(q^{-1/2}\sum_{i=1}^{q}z_i\big)$ is the block-variance proxy; when the autocovariances are absolutely summable, $\nu_q^2\le C'\sigma_{\mathrm{abs}}^2$ with the summable-autocovariance constant $\sigma_{\mathrm{abs}}^2=\sum_h|\gamma(h)|$ of Appendix~C (hypothesis (B2a)) --- the signed long-run variance $\sigma_{\mathrm{LR}}^2\le\sigma_{\mathrm{abs}}^2$ of Appendix~C is the asymptotic variance of window means and plays no role in these non-asymptotic bounds. Uniformity over $n=1,\dots,N$ follows by replacing $\delta$ with $\delta/N$ (a $\log N$ inflation inside the brackets).

\emph{(Expectation version, weaker hypotheses.)} If instead $\{z_n\}$ is only $\alpha$-mixing with $\sum_{h\ge1}\alpha(h)^{\eta/(2+\eta)}<\infty$ and $\|z_0\|_{2+\eta}<\infty$ for some $\eta>0$, then
\begin{equation}
\mathbb{E}\big[\hat g_n-(\ell_n+m_T)\big]^2\;\le\;C\,L_a^2W^{2\beta}+C\,\sigma_{\mathrm{abs}}^2/W. \tag{D.3}
\end{equation}

\paragraph{Proof.} \emph{Bias.} $\hat g_n-(\ell_n+m_T)=\big[W^{-1}\sum_k\ell_k-\ell_n\big]+W^{-1}\sum_k z_k$; the first bracket is bounded by $C_\beta L_aW^\beta$ by (D.1).

\emph{Fluctuation, high probability.} Partition the window $W_n$ into $m=\lceil W/q\rceil$ consecutive blocks of length $q$ and separate them into odd- and even-indexed families. By Berbee's coupling lemma~\cite{Berbee1979} (see also Rio~\cite{Rio2017}), on an enlarged probability space each family can be replaced by a sequence of \emph{independent} blocks with the same marginals, at total-variation cost at most $m\,\beta(q)\le \beta_0\,W\exp[-(q/b)^{\kappa}]$. The choice $q=\lceil b[\log(4W/\delta)]^{1/\kappa}\rceil$ makes this cost at most $\delta/2$ (adjusting $\beta_0$ into $C$). On the coupling event, each family is a sum of independent block sums $S_j=\sum_{i\in\text{block }j}z_i$ with $\|S_j\|_{\psi_1}\le CqM$ (triangle inequality for the $\psi_1$ norm) and $\operatorname{Var}(S_j)=q\,\nu_q^2$. Bernstein's inequality for sums of independent sub-exponential variables (e.g.\ Vershynin~\cite{Vershynin2018}, Thm.~2.8.1) applied to each family, with a union bound over the two families (whence the $4/\delta$), yields the bracketed term in (D.2) with the characteristic Gaussian-plus-linear tail: a $\sqrt{\nu_q^2 q\log(4/\delta)/W}$ moderate-deviation term and an $Mq\log(4/\delta)/W$ heavy-tail correction. When $\sum_h|\operatorname{Cov}(z_0,z_h)|<\infty$, $\nu_q^2=\operatorname{Var}(z_0)+2\sum_{h=1}^{q-1}(1-h/q)\operatorname{Cov}(z_0,z_h)\le C'\sigma_{\mathrm{abs}}^2$.

\emph{Fluctuation, expectation.} Davydov's covariance inequality~\cite{Davydov1968} gives $|\operatorname{Cov}(z_0,z_h)|\le 8\,\alpha(h)^{\eta/(2+\eta)}\|z_0\|_{2+\eta}^2$; the summability hypothesis makes the autocovariance series absolutely summable, so $\operatorname{Var}(W^{-1}\sum_{k\in W_n}z_k)\le \bar\sigma^2/W$ with $\bar\sigma^2\le C\sigma_{\mathrm{abs}}^2$. Combining with the squared bias via $(x+y)^2\le2x^2+2y^2$ gives (D.3). \hfill$\blacksquare$

\begin{namedthm}{Corollary D.2 (optimized rate = Theorem C's high-photon case)}
Minimizing the right side of (D.3) over $W$ gives
\begin{equation}
\begin{gathered}
W^{*}\asymp\big(\sigma_{\mathrm{abs}}^2/L_a^2\big)^{1/(2\beta+1)},\\
\mathrm{MSE}^{*}\asymp L_a^{2/(2\beta+1)}\,\sigma_{\mathrm{abs}}^{4\beta/(2\beta+1)},
\end{gathered}
\tag{D.4}
\end{equation}
valid when $W^{*}\in[1,N]$ and $W^{*}$ is shorter than the gain coherence time; otherwise the boundary window $W=1$ or $W=N$ applies. For $\beta=1$ with per-frame drift scale $L_1\asymp s\rho_{\mathrm{pair}}$ (adjacent-frame decorrelation convention --- \emph{not} $\rho_{\mathrm{bw}}$), this is $\mathrm{MSE}^{*}\asymp\sigma_{\mathrm{abs}}^{4/3}(s\rho_{\mathrm{pair}})^{2/3}$; for random carriers the carrier is white, $\sigma_{\mathrm{abs}}^2=\sigma_{\mathrm{LR}}^2\approx(\sigma_I/\mu_I)^2/K_{\mathrm{eff}}$ under the spikiness condition $K_\infty\to\infty$ of Prop.~1.
\end{namedthm}

\paragraph{What this does not claim.} The constants $C,C_\beta$ are not tracked and depend on the mixing constants $(\beta_0,b,\kappa)$ and $M$. The bound is for the \emph{relative} gain $\ell_n+m_T$; the scalar $m_T+\text{gauge}(\ell)$ is not estimable (Theorem~B). Nothing is claimed for $\beta>2$ (would need local polynomials), for non-stationary carriers, or about \emph{adaptive} window selection when $(\beta,L_a)$ are unknown. The step $\nu_q^2\lesssim\sigma_{\mathrm{abs}}^2$ requires summable autocovariances, which the stated mixing plus moment hypotheses supply.

\subsection{D.2 Theorem D (minimax lower bounds; rate level, constants open)}

\paragraph{Statement.} \emph{(H\"older / Assouad--two-point.)} Consider the i.i.d.\ Gaussian submodel $z_n\sim\mathcal N(0,\sigma^2)$ --- a member of the model class of Lemma D.1. Then for every $n$ interior to the design,
\begin{equation}
\inf_{\hat\ell}\ \sup_{\ell\in H_\beta(L_a)}\ \mathbb{E}\big(\hat\ell_n-\ell_n\big)^2\;\ge\;c_\beta\,L_a^{2/(2\beta+1)}\,\sigma^{4\beta/(2\beta+1)}, \tag{D.5}
\end{equation}
the infimum over all measurable estimators of the gauge-fixed log-gain. In normalized continuous time ($N$ samples on $[0,1]$, bandwidth $h=W/N$) the equivalent form is $\inf\sup\mathbb{E}(\hat\ell(t)-\ell(t))^2\ge c\,(\sigma^2/N)^{2\beta/(2\beta+1)}L^{2/(2\beta+1)}$.

\emph{(Sobolev / van Trees--Pinsker.)} In the Gaussian sequence formulation --- expand $\ell(t)=\sum_j\theta_j\varphi_j(t)$ in a Fourier basis and impose the Sobolev ellipsoid $\sum_j j^{2\beta}\theta_j^2\le L^2$ --- the integrated risk obeys
\begin{equation}
\begin{aligned}
&\inf_{\hat\ell}\ \sup_{\|\ell\|_{H^\beta}\le L}\ \frac1N\sum_n\mathbb{E}(\hat\ell_n-\ell_n)^2\\
&\qquad\ge\;c_\beta\,L^{2/(2\beta+1)}(\sigma^2/N)^{2\beta/(2\beta+1)},
\end{aligned}
\tag{D.6}
\end{equation}
with the gauge direction (the constant mode, confounded with $m_T$) excluded from the risk. Consequently the windowed estimator of Lemma D.1 is \textbf{minimax-rate optimal} over $H_\beta(L_a)$ for $\beta\le1$ with one-sided windows and for $\beta\le 2$ with centered symmetric windows. \textbf{Sharp minimax constants are open} (this is Theorem~D of the main text).

\paragraph{Proof.} \emph{Two-point (pointwise).} Fix $n_0$ and let $\phi$ be a fixed $C^\infty$ bump, $\mathrm{supp}\,\phi\subset[-1,1]$, $\phi(0)=1$, $\|\phi\|_{H_\beta}\le1$. Take the hypotheses $\ell^{(0)}\equiv0$ and $\ell^{(1)}_k=L_a W_0^{\beta}\phi\big((k-n_0)/W_0\big)$; both lie in $H_\beta(L_a)$ (up to an absolute constant absorbed into $c_\beta$). Since the noise is i.i.d.\ Gaussian, the Kullback--Leibler divergence between the two data laws is $\mathrm{KL}=\tfrac{1}{2\sigma^2}\sum_k(\ell^{(1)}_k)^2\le C\,L_a^2W_0^{2\beta+1}/\sigma^2$. Choosing $W_0\asymp(\sigma^2/L_a^2)^{1/(2\beta+1)}$ makes $\mathrm{KL}\le c<\infty$, and the two hypotheses are separated at $n_0$ by $|\ell^{(1)}_{n_0}-\ell^{(0)}_{n_0}|=L_aW_0^{\beta}$. Le Cam's two-point lemma~\cite{Tsybakov2009} then gives (D.5). The nuisance scalar $m_T$ only enlarges the model, so the bound holds a fortiori; the bumps can be taken mean-centered over the record so that the two hypotheses differ in the gauge-invariant component $\ell-\bar\ell$ itself. For the integrated version, tile $[1,N]$ with $\sim N/W_0$ disjoint bumps carrying independent signs and apply Assouad's lemma~\cite{Tsybakov2009}.

\emph{van Trees / Pinsker.} In the sequence model $y_j=\theta_j+ (\sigma/\sqrt N)\epsilon_j$ (exact for equispaced discrete-Fourier regression with white Gaussian noise), place independent priors $\theta_j\sim\mathcal N(0,\tau_j^2)$ truncated inside the ellipsoid. The van Trees inequality~\cite{GillLevit1995} bounds the Bayes risk of each coordinate by the reciprocal of (data Fisher $+$ prior Fisher): $\text{BayesRisk}\ge\sum_j(N/\sigma^2+\tau_j^{-2})^{-1}$, the constant mode omitted (gauge). Maximizing the right side over prior spectra supported in the ellipsoid is Pinsker's calculation~\cite{Pinsker1980} and yields the scaling (D.6). Since minimax risk dominates Bayes risk, (D.6) follows. \hfill$\blacksquare$

\paragraph{What this does not claim.} (i) Constants: neither $c_\beta$ nor a Pinsker-type sharp constant is claimed for our model; the exact Pinsker constant transfers only where the acquisition is genuinely the Gaussian sequence model --- for mixing, non-Gaussian carriers the statement is at rate level only (a Brown--Low-type asymptotic-equivalence argument~\cite{BrownLow1996} would be needed for more, and is not attempted). (ii) The lower bound is proven inside the i.i.d.\ Gaussian \emph{submodel}; this suffices for minimaxity of the full class (sup over the class dominates sup over the submodel, while the upper bound (D.3) holds over the full class), but no lower bound is claimed \emph{within} every fixed mixing law. (iii) No adaptivity claim: rate optimality is for known $(\beta,L_a)$.

\subsection{D.3 Proposition D.4 (quotient Fisher geometry: singularity \texorpdfstring{$\Leftrightarrow$}{<=>} the Theorem-A ambiguity)}

\paragraph{Statement.} \emph{(a) Linear-Gaussian window model.} Let $\ell=U\theta$ with $U\in\mathbb R^{N\times p'}$, and observe $Y=U\theta+m\mathbf 1+z$, $z\sim\mathcal N(0,\Sigma)$, $\Sigma\succ0$ known, $m\in\mathbb R$ the unknown carrier log-mean (confounded with global gain scale). The Fisher information for $\theta$ after eliminating the nuisance $m$ is
\begin{equation}
\begin{gathered}
I_\theta\;=\;U^{\top}\Sigma^{-1/2}P_\perp\Sigma^{-1/2}U,\\
P_\perp=I-\frac{\Sigma^{-1/2}\mathbf 1\mathbf 1^{\top}\Sigma^{-1/2}}{\mathbf 1^{\top}\Sigma^{-1}\mathbf 1}.
\end{gathered}
\tag{D.7}
\end{equation}
For a linear functional $L\theta$: if $\mathrm{row}(L)\subseteq\mathrm{range}(I_\theta)$, every unbiased estimator obeys $\operatorname{Cov}(L\hat\theta)\succeq L\,I_\theta^{+}L^{\top}$ (Moore--Penrose CRB for singular information~\cite{RaoMitra1971,StoicaMarzetta2001}); if $L$ has a nonzero component along $\ker I_\theta$, \textbf{no unbiased estimator of $L\theta$ exists and the CRB is infinite} --- the functional is not locally estimable.

\emph{(b) Quotient statement.} Let $\Theta$ be the parameter manifold, $G$ the global-scale group $(a,T)\mapsto(ca,T/c)$, and $P_\theta$ a dominated model differentiable in quadratic mean. Define the quotient score operator $S_\theta:T_{[\theta]}(\Theta/G)\to L_0^2(P_\theta)$ and $I_\theta=S_\theta^{*}S_\theta$. Then: $I_\theta\succ0$ on $T_{[\theta]}(\Theta/G)$ \textbf{iff} $[\theta]\mapsto P_\theta$ is an immersion, i.e.\ iff the model is \emph{locally differentially identifiable} modulo scale. Any $C^1$ ambiguity curve $[\theta(t)]$ with $P_{\theta(t)}=P_{\theta(0)}$ satisfies $S_\theta\theta'(0)=0$, hence produces a Fisher null direction. For the Theorem-A ambiguity this implication is \emph{exact}: the explicit data-preserving path $\ell'=\ell+ts$, $T'=M^{-1}(c\odot e^{-ts})$, $s\in S\setminus\mathrm{span}\{1\}$ --- the \emph{exact collision curve} $\ell_u=\ell+u$, $T_u=M^{-1}D_{e^{-u}}MT$ in the notation of Appendix~A --- is a genuine ambiguity curve for a square invertible design, so the Fisher matrix in the full $(\ell,T)$ parameterization is singular along its tangent, and the CRB is $+\infty$ for every functional that varies along it. For a \emph{tall} full-column-rank design the exact collision exists for a given $u\in S$ iff the range condition $(I-P_M)D_{e^{-u}}MT=0$ holds ($P_M$ the column-space projector); under the nonvanishing-carrier assumption $T_u\propto T$ only when $u$ is constant, and the quotient tangent null space is $\mathcal N_{\ell,T}=\{(u,h):D_{MT}u+Mh=0\}\supseteq\operatorname{span}\{(\mathbf1,-T)\}$, with local differential identifiability modulo scale exactly when equality holds.

\emph{(c) Local-window corollary.} For a window over which $\ell$ is approximately constant, the Gaussian CRB reduces to $\operatorname{Var}(\hat\ell_n)\ge\sigma_z^2/W_{\mathrm{eff}}$ with $W_{\mathrm{eff}}=W/(1+2\sum_h\rho_z(h))$ under weak dependence (equivalently, $\sigma_z^2$ replaced by the signed long-run variance $\sigma_{\mathrm{LR}}^2$, which is the exact asymptotic constant here) --- the Fisher-side counterpart of the $\sigma_{\mathrm{abs}}^2/W$ upper bound in (D.3).

\paragraph{Proof.} \emph{(a)} The joint Fisher matrix of $(\theta,m)$ in the Gaussian model is the block matrix $\begin{pmatrix}U^{\top}\Sigma^{-1}U & U^{\top}\Sigma^{-1}\mathbf 1\\ \mathbf 1^{\top}\Sigma^{-1}U & \mathbf 1^{\top}\Sigma^{-1}\mathbf 1\end{pmatrix}$. Nuisance elimination is the Schur complement with respect to the $m$-block~\cite{LehmannCasella1998}, which is exactly (D.7). The pseudoinverse CRB for singular FIM is Stoica--Marzetta~\cite{StoicaMarzetta2001}. The infinite-CRB claim is cleanest without CRB machinery: a direction $u\in\ker I_\theta$ that is tangent to an exact ambiguity curve leaves the data law invariant, so any estimator has \emph{identical} distribution under $\theta$ and $\theta+tu$; unbiasedness for a functional with $Lu\neq0$ would force $L\theta=L(\theta+tu)$, a contradiction --- hence no unbiased finite-variance estimator exists.

\emph{(b)} Differentiability in quadratic mean makes $S_\theta$ well defined and $I_\theta=S_\theta^*S_\theta$ the pullback metric \cite[Ch.~7]{vanderVaart1998}. $I_\theta\succ0$ iff $S_\theta$ is injective iff the differential of $[\theta]\mapsto\sqrt{dP_\theta}$ is injective, i.e.\ an immersion. The ambiguity-curve implication is the chain rule: $t\mapsto P_{\theta(t)}$ constant $\Rightarrow$ its $L^2$ derivative $S_\theta\theta'(0)$ vanishes. For Theorem~A the curve is explicit (see Appendix~A), so no genericity is needed. \hfill$\blacksquare$

\paragraph{What this does not claim.} The converse direction --- Fisher singularity implying an actual \emph{ambiguity manifold} --- requires constant-rank (or analytic) regularity so that the rank defect integrates (rank theorem); Fisher singularity by itself only defeats \emph{local differential} identifiability. Nonsingular Fisher does not exclude isolated \emph{global} aliases. Part (a) assumes Gaussian noise with known $\Sigma$; part (c) is a Gaussian-model computation used as an order-of-magnitude floor, not a bound proven for the actual log-carrier distribution. No efficiency claim is made for the windowed estimator beyond rate (constants open, Theorem~D).

\subsection{D.4 Theorem C (low-photon calibrated soft-log rate and the conditional \texorpdfstring{$1/(W\bar\lambda)$}{1/(W lambda-bar)} law)}

\paragraph{Statement (exact finite-window identity with known carriers, weighted form).} Let counts follow the calibrated model $C_n\mid \ell_n\sim\mathrm{Pois}(\lambda_n)$, $\lambda_n=\Lambda_0e^{\ell_n}b_n+d$, with $d\ge0$ dark counts and $\{b_n\}$ the \emph{known} per-frame carrier intensities ($b_n>0$; in the Fig.~5 protocol these are \emph{oracle} values supplied from the simulation truth --- a flat-field-calibrated benchmark, not programmed design intensities --- and $d=0$ --- Remark D.4.1; the unknown-random-carrier and known-law variants are Remarks D.4.3 and D.4.5). Fix $\alpha>0$, let $\psi_\alpha(c)=\log(c+\alpha)$, and let $m_\alpha(\lambda)=\mathbb{E}[\psi_\alpha(C)]$ for $C\sim\mathrm{Pois}(\lambda)$ --- the homogeneous calibration curve, strictly increasing. For a fixed window $n$, let $\mathcal I_n$ be the set of \emph{distinct} frame indices used by the estimator and let $w_{nk}\ge0$, $\sum_{k\in\mathcal I_n}w_{nk}=1$, be its weights (for the implementation's windows, $\mathcal I_n=W_n$ and $w_{nk}=1/W$). If a padding convention repeats a frame, all repeated occurrences are first aggregated into its single weight $w_{nk}$. Put
\[
q_k(t)=m_\alpha(\Lambda_0e^t b_k+d),
\]
\[
v_k(t)=\operatorname{Var}\!\left[
\psi_\alpha(\operatorname{Pois}(\Lambda_0e^t b_k+d))\right],
\]
\[
M_n(t)=\sum_{k\in\mathcal I_n}w_{nk}q_k(t),
\qquad
y_n=\sum_{k\in\mathcal I_n}w_{nk}\psi_\alpha(C_k),
\]
and define
\[
B_n=\sum_{k\in\mathcal I_n}w_{nk}
\{q_k(\ell_k)-q_k(\ell_n)\},
\qquad
V_n=\sum_{k\in\mathcal I_n}w_{nk}^2v_k(\ell_k).
\]
Fix a compact operating range $\Theta=[\theta_{\min},\theta_{\max}]$ and let
\[
\underline\kappa_n=\inf_{t\in\Theta}M_n'(t)>0,
\qquad
D_\Theta=\theta_{\max}-\theta_{\min}.
\]
In the implementation $\Theta$ is the \emph{calibration-safe interval}: if the calibration curve is certified only for $\lambda\in[\lambda_-,\lambda_+]$ and $d<\lambda_-$, the largest safe interval before intersection with the physical parameter range is
\[
\left[
\max_{k\in\mathcal I_n}\log\frac{\lambda_--d}{\Lambda_0b_k},
\quad
\min_{k\in\mathcal I_n}\log\frac{\lambda_+-d}{\Lambda_0b_k}
\right],
\]
and \textbf{nonemptiness of this interval is an assumption of the theorem} --- a clamp cannot repair an empty calibration-safe interval. The estimator is the \emph{endpoint-clamped generalized inverse} applied to the window statistic:
\begin{equation}
\hat\theta_n\;=\;G_n(y_n),
\qquad
G_n(y)\;=\;M_n^{-1}\!\big[\Pi_{M_n(\Theta)}\,y\big], \tag{D.8a}
\end{equation}
where $\Pi_{M_n(\Theta)}$ projects onto the compact interval $M_n(\Theta)$ before inversion, so $G_n$ returns the corresponding endpoint of $\Theta$ when $y_n$ falls outside $M_n(\Theta)$ (an all-zero window has $y_n=\log\alpha$ and returns $\theta_{\min}$; without the clamp the inverse would diverge to $-\infty$ there --- the clamp is part of the estimator, not an implementation afterthought). Assume:
\begin{enumerate}
\item \emph{(Operating range; per-frame sensitivities and noise)} for every $k$ and every $\theta\in\Theta$, $\lambda_k(\theta)=\Lambda_0e^{\theta}b_k+d\in[\lambda_-,\lambda_+]\subset(0,\infty)$; consequently the per-frame sensitivity bound $\overline\kappa_k=\sup_{t\in\Theta}q_k'(t)$ and the per-frame noise $v_k(t)$ are finite, with $v_k(t)\le\bar\sigma_\alpha^2:=\sup_{\lambda\in[\lambda_-,\lambda_+]}\operatorname{Var}[\psi_\alpha(\mathrm{Pois}(\lambda))]<\infty$, $\kappa_{\max}:=\sup_k\overline\kappa_k<\infty$, and $\underline\kappa_n>0$;
\item \emph{(Smoothness)} $\ell\in H_\beta(L_a)$ with the \emph{pointwise} modulus $|\ell_k-\ell_n|\le L_a|k-n|^{\beta}$ and $\beta\in(0,1]$ (the paper's operative case is $\beta=1$; the precise obstruction above one and its partial repair are Remark D.4.2);
\item \emph{(Membership)} $\ell_n\in\Theta$ for all $n$. (A positive margin $\Delta:=\min_n\operatorname{dist}(\ell_n,\partial\Theta)>0$ enters only the clamp-\emph{probability} diagnostic of Remark D.4.4; the risk bound below does not use it.)
\end{enumerate}
Then, conditionally on the gain path and known carriers,
\[
\operatorname{Var}\!\left[
\sum_{k\in\mathcal I_n}w_{nk}
\{\psi_\alpha(C_k)-q_k(\ell_k)\}\right]
=V_n
\]
\emph{exactly}, and no mixing or stationarity assumption is required. Then, conditionally on the gain path and the known carriers,
\begin{equation}
\mathbb E\big[(\hat\theta_n-\ell_n)^2\mid\ell,b\big]
\le
\frac{B_n^2+V_n}{\underline\kappa_n^2}.
\tag{D.8}
\end{equation}
There is \textbf{no separate clamp-event risk term}: the endpoint-clamped generalized inverse $G_n$ is globally $\underline\kappa_n^{-1}$-Lipschitz on all of $\mathbb R$ (proof below), so all-zero and out-of-range windows are already covered by (D.8); the clamp probability survives only as a diagnostic (Remark D.4.4). This is a finite-window statement: it requires neither stationarity nor mixing of the known carrier sequence.
Moreover, the pointwise modulus gives
\[
|B_n|
\le
L_a\sum_{k\in\mathcal I_n}
w_{nk}\overline\kappa_k|k-n|^\beta.
\]
For an unpadded average of $W$ distinct frames, $w_{nk}=1/W$ and hence
\[
V_n=\frac{\bar v_{n,\alpha}}{W},
\qquad
\bar v_{n,\alpha}
=\frac1W\sum_{k\in W_n}v_k(\ell_k)
\le\bar\sigma_\alpha^2,
\]
while $|B_n|\le C\,\kappa_{\max}L_aW^{\beta}$, so (D.8) specializes to
\[
\mathbb{E}\big(\hat\theta_n-\ell_n\big)^2\;\le\;C\Big(\tfrac{\kappa_{\max}}{\underline\kappa_n}\Big)^{2}L_a^2W^{2\beta}
+\frac{\bar\sigma_\alpha^2}{\underline\kappa_n^2\,W}
\]
(no clamp term: the global Lipschitz property of $G_n$ already covers clamped windows). Optimizing over $W$,
\begin{equation}
\mathrm{MSE}_\alpha^{*}\;\le\;C\,\underline\kappa^{-2}\,\big[\kappa_{\max}L_a\big]^{2/(2\beta+1)}\,\bar\sigma_\alpha^{4\beta/(2\beta+1)} \tag{D.9}
\end{equation}
--- Theorem~C of the main text, with $\underline\kappa:=\inf_n\underline\kappa_n$. \textbf{The loss throughout is the log-domain squared error}: $\hat\theta_n-\ell_n$ is an increment of log-gain (equivalently log-intensity), the record-mean centering fixes the additive gauge, and this is the metric in which the Fisher comparison below is stated (and which Fig.~5 reports). In the photon-starved regime the variance-to-sensitivity ratio obeys the \emph{general} window law
\[
\frac{V_n}{M_n'(\ell_n)^2}
=\frac{\sum_kw_{nk}^2\lambda_k(\ell_k)}{\big[\sum_kw_{nk}\lambda_k(\ell_n)\big]^2}\{1+O_\alpha(\varepsilon)\}
\]
(low-photon corollary below), with the \emph{actual} intensities $\lambda_k(\ell_k)$ in the numerator and the \emph{anchor} intensities $\lambda_k(\ell_n)$ in the sensitivity denominator; it simplifies to $\{1+O_\alpha(\varepsilon)\}/(W\bar\lambda_n)$ only for distinct equal-weight frames under \emph{local gain flatness} (see the counterexample below), and in that flat regime it is minimax-\emph{order} sharp \emph{in this log-domain loss at $d=0$} by the gauge-invariant lower bound following the proof. Zeros only reduce Fisher information; with the clamp (D.8a) they never make the estimator diverge.

\paragraph{Proof.} \emph{Sensitivity and variance of the soft-log.} The Poisson derivative identity $\frac{d}{d\lambda}\mathbb{E}f(C)=\mathbb{E}[f(C{+}1)-f(C)]$ (differentiate $\sum_c e^{-\lambda}\lambda^c f(c)/c!$ term by term) gives
\[
\kappa_\alpha(\lambda):=\frac{d}{d\log\lambda}\mathbb{E}\log(C+\alpha)=\lambda\,\mathbb{E}\log\frac{C+1+\alpha}{C+\alpha},
\]
with the expansions $\kappa_\alpha(\lambda)=1+O(1/\lambda)$ for $\lambda\gg1$ and $\kappa_\alpha(\lambda)=c_\alpha\lambda+O(\lambda^2)$ for $\lambda\ll1$, where $c_\alpha=\log(1+1/\alpha)$. Since $q_k'(t)=\kappa_\alpha(\lambda_k(t))\cdot\Lambda_0e^{t}b_k/\lambda_k(t)$ ($=\kappa_\alpha(\lambda_k(t))$ at $d=0$), hypothesis 1 holds on any compact operating range. No uniform identification $\underline\kappa_n\asymp c_\alpha\bar\lambda_n$ is claimed: under unrestricted heteroscedastic carriers the infimum of $M_n'$ over $\Theta$ can sit far below $c_\alpha\bar\lambda_n$ (two-frame example $\lambda_1=\varepsilon^2$, $\lambda_2=2\varepsilon-\varepsilon^2$), so $\underline\kappa_n$-based displays are conservative bounds only; the correct low-photon law is the window-average corollary below. The Poisson Poincar\'e inequality~\cite{Klaassen1985,BobkovLedoux1998}, $\operatorname{Var}f(C)\le\lambda\,\mathbb{E}[(f(C{+}1)-f(C))^2]$, yields $v_k=O(1/\lambda)$ at high $\lambda$ and $v_k=c_\alpha^2\lambda+O(\lambda^2)$ at low $\lambda$, whence $\bar\sigma_\alpha^2<\infty$ on $[\lambda_-,\lambda_+]$.

\emph{Risk bound.} Write
\[
\begin{aligned}
y_n-M_n(\ell_n)&=B_n+Z_n,\\
Z_n&=\sum_{k\in\mathcal I_n}w_{nk}
\{\psi_\alpha(C_k)-q_k(\ell_k)\}.
\end{aligned}
\]
The variables in $Z_n$ are independent and centered conditionally on the gain path and known carriers (Poisson counts are independent across frames given their intensities --- the carrier is known, so no carrier randomness remains; this \emph{conditional count independence} is an assumption of the theorem), so $\mathbb EZ_n=0$ and $\operatorname{Var}(Z_n)=V_n$ \emph{exactly} --- a finite-window identity requiring no mixing and no long-run variance.

\emph{Global Lipschitz property of the clamped inverse.} $M_n$ is continuous and strictly increasing on the compact $\Theta$ with $M_n'\ge\underline\kappa_n>0$, so its ordinary inverse is $\underline\kappa_n^{-1}$-Lipschitz on $M_n(\Theta)$. Projection onto the interval $M_n(\Theta)$ is nonexpansive, hence the composition $G_n=M_n^{-1}\circ\Pi_{M_n(\Theta)}$ satisfies
\[
|G_n(y_2)-G_n(y_1)|\le\frac{|y_2-y_1|}{\underline\kappa_n}
\qquad\text{for \emph{all} }y_1,y_2\in\mathbb R .
\]
Since $\ell_n\in\Theta$ gives $G_n\{M_n(\ell_n)\}=\ell_n$,
\[
|\hat\theta_n-\ell_n|=|G_n(y_n)-G_n\{M_n(\ell_n)\}|
\le\frac{|B_n+Z_n|}{\underline\kappa_n}
\]
with no case split on whether $y_n$ lands inside $M_n(\Theta)$: taking conditional second moments and using $\mathbb E[(B_n+Z_n)^2\mid\ell,b]=B_n^2+V_n$ gives (D.8) directly. The bound on $B_n$ follows by applying the mean value theorem to each $q_k$ and then the pointwise modulus; for equal weights the triangle inequality gives $|B_n|\le C\kappa_{\max}L_aW^{\beta}$ (valid for the pointwise modulus with $\beta\le1$; no cancellation across the window is invoked --- cf.\ Remark D.4.2), and the standard bias--variance optimization gives (D.9). \hfill$\blacksquare$

\paragraph{Approximate calibration and numerical inversion.} Let $\widetilde M_n$ be a continuous strictly increasing approximation to $M_n$ on the same interval $\Theta$ with
$\sup_{t\in\Theta}|\widetilde M_n(t)-M_n(t)|\le\varepsilon_{\mathrm{cal},n}$,
let $\widetilde G_n$ be its clamped generalized inverse (same endpoint convention), and suppose the numerical solver returns $\hat\theta_n^{\mathrm{num}}$ with $|\hat\theta_n^{\mathrm{num}}-\widetilde G_n(y_n)|\le\varepsilon_{\mathrm{bis},n}$, where $\varepsilon_{\mathrm{bis},n}$ is measured in the \emph{log-parameter} domain. Inverse bracketing under the shared endpoint convention gives $G_n(y-\varepsilon_{\mathrm{cal},n})\le\widetilde G_n(y)\le G_n(y+\varepsilon_{\mathrm{cal},n})$, hence $\sup_y|\widetilde G_n(y)-G_n(y)|\le\varepsilon_{\mathrm{cal},n}/\underline\kappa_n$, and the conservative implementation bound is the RMSE (Minkowski) form
\begin{equation}
\big\{\mathbb E[(\hat\theta_n^{\mathrm{num}}-\ell_n)^2\mid\ell,b]\big\}^{1/2}
\le\frac{\sqrt{B_n^2+V_n}}{\underline\kappa_n}
+\frac{\varepsilon_{\mathrm{cal},n}}{\underline\kappa_n}
+\varepsilon_{\mathrm{bis},n}.
\tag{D.8b}
\end{equation}
The quantities with the same (log-parameter) units are $\varepsilon_{\mathrm{cal},n}/\underline\kappa_n$ and $\varepsilon_{\mathrm{bis},n}$; a response-domain calibration error and a log-parameter bisection error must \emph{not} be added before division by the sensitivity. If tabulation and interpolation have separately certified forward-error budgets, replace $\varepsilon_{\mathrm{cal},n}$ by their sum. For $J$ midpoint bisections on a log-parameter interval of diameter $D_n$, the numerical error is certified by $\varepsilon_{\mathrm{bis},n}\le D_n/2^{J+1}$.

\paragraph{A theorem-level interpolation certificate.} A midpoint probe of the interpolation error is an empirical diagnostic, not by itself a uniform bound over the calibration interval. Put $x=\log\lambda$ and $f_\alpha(x)=m_\alpha(e^x)$. For a log-intensity grid of maximum spacing $\Delta_x$, piecewise-linear interpolation from exact node values obeys
\begin{equation}
\|f_\alpha-I_{\Delta_x}f_\alpha\|_\infty
\le\frac{\Delta_x^2}{8}\,\|f_\alpha''\|_\infty ,
\tag{D.8c}
\end{equation}
and with $\Delta g(c)=g(c{+}1)-g(c)$ the Poisson difference identities give
\[
f_\alpha''(x)=\lambda\,\mathbb E\Delta\psi_\alpha(C)+\lambda^2\,\mathbb E\Delta^2\psi_\alpha(C),
\qquad\lambda=e^x .
\]
A certified bound $K_{\alpha,2}\ge\sup_{\lambda\in[\lambda_-,\lambda_+]}|f_\alpha''|$ together with a certified node-evaluation error $\varepsilon_{\mathrm{tab}}$ yields $\varepsilon_{\mathrm{cal}}\le\varepsilon_{\mathrm{tab}}+K_{\alpha,2}\Delta_x^2/8$; the expectations and Poisson tails may be enclosed analytically or by interval arithmetic. \textbf{Until such an enclosure is generated, the measured midpoint maximum of Remark D.4.1 must be labeled empirical}, and the manuscript so labels it.

\paragraph{Low-photon corollary (the general leading term).} Let $d=0$ and define the \emph{anchor} and \emph{actual} per-frame intensities
\[
\lambda_{k,n}^{0}=\Lambda_0e^{\ell_n}b_k=\lambda_k(\ell_n),
\qquad
\lambda_k^\star=\Lambda_0e^{\ell_k}b_k=\lambda_k(\ell_k),
\]
\[
A_n=\sum_kw_{nk}\lambda_{k,n}^{0},\qquad
D_n=\sum_kw_{nk}^2\lambda_k^\star .
\]
Fix $\alpha>0$, put $c_\alpha=\log(1+1/\alpha)$, and suppose $\max_{k\in\mathcal I_n}\{\lambda_{k,n}^{0},\lambda_k^\star\}\le\varepsilon$. The fixed-$\alpha$ Poisson expansions of the proof give, uniformly over arbitrary positive carrier heterogeneity,
\[
M_n'(\ell_n)=c_\alpha A_n\{1+O_\alpha(\varepsilon)\},\qquad
V_n=c_\alpha^2D_n\{1+O_\alpha(\varepsilon)\},
\]
and therefore
\begin{equation}
\begin{aligned}
\frac{V_n}{M_n'(\ell_n)^2}
&=\frac{D_n}{A_n^2}\{1+O_\alpha(\varepsilon)\}\\
&=\frac{\sum_kw_{nk}^2\lambda_k(\ell_k)}
      {\big[\sum_kw_{nk}\lambda_k(\ell_n)\big]^2}\{1+O_\alpha(\varepsilon)\}.
\end{aligned}
\tag{D.10}
\end{equation}
For $m=|\mathcal I_n|$ equal-weight \emph{distinct} frames, writing $I_{n,\mathrm{ref}}=\sum_k\lambda_{k,n}^{0}$ and $I_{n,\mathrm{act}}=\sum_k\lambda_k^\star$, (D.10) reads $I_{n,\mathrm{act}}/I_{n,\mathrm{ref}}^2\cdot\{1+O_\alpha(\varepsilon)\}$, and with $h_{\ell,n}=\max_k|\ell_k-\ell_n|$,
\[
e^{-h_{\ell,n}}\le I_{n,\mathrm{act}}/I_{n,\mathrm{ref}}\le e^{h_{\ell,n}} .
\]
Only under \emph{local gain flatness} $h_{\ell,n}=o(1)$ may the leading term be simplified to $\{1+O_\alpha(\varepsilon+h_{\ell,n})\}/I_{n,\mathrm{ref}}$; when the gain is constant in the window, $I_{n,\mathrm{act}}=I_{n,\mathrm{ref}}=m\bar\lambda_n$ and only then does (D.10) become $\{1+O_\alpha(\varepsilon)\}/(m\bar\lambda_n)$ --- the $1/(W\bar\lambda)$ statement quoted in the main text. Carrier comparability is \emph{not} needed; local gain flatness \emph{is}. These expansions are stated for fixed $\alpha$; $c_\alpha\to\infty$ as $\alpha\to0$, so the fixed-$\alpha$ small-$\lambda$ statement must not be read as a triangular $\alpha\to0$ asymptotic.

\paragraph{Counterexample to the unqualified $1/(W\bar\lambda)$ claim.} Take two equal-weight frames with equal carriers, target gain $\ell_n=0$, the other frame's gain $\ell=\log2$, and reference intensity $\varepsilon$. Then $I_{n,\mathrm{ref}}=2\varepsilon$ but $I_{n,\mathrm{act}}=3\varepsilon$, so (D.10) gives $3/(4\varepsilon)$ whereas the unqualified formula gives $1/(2\varepsilon)$: the ratio is $3/2$ and does \emph{not} converge to one as $\varepsilon\downarrow0$. The unconditional window-average law asserted in earlier drafts is therefore false as stated and is withdrawn; (D.10) is the correct general statement.

\paragraph{From a variance factor to estimator asymptotics.} Display (D.10) is an algebraic expansion of the local variance-to-sensitivity factor. To identify it with the estimator's leading MSE one additionally needs a local variance regime --- $I_{n,\mathrm{act}}\to\infty$, $|B_n|=o(\sqrt{V_n})$, a Lindeberg condition, negligible clamp probability, and inverse-map linearization such as $\sup_{\Theta}|M_n''|\sqrt{V_n}/M_n'(\ell_n)^2\to0$. For bounded total information the compact clamp makes the risk saturate; one cannot globally identify the MSE with $1/I_n$.

\paragraph{Gauge-invariant low-photon lower bound.}
The loss corresponding to the unidentifiable global scale is
\[
\begin{aligned}
d_{\mathrm{quot}}^2(\hat g,g)
&=\inf_{a\in\mathbb R}\frac1N\|\log\hat g-\log g-a\mathbf1\|_2^2\\
&=\frac1N\|P_N(\log\hat g-\log g)\|_2^2,
\quad
P_N=I-\tfrac1N\mathbf1\mathbf1^\top.
\end{aligned}
\]
This is the centered-log-gain loss reported in Fig.~5. Multiplying
either gain profile by an arbitrary positive scalar leaves the loss
unchanged.

\emph{Pointwise lower bound (Le Cam two-point).} Let $0<\beta\le1$ and, conditionally on known positive carriers $q_1,\dots,q_N$, observe independent counts $C_i\sim\operatorname{Pois}(q_ie^{\ell_i})$. Estimate the quotient coordinate $P_N\ell$. Fix an interior index $n_0$, a baseline $\ell^{(0)}$ lying strictly inside both the amplitude and H\"older constraints (e.g.\ $\|\ell^{(0)}\|_\infty\le a_0/2$ and seminorm at most $L_a/2$), and a window $A$ of order $W$ containing $n_0$. Choose a contrast $\psi^{(W)}$ with
\[
\operatorname{supp}\psi^{(W)}\subseteq A,\quad
\mathbf1^\top\psi^{(W)}=0,\quad
\|\psi^{(W)}\|_\infty\le1,
\]
\[
\psi^{(W)}_{n_0}=1,\quad
\|\psi^{(W)}\|_2^2\ge cW,
\]
\[
|\psi^{(W)}_i-\psi^{(W)}_j|\le CW^{-\beta}|i-j|^\beta .
\]
Put $\lambda_i^{(0)}=q_ie^{\ell_i^{(0)}}$ and $I_A(\psi)=\sum_{i\in A}\lambda_i^{(0)}\psi_i^2\le W\bar\lambda_A$ with $\bar\lambda_A=W^{-1}\sum_{i\in A}\lambda_i^{(0)}$. For amplitude
$a\le c_0\min\{a_0,\,L_aW^\beta,\,I_A(\psi)^{-1/2}\}$,
both $\ell^{(0)}$ and $\ell^{(1)}=\ell^{(0)}+a\psi^{(W)}$ lie in a rescaled $H_\beta(L_a)$ ball; the Poisson KL obeys
\[
\operatorname{KL}(P_0,P_1)
=\sum_i\lambda_i^{(0)}\{e^{a\psi_i}-1-a\psi_i\}
\le\tfrac{e^{a_0}}{2}a^2I_A(\psi),
\]
kept below a numerical constant by the amplitude choice. Because $\mathbf1^\top\psi=0$, quotient projection does not reduce the separation at $n_0$, and Pinsker followed by Le Cam's two-point lemma~\cite{Tsybakov2009} gives
\begin{equation}
R_{n_0}^{\mathrm{quot}}
\ge c\min\left\{a_0^2,\;L_a^2W^{2\beta},\;\frac1{W\bar\lambda_A}\right\}.
\tag{D.11}
\end{equation}
When local photon rates have common order $\bar\lambda$, optimizing $W$ yields
\[
R_{n_0}^{\mathrm{quot}}\gtrsim L_a^{2/(2\beta+1)}\bar\lambda^{-2\beta/(2\beta+1)},
\]
the direct Poisson counterpart of the upper rate (D.9); these optimized displays assume $1\lesssim W_*\lesssim N$, and otherwise (D.11) stays in the one-frame or finite-record saturation regime. Thus the $1/(W\bar\lambda)$ law is a \emph{local-minimax} statement in the regime $W\bar\lambda\to\infty$; for bounded total information the risk saturates at the squared diameter of the local parameter neighborhood. It does \emph{not} route through Theorem~D's Gaussian bound (a Gaussian model is not a submodel of the mixed-Poisson experiment).

\emph{Integrated lower bound (guarded Assouad hypercube).} Le Cam's two-point argument controls the pointwise loss only; the integrated quotient loss needs a hypercube. Partition an interior fraction of $\{1,\dots,N\}$ into $m\asymp N/W$ cells of length $CW$; inside each cell reserve a guard band of order $W$ and place a translated, rescaled copy $\psi_j$ of one fixed compactly supported, mean-zero $H_\beta$ template that vanishes on the cell boundary, with active supports separated by at least $cW$ and $\|\psi_j\|_2^2\ge cW$. Set $\ell^\omega=\ell^{(0)}+a\sum_j\omega_j\psi_j$, $\omega\in\{0,1\}^m$, with
\[
\bar\lambda_*=\max_j|A_j|^{-1}\sum_{i\in A_j}\lambda_i^{(0)},
\]
\[
a\le c_0\min\{a_0,\,L_aW^\beta,\,(W\bar\lambda_*)^{-1/2}\}.
\]
The guard bands and the vanishing boundary keep the H\"older constant of the sum bounded independently of $m$, and the interior baseline keeps every vertex in the ball. Adjacent vertices differ on one block with $\operatorname{KL}\le Ca^2W\bar\lambda_*\le\kappa<1$; all pairwise differences are mean-zero, so quotient projection loses no distance, and Assouad's lemma~\cite{Tsybakov2009} gives
\begin{equation}
\begin{aligned}
R_{\mathrm{int}}^{\mathrm{quot}}
&=\inf_{\hat\ell}\sup_\omega\frac1N\mathbb E_\omega\|P_N(\hat\ell-\ell^\omega)\|_2^2\\
&\ge c\min\left\{a_0^2,\;L_a^2W^{2\beta},\;\frac1{W\bar\lambda_*}\right\}.
\end{aligned}
\tag{D.12}
\end{equation}

\emph{Random carriers: choose alternatives before observing $B$.} If $q_i=\Lambda_0B_i$ is random with a parameter-independent law and $\mathbb EB_i<\infty$, the alternatives and their amplitude must be chosen \emph{before} observing $B$, using the \emph{expected} block information $\overline{\mathcal I}_A=\mathbb E_B\sum_{i\in A}\lambda_i^{(0)}\psi_i^2$ --- not the realized one. The calculation then applies verbatim to the carrier-observed \emph{joint oracle} experiment $(C,B)$, whose KL is $\operatorname{KL}(P_0^{C,B},P_1^{C,B})=\mathbb E_B\operatorname{KL}(P_0^{C\mid B},P_1^{C\mid B})$; marginalization gives $\operatorname{KL}(P_0^{C},P_1^{C})\le\operatorname{KL}(P_0^{C,B},P_1^{C,B})$, so the oracle lower bound transfers to the less-informative count-only experiment. This transfers \emph{hardness}; it does not construct a blind estimator. For dark count $d>0$, replace the total count information by
\[
\overline{\mathcal I}_A(\psi)=\mathbb E_B\sum_{i\in A}\frac{s_i^2}{s_i+d}\,\psi_i^2,
\qquad s_i=\Lambda_0e^{\ell_i}b_i,
\]
the local quadratic-KL information in the contrast direction (the bounded-perturbation remainder must be controlled uniformly); the identity $I(\theta)=\lambda$ does not hold for log gain in the presence of additive dark counts.

\paragraph{Oracle local-shift information (the Fig.~5 reference line).} For a fixed window $A$, suppose independently $C_k\sim\operatorname{Pois}\{\lambda_k(t)\}$ with $\lambda_k(t)=s_ke^t+d_k$, where $s_k>0$ and $d_k\ge0$ are \emph{known}. The score and Fisher information for the one-dimensional common shift $t$ are
\[
S_t(C)=\sum_{k\in A}\{C_k-\lambda_k(t)\}\frac{s_ke^t}{\lambda_k(t)},
\]
\begin{equation}
I_A(t)=\sum_{k\in A}\frac{s_k^2e^{2t}}{s_ke^t+d_k},
\qquad
I_A(t)=\sum_{k\in A}\lambda_k(t)\ \ \text{at }d_k=0 .
\tag{D.13}
\end{equation}
Therefore $1/I_A(t)$ is the scalar Cram\'er--Rao reference for a regular unbiased estimator of a \emph{common log-intensity shift} with carriers, offsets, and all other path coordinates held fixed --- an \textbf{oracle common-shift diagnostic}, \emph{not} the nuisance-path quotient CRB. For bias $b(t)=\mathbb E_t\hat t-t$ the biased CR inequality reads
\[
\operatorname{MSE}_t(\hat t)\ge b(t)^2+\frac{\{1+b'(t)\}^2}{I_A(t)} :
\]
MSE \emph{below} $1/I_A$ is compatible with shrinkage bias, and MSE \emph{above} $1/I_A$ does not establish unbiasedness. For a path model $\ell=U\theta$ at $d=0$ with estimable contrast coordinates separated from the constant direction and $D_\lambda=\operatorname{diag}(\lambda_1,\dots,\lambda_N)$, the information is $I_\theta=U^\top D_\lambda U$; eliminating a common nuisance scale $m$ entering as $\ell+m\mathbf1$ gives the Schur complement
\[
I_{\theta\mid m}=U^\top D_\lambda U
-U^\top D_\lambda\mathbf1(\mathbf1^\top D_\lambda\mathbf1)^{-1}\mathbf1^\top D_\lambda U,
\]
the quotient Fisher object of this known-carrier parametric path model. It does not include unknown-object or unknown-carrier nuisance parameters; a contrast outside the range of the information matrix has no finite CR bound, and overlapping windows create cross-window covariance, so an average of $1/I_n$ is not an exact whole-record centered-loss CRB. Sharp minimax constants beyond these local calculations are not claimed.

\paragraph{Remark D.4.1 (instantiation in the Fig.~5 protocol).} There $d=0$, the carriers $b_k$ are \emph{oracle} values supplied from the simulation truth --- the noiseless bucket carrier of the simulated object, times the photon budget; operationally a flat-field-calibrated benchmark, \emph{not} programmed design intensities --- and $\Theta$ is the safe tabulated operating range $[\lambda_-/\min_kb_k,\lambda_+/\max_kb_k]$ with $\Lambda_0$ fixed by convention (the record-mean centering absorbs the global gauge). The implemented estimator solves $M_n(\hat\theta_n)=y_n$ per window by monotone bisection on the exactly tabulated homogeneous $m_\alpha$ --- for $C\sim\mathrm{Pois}(\theta b)$, $\mathbb{E}\psi_\alpha(C)=m_\alpha(\theta b)$ with the \emph{same} homogeneous curve, so no carrier-averaged table is needed --- with the clamp of (D.8a); its windows are truncated and renormalized to their distinct in-record members at the record boundaries (the equal-weight instance of the weighted statement; no replicate padding; requested width $W=64$ realizes centered odd effective width $65$ in the interior). Calibration accuracy: between grid points the table is evaluated by linear interpolation of $m_\alpha$ against $\log\lambda$; the measured midpoint interpolation error on the dense check grid of all 2999 geometric midpoints is $\varepsilon_{\mathrm{cal}}=1.6410165\times10^{-6}$ (attained near $\lambda\approx0.92$) --- an \textbf{empirical midpoint maximum, not a certified uniform bound over the continuous interval}; the certificate route is (D.8c) --- and the 60-step bisection contributes $\varepsilon_{\mathrm{bis}}\le1.5396755\times10^{-17}$ in \emph{log-parameter} units. Per (D.8b) these enter the RMSE bound as $\varepsilon_{\mathrm{cal}}/\underline\kappa_n+\varepsilon_{\mathrm{bis}}$ (the two terms have the same units only \emph{after} the calibration error is divided by the sensitivity; they must not be summed in the response domain and jointly divided). The earlier (r3) arm inverted the homogeneous curve at the window mean, $m_\alpha^{-1}(y_n)$: a \emph{mean-Poisson calibrated proxy} that ignores the within-window carrier variation and coincides with (D.8a) only for a constant carrier. Both arms are reported in Fig.~5; at the protocol's small carrier coefficient of variation they agree numerically, but only (D.8a) is the estimator this theorem is about.

\paragraph{Remark D.4.2 (the $\beta\le1$ theorem and the precise
obstruction above one).}
For $\beta\in(0,1]$, the pointwise modulus
$|\ell_k-\ell_n|\le L_a|k-n|^\beta$ and the mean value theorem give
\[
|q_k(\ell_k)-q_k(\ell_n)|\le\overline\kappa_kL_a|k-n|^\beta,
\]
so the triangle-inequality bias bound used in Theorem C is valid for
arbitrary known carriers.

For $\beta\in(1,2]$, standard smoothness has the local form
\[
\ell_k-\ell_n=v_nu_k+r_{nk},\quad
u_k=\frac{k-n}{N},\quad
|r_{nk}|\le L|u_k|^\beta,
\]
rather than a pointwise modulus of order $\beta$ on the function
itself. Let $w_k\ge0$ be the window weights and write
$q_k(t)=m_\alpha(\Lambda_0e^tb_k+d)$. Taylor's theorem gives
\[
\begin{aligned}
|B_n|\le{}&
|v_n|\left|\sum_kw_kq_k'(\ell_n)u_k\right|
+Q_1L\sum_kw_k|u_k|^\beta\\
&+Q_2\left\{v_n^2\sum_kw_ku_k^2+L^2\sum_kw_k|u_k|^{2\beta}\right\},
\end{aligned}
\tag{D.14}
\]
where $Q_1=\sup_{k,t\in\Theta}q_k'(t)$ and
$Q_2=\sup_{k,t\in\Theta}|q_k''(t)|$.

The first term in (D.14) is the load-bearing obstruction: an
$O(h_n^\beta)$ deterministic bias requires the
\emph{sensitivity-weighted first-moment balance}
\[
|v_n|\left|\sum_kw_kq_k'(\ell_n)u_k\right|\lesssim\underline\kappa_n(W/N)^\beta .
\]
An \emph{internal symmetric equal-weight window with a homogeneous
carrier} satisfies it exactly ($q_k'$ constant and
$\sum_kw_ku_k=0$ --- the interior identity). A heterogeneous carrier
generally does not: ordinary centering cancels $\sum_kw_ku_k$ but not
the carrier-weighted moment, which is then $O(W/N)$, so an
$O((W/N)^\beta)$ claim is false for $\beta>1$ under the present
assumptions. Crucially, a \textbf{truncated one-sided edge window
fails the balance even with a homogeneous carrier}: for the
implementation's first truncated window $k=0,\dots,32$ with target
$n=0$,
\[
\frac1{33}\sum_{k=0}^{32}\frac kN=\frac{16}{N},
\]
first order in the window radius, so the current local-constant
estimator cannot support an all-frame uniform $\beta>1$ rate ---
the $\beta>1$ statement is restricted to \emph{internal balanced
windows}.

Under the balance condition the remaining nonlinear term
is $O((W/N)^2)$ and is therefore no larger than $O((W/N)^\beta)$ for
$1<\beta\le2$, provided the local slope is bounded. Thus the
calibrated inverse recovers the $\beta$-order deterministic bias
on internal windows under homogeneous or sensitivity-balanced
carriers, with a constant depending also on the slope bound and on
$Q_2$.

A second-order Taylor expansion of $M_n^{-1}$ by itself does not
remove an unbalanced first moment. For arbitrary carriers, or for
edge frames, recovery of the $\beta>1$ rate requires carrier-adapted
moment conditions, carrier-adapted local-linear weights, or a
separately proved local-polynomial edge estimator (for $\beta>2$,
moments through order $\lfloor\beta\rfloor$ must be cancelled).
Theorem C therefore remains restricted to
$\beta\le1$, while the internal-window homogeneous-carrier and
sensitivity-balanced extensions above are valid weaker statements.

\paragraph{Remark D.4.3 (unknown random carrier; nonstationary
finite-$W$ covariance bound).}
Let $\{B_k\}$ be an unobserved stationary random carrier with known
law, independent of the parameter, and let
\[
\bar m_{\alpha,k}(t)
=\mathbb E_B\,m_\alpha(\Lambda_0e^tB_k+d).
\]
Because the gain path $\ell_k$ may vary, stationarity of $\{B_k\}$
does not in general make the soft-log sequence stationary. Define
\[
z_k=\psi_\alpha(C_k)-\mathbb E\psi_\alpha(C_k),
\qquad
g_k(b)=m_\alpha(\Lambda_0e^{\ell_k}b+d).
\]
Conditional independence of the Poisson randomization gives, for
$i\ne j$,
\[
\operatorname{Cov}(z_i,z_j)
=\operatorname{Cov}(g_i(B_i),g_j(B_j)).
\]
Consequently the exact finite-window identity is
\[
\begin{aligned}
\operatorname{Var}\!\left(W^{-1}\sum_{i=1}^Wz_i\right)
={}&W^{-2}\sum_{i=1}^W\operatorname{Var}(z_i)\\
&+2W^{-2}\sum_{h=1}^{W-1}
\sum_{i=1}^{W-h}\operatorname{Cov}(z_i,z_{i+h}).
\end{aligned}
\tag{D.15}
\]
Let $p=2+\eta$, $r_\eta=\eta/(2+\eta)$, and assume
\[
v_*:=\sup_i\operatorname{Var}(z_i)<\infty,
\]
\[
G_p:=\sup_i\|g_i(B_i)-\mathbb Eg_i(B_i)\|_p<\infty.
\]
If the carrier is strongly mixing with coefficient $\alpha_B(h)$,
Davydov's inequality~\cite{Davydov1968} gives
\[
|\operatorname{Cov}(z_i,z_{i+h})|
\le8\,\alpha_B(h)^{r_\eta}G_p^2.
\]
Substitution into (D.15) yields the explicit finite-$W$ bound
\[
\begin{aligned}
&\operatorname{Var}\!\left(W^{-1}\sum_{i=1}^Wz_i\right)\\
&\quad\le
\frac1W\left[
v_*+16G_p^2\sum_{h=1}^{W-1}
\left(1-\frac hW\right)\alpha_B(h)^{r_\eta}
\right]\\
&\quad\le
\frac{v_*+16G_p^2\sum_{h\ge1}\alpha_B(h)^{r_\eta}}{W}.
\end{aligned}
\]
The factor $16$ is the constant obtained from Davydov's constant $8$
and the two covariance triangles. It is sharper to apply Davydov to
the conditional means $g_i(B_i)$ than to $z_i$ itself, because
independent Poisson shot noise contributes only to the diagonal
variance $v_*$.

If $\ell_k$ is constant and the resulting marked process is strictly
stationary, write $\gamma_\psi(h)=\operatorname{Cov}(z_0,z_h)$. Then
\[
\begin{aligned}
\operatorname{Var}\!\left(W^{-1}\sum_{k=1}^Wz_k\right)
&=
W^{-1}\sum_{|h|<W}\left(1-\frac{|h|}{W}\right)\gamma_\psi(h)\\
&\le\frac{\sum_h|\gamma_\psi(h)|}{W},
\end{aligned}
\]
and when $\sum_h|\gamma_\psi(h)|<\infty$,
\[
\begin{aligned}
\operatorname{Var}\!\left(W^{-1}\sum_{k=1}^Wz_k\right)
&=\frac{\sigma_{\psi,\mathrm{LR}}^2}{W}+o(W^{-1}),\\
\sigma_{\psi,\mathrm{LR}}^2
&=\sum_h\gamma_\psi(h).
\end{aligned}
\]
The multiplicative notation
$\sigma_{\psi,\mathrm{LR}}^2W^{-1}\{1+o(1)\}$ is valid only when
$\sigma_{\psi,\mathrm{LR}}^2>0$. For
$z_n=\epsilon_n-\epsilon_{n-1}$ the signed long-run variance is zero
while the finite-window variance is $2\sigma_\epsilon^2/W^2$.

\paragraph{Remark D.4.4 (clamp probability: a diagnostic, not an
extra risk term).}
By the global Lipschitz property of the clamped inverse, (D.8)
already covers clamped, all-zero, and out-of-range windows; no term
$D_\Theta^2\,\mathbb P(\mathcal E_n^c)$ is added to the risk. The
clamp probability remains useful as a \emph{diagnostic} for whether
the localized variance regime of the low-photon expansion applies
(low total information drives clamp saturation, in which case the
inverse-information reading of (D.10) is off scope). With
$\mathcal E_n=\{y_n\in M_n(\Theta)\}$, use the notation
$B_n,V_n,\underline\kappa_n$ of the finite-window proof and put
\[
r_n=
\min\{M_n(\ell_n)-M_n(\theta_{\min}),
M_n(\theta_{\max})-M_n(\ell_n)\}.
\]
If $\ell_n$ has margin $\Delta_n$ from $\partial\Theta$, then
$r_n\ge\underline\kappa_n\Delta_n$. Since
$y_n-M_n(\ell_n)=B_n+Z_n$, with
$\mathbb EZ_n=0$ and $\operatorname{Var}(Z_n)=V_n$,
\[
\mathcal E_n^c\subseteq\{|B_n+Z_n|\ge r_n\}.
\]
Therefore
\[
\mathbb P(\mathcal E_n^c)
\le\min\left\{1,\frac{B_n^2+V_n}{r_n^2}\right\}
\le\min\left\{1,\frac{B_n^2+V_n}{(\underline\kappa_n\Delta_n)^2}\right\}.
\]
If $|B_n|<r_n$, the sharper centered bound is
\[
\mathbb P(\mathcal E_n^c)
\le\min\left\{1,\frac{V_n}{(r_n-|B_n|)^2}\right\}.
\]
An exponential Bernstein bound requires a positive residual margin
$r_n-|B_n|$; it is not available merely from a positive parameter
margin when the deterministic smoothing bias can reach the boundary.

On $\mathcal E_n^c$, both the clamped estimate and the target lie in
$\Theta$, so the clamp-event slice of the MSE is trivially at most
$D_\Theta^2\,\mathbb P(\mathcal E_n^c)$ --- but this bound is
\emph{never added} to (D.8), which already dominates it through the
global Lipschitz argument; it is recorded only to size the diagnostic.

For a window containing $W$ distinct, conditionally independent
Poisson counts,
\[
\begin{aligned}
&\mathbb P(C_k=0\ \hbox{for all }k\in W_n)\\
&\quad=\exp\!\left(-\sum_{k\in W_n}\lambda_k(\ell_k)\right)
\le e^{-W\lambda_-}.
\end{aligned}
\]
If a padding convention repeats observations, the sum and the
cardinality in this formula are over distinct frames after repeated
weights have been aggregated; in that case the exponent need not be
$W\lambda_-$ (the first width-65 replicate-padded window of the
former convention had only 33 distinct frames, so its all-zero
probability was $e^{-33\lambda}$, not $e^{-65\lambda}$ --- at
$\lambda=0.25$, $2.6\times10^{-4}$ versus $8.8\times10^{-8}$).

\paragraph{Remark D.4.5 (scope: oracle carriers, known carrier law, and blind recovery --- three distinct estimators).}
Theorem C is conditional on the realized per-frame carriers $(b_k)$ being \emph{known}: this is precisely the oracle-known-carrier arm implemented in the low-photon experiment (Fig.~5), and it does \textbf{not} establish blind low-photon recovery from an unknown carrier law. Theorem~B, by contrast, is the \emph{blind high-count} stationary-carrier result; the two must not be merged. If an unobserved carrier $B$ has a fully \emph{known law} $F_B$, one may instead calibrate the different marginal curve
\[
\bar q_{F_B}(t)=\mathbb E_{B\sim F_B}\,m_\alpha(\Lambda_0e^tB+d),
\]
a \emph{third}, law-calibrated estimator that needs its own variance or long-run-variance theorem (Remark D.4.3 supplies the covariance tool); it is not the per-frame known-carrier estimator above. If $F_B$ is \emph{unknown}, $\bar q_{F_B}$ is an unknown nonlinear function which, unlike the high-count logarithmic model, cannot generally be absorbed into one additive gauge constant: with $s=\Lambda_0e^t$,
\[
m_\alpha(sB)=\log\alpha+c_\alpha sB+a_{2,\alpha}s^2B^2+O(s^3B^3),
\]
\[
a_{2,\alpha}=\tfrac12\log\frac{\alpha(\alpha+2)}{(\alpha+1)^2}\ne0,
\]
so the laws $B\equiv1$ and $\Pr(B=\tfrac12)=\Pr(B=\tfrac32)=\tfrac12$ --- same mean, different second moments --- produce \emph{different} soft-log response curves. The known-carrier theorem plus one homogeneous calibration inverse therefore does not yield an unknown-law blind theorem (though this does not rule out every joint semiparametric method).

\paragraph{Centered-log quotient risk.}
Let $e=\hat\ell-\ell$ and $P_N=I-N^{-1}\mathbf1\mathbf1^\top$. Since $P_N$ is an orthogonal projection,
\[
\frac1N\,\mathbb E\|P_Ne\|_2^2\le\frac1N\sum_{n=1}^N\mathbb Ee_n^2 ,
\]
so the pointwise bounds transfer directly to the record-centered log-gain loss of Fig.~5; overlapping windows need not be independent for this expectation bound.

\begin{statusblock}\textbf{Shrinkage-bias caveat (load-bearing, per Sec.~6).} For $\bar\lambda<1$ the windowed soft-log statistic is \emph{bias-fragile}: the mass at $C=0$ pulls the weighted mean of $\psi_\alpha$ toward $\log\alpha$, and the amplification factor $(\kappa_{\max}/\underline\kappa_n)^2$ in the specialized bias term of (D.8) blows up as the window sensitivity vanishes with $\bar\lambda_n\to0$ (at $d=0$; window-average form $M_n'(\ell_n)\approx c_\alpha\bar\lambda_n$). Empirically the MSE of the \emph{uncalibrated} proxy (and of Anscombe) can then sit \emph{below} the Fisher reference curve --- that is bias, not super-efficiency, and it does not contradict the lower bound (which constrains variance-limited, locally unbiased behavior). The $1/(W\bar\lambda)$ law is asserted only in the variance-limited regime above the $\bar\lambda\sim1$ crossover (and under the local gain flatness of (D.10)); below it, honest accounting reports the bias floor, and all Fisher comparisons are made in the log-domain loss of this theorem (Fig.~5's \texttt{log\_gain\_mse} metric, against the \emph{oracle common-shift diagnostic} $\mathrm{mean}_n[1/I_n]$ with $I_n=\sum_{k\in W_n}\lambda_k$ at $d=0$, per (D.13) --- \emph{not} a nuisance-path quotient CRB) --- the gain-domain relMSE column retained for continuity is a different loss and is not Fisher-comparable. Recovering the rate below the crossover requires the offset-design, Anscombe~\cite{Anscombe1948} ($2\sqrt{C+3/8}$), or full Poisson-mixture MLE variants; for the mixture MLE the claim $\mathrm{MSE}\approx1/(WJ(\theta))$ with $J\sim\mathbb{E}\lambda$ at low photons is standard parametric asymptotics under regularity of the mixture likelihood and is \emph{not} re-proven here.\end{statusblock}

\paragraph{What this does not claim.} No absolute gain scale is estimated (gauge remains). The calibration curve $m_\alpha$ and the carriers $b_k$ are assumed known; in the Fig.~5 protocol the carriers are \emph{oracle} values from the simulation truth (a flat-field-calibrated benchmark), so Fig.~5 upper-bounds what a real flat-field calibration could achieve rather than demonstrating blind carrier recovery --- Theorem C is an oracle-known-carrier result, Theorem~B is the blind high-count result, and the known-law estimator of Remark D.4.5 is a third object; none substitutes for the others. Miscalibration beyond the measured interpolation term ($\varepsilon_{\mathrm{cal}}/\underline\kappa_n+\varepsilon_{\mathrm{bis}}$ of (D.8b), with $\varepsilon_{\mathrm{cal}}$ an empirical midpoint maximum pending the (D.8c) certificate) adds a bias not covered by (D.8). Constants in (D.8)--(D.9) are not tracked. Smoothness beyond $\beta=1$ is claimed only in the internal-window homogeneous-carrier / sensitivity-balanced form of Remark D.4.2 (display (D.14)); for arbitrary carriers the carrier-weighted first moment obstructs every $\beta>1$ rate, and the implementation's truncated edge windows are not covered even under a homogeneous carrier. The unconditional $1/(W\bar\lambda)$ law is withdrawn in favor of (D.10); the flat-gain specialization is the only $1/(W\bar\lambda)$ statement retained. The minimax matching of (D.9) is at rate level only, \emph{in the log-domain quotient loss and at $d=0$}, proved by the direct Poisson Le~Cam/Assouad bounds (D.11)--(D.12) --- not inherited from Theorem~D, whose Gaussian model is not a submodel of the mixed-Poisson experiment; for $d>0$ every information statement must replace $W\bar\lambda$ by $\sum_k\mathbb{E}_B\,s_k^2/(s_k+d)$. The plotted reference $\mathrm{mean}_n[1/I_n]$ is the oracle common-shift diagnostic of (D.13): MSE above it does not establish unbiasedness, and MSE below it signals shrinkage bias, not super-efficiency. No optimality of any kind is claimed in the gain-domain (linear) metric, and no claim is made that the calibrated estimator beats bias-accepting transforms (e.g.\ Anscombe) in MSE below the $\bar\lambda\sim1$ crossover --- empirically it does not (Sec.~9). As $\Phi\to0$ with no offset and no reference, Fisher information for the gain vanishes --- an information limit no estimator evades (Sec.~10).

\section{SRHT whitening: exact Walsh condition and permutation bounds}

This appendix proves the claims of Sec.~7: the exact covariance identity $\operatorname{Cov}_D(Z_g,Z_h)=\hat w(g+h)$, the necessary-and-sufficient Walsh-flatness conditions for sign-only whitening (pairwise, spectral-window, and window-average forms), the Bernstein--Serfling permutation bound and its union-bound corollary, the permutation-alone carrier-variance identity with its flat-object counterexample, and the Hanson--Wright window-energy bound.

\subsection{E.0 Setting and standing assumptions}

\begin{itemize}
\item \textbf{(E1)} $K=2^m$ and pixels are indexed by the Walsh group $G=\mathbb{F}_2^m$, with characters $\chi_g(j)=(-1)^{g\cdot j}$, unnormalized Hadamard matrix $H_{g,j}=\chi_g(j)$, and orthonormal carrier $U=K^{-1/2}H$. For every $g\neq 0$, $\chi_g$ is balanced: exactly $K/2$ entries equal $+1$.
\item \textbf{(E2)} The object $T\in\mathbb{R}^K$ is deterministic and nonzero. Write $w_j=T_j^2$; $S_1$, $S_2$, and the spread functionals $K_{\mathrm{eff}}$, $K_\infty$, $K_4$ are as in Appendix~C. The Walsh transform of $w$ is $\hat w(h)=\sum_j \chi_h(j)\,w_j$, so $\hat w(0)=S_2$.
\item \textbf{(E3)} The SRHT coefficient vector is $x=UDPT$, where $D=\mathrm{diag}(\eta_1,\dots,\eta_K)$ with i.i.d.\ Rademacher signs $\eta_j$, and $P$ is a uniformly random pixel permutation independent of $D$. The normalized carrier is $Z_g=\sqrt{K}\,x_g=\sum_j \chi_g(j)\,\eta_j\,(PT)_j$. Physical nonnegative patterns are the offset form $B_g=\mu S_1+\sigma Z_g$ with positivity margin $\mu S_1>\sigma|Z_g|$; the bucket record is $R_n=a_nB_n$ with $\ell=\log a$ as in the main text. A \emph{window} is a subset $A\subset G$ of $W=|A|$ rows.
\end{itemize}

Permutation leaves $S_1$, $S_2$, $K_{\mathrm{eff}}$, $K_4$, $K_\infty$ invariant; it only reorders $w$, i.e.\ replaces $\hat w$ by $\widehat{w^P}$ with $w^P_j=T^2_{Pj}$.

\subsection{E.1 The exact covariance identity}

\begin{namedthm}{Lemma E.1}
Under (E1)--(E3) with $P$ fixed (in particular $P=\mathrm{id}$), $\mathbb{E}_D Z_g=0$ and for all $g,h\in G$
\begin{equation}
\operatorname{Cov}_D(Z_g,Z_h)\;=\;\sum_j \chi_{g+h}(j)\,(PT)_j^2\;=\;\widehat{w^P}(g+h). \tag{E.1}
\end{equation}
\end{namedthm}

\noindent\textbf{Proof.} $\mathbb{E}\eta_j=0$ gives the mean. Since $\mathbb{E}\,\eta_i\eta_j=\delta_{ij}$,
\[
\begin{aligned}
\mathbb{E}_D Z_gZ_h&=\sum_{i,j}\chi_g(i)\chi_h(j)(PT)_i(PT)_j\,\mathbb{E}\eta_i\eta_j\\
&=\sum_j \chi_g(j)\chi_h(j)(PT)_j^2,
\end{aligned}
\]
and $\chi_g\chi_h=\chi_{g+h}$ by the character property of the Walsh group. \hfill$\blacksquare$

Consequently the covariance of any window $A$ is the pattern matrix $\Gamma_A(g,h)=\widehat{w^P}(g+h)$, and, normalized by the common marginal variance $S_2$, $\Gamma_A/S_2=I_A+R_A$ with off-diagonal entries $R_A(g,h)=\widehat{w^P}(g+h)/S_2$, $g\neq h$. The entire gain--object coupling of sign-randomized acquisition is therefore carried by the non-DC Walsh spectrum of the squared (permuted) object --- nothing else.

\subsection{E.2 Sign-only whitening: necessary and sufficient conditions}

\begin{namedthm}{Theorem E.2 (exact Walsh condition)}
Under (E1)--(E3) with no permutation ($P=\mathrm{id}$):
\begin{enumerate}
\item \emph{(exact pairwise)} $\operatorname{Cov}_D(Z_g,Z_h)=S_2\,\delta_{g,h}$ for all $g,h$ \textbf{iff} $\hat w(q)=0$ for every $q\neq0$, equivalently \textbf{iff} $T_j^2=S_2/K$ for all $j$ ($T^2$ flat).
\item \emph{(pairwise $\varepsilon$-whitening)} $|\operatorname{Cov}_D(Z_g,Z_h)|\le\varepsilon S_2$ for all $g\neq h$ \textbf{iff} $\max_{q\neq0}|\hat w(q)|\le\varepsilon S_2$.
\item \emph{(spectral window)} For a window class $\mathcal{A}$, $\|\Gamma_A/S_2-I_A\|_{\mathrm{op}}\le\varepsilon$ for all $A\in\mathcal{A}$ \textbf{iff} $\sup_{A\in\mathcal{A}}\|R_A\|_{\mathrm{op}}\le\varepsilon$.
\item \emph{(window average)} With $\bar Z_A=W^{-1}\sum_{g\in A}Z_g$ and multiplicity $m_A(q)=\#\{(g,h)\in A\times A: g\neq h,\ g+h=q\}$, the variance $\operatorname{Var}(\bar Z_A)$ lies in $[(1\mp\varepsilon)S_2/W]$ for every $A\in\mathcal{A}$ \textbf{iff} $\sup_{A\in\mathcal{A}}\big|\sum_{q\neq0}m_A(q)\,\hat w(q)\big|\le\varepsilon\,W S_2$.
\end{enumerate}
\end{namedthm}

\noindent\textbf{Proof.} (1) If whitening holds, then for any $q\neq0$ pick $g,h$ with $g+h=q$ (always possible in $G$); (E.1) forces $\hat w(q)=0$. Conversely, vanishing non-DC coefficients make (E.1) equal $S_2\delta_{g,h}$. The flatness equivalence is Walsh inversion: the transform is invertible, so $\hat w$ supported on $\{0\}$ iff $w$ is constant, i.e.\ $T_j^2=S_2/K$. (2) Immediate from (E.1) after dividing by $S_2$. (3) The normalized window covariance \emph{is} $I_A+R_A$; the statement is definitional. (4) Direct expansion:
\[
\begin{split}
\operatorname{Var}(\bar Z_A)&=W^{-2}\!\!\sum_{g,h\in A}\hat w(g+h)\\
&=\frac{S_2}{W}+W^{-2}\sum_{q\neq0}m_A(q)\,\hat w(q),
\end{split}
\]
and the two-sided requirement is exactly the stated inequality. \hfill$\blacksquare$

Two standard reductions relate the forms. If $\mathcal{A}$ contains all two-point windows, (3) implies (2), and (2) is the sharpest scalar condition. In the converse direction, since $R_A$ has zero diagonal, Gershgorin's disc theorem gives $\|R_A\|_{\mathrm{op}}\le(W-1)\max_{q\neq0}|\hat w(q)|/S_2$, so pairwise condition (2) at tolerance $\varepsilon/(W_{\max}-1)$ implies spectral condition (3) for all windows with $W\le W_{\max}$. This Gershgorin step is a (generally loose) sufficient bound, not an equivalence.

The obstruction is sharp: if $T^2$ is aligned with the positive support of a single character $\chi_q$, then $|\hat w(q)|=S_2$ and the row pair $(g,h)$ with $g+h=q$ is \emph{perfectly} correlated under sign randomization. Signs cannot remove this; it is a property of the object's pixel ordering.

\paragraph{What this does not claim.} Theorem E.2 characterizes second-moment decorrelation over the sign draw for a \emph{fixed} object and ordering; it does not assert Gaussianity of $Z_g$ (a Rademacher sum, whose normal approximation requires $K_\infty\to\infty$ with Berry--Esseen error $\lesssim K_\infty^{-1/2}$), independence of rows, or anything about noise, gain drift, or estimation error --- those enter through Secs.~5--6 and Appendix~F. It is also a statement about the design/object pair, not a high-probability event: no randomness beyond $D$ is involved.

\subsection{E.3 Random permutation as a probabilistic Walsh-flattener}

By Lemma E.1, Theorem E.2 holds verbatim for a realized permutation $P$ with $w$ replaced by $w^P$: the \emph{deterministic} necessary-and-sufficient condition remains Walsh-flatness, now of the permuted squared object. A random $P$ makes that condition likely.

\begin{namedthm}{Theorem E.3 (Bernstein--Serfling permutation bound)}
Under (E1)--(E3), there exist universal constants $c,C>0$ such that for every fixed $q\neq0$ and $\varepsilon>0$,
\begin{equation}
\Pr_P\!\big(|\widehat{w^P}(q)|\ge\varepsilon S_2\big)\;\le\;2\exp\!\big[-c\min(\varepsilon^2K_4,\ \varepsilon K_\infty)\big], \tag{E.2}
\end{equation}
and by the union bound over the $K-1$ nonzero frequencies,
\begin{equation}
\begin{aligned}
&\Pr_P\!\big(\max_{q\neq0}|\widehat{w^P}(q)|\ge\varepsilon S_2\big)\\
&\qquad\le\;2(K-1)\exp\!\big[-c\min(\varepsilon^2K_4,\ \varepsilon K_\infty)\big].
\end{aligned}
\tag{E.3}
\end{equation}
In particular, pairwise $\varepsilon$-whitening (Theorem E.2(2) for $w^P$) holds with probability at least $1-\delta$ whenever
\begin{equation}
\min(\varepsilon^2K_4,\ \varepsilon K_\infty)\;\ge\;C\log(K/\delta). \tag{E.4}
\end{equation}
\end{namedthm}

\noindent\textbf{Proof.} Fix $q\neq0$ and let $A_q=\{j:\chi_q(j)=+1\}$, so $|A_q|=K/2$ by (E1). Then
\[
\widehat{w^P}(q)=\sum_{j\in A_q}w^P_j-\sum_{j\notin A_q}w^P_j=2\sum_{j\in A_q}w^P_j-S_2,
\]
a centered sum of $n=K/2$ draws \emph{without replacement} from the finite population $\{w_1,\dots,w_K\}$ with mean $\mu_w=S_2/K$. Its exact variance is $n\sigma_w^2\frac{K-n}{K-1}\le\frac{K}{4}\sigma_w^2\le\frac14\sum_j w_j^2=\frac{S_2^2}{4K_4}$, where $\sigma_w^2=K^{-1}\sum_j(w_j-\mu_w)^2$; and each summand satisfies $0\le w_j\le\|T\|_\infty^2=S_2/K_\infty$. The Bernstein inequality for sampling without replacement (Serfling~\cite{Serfling1974}; the Bernstein--Serfling form of Bardenet--Maillard~\cite{BardenetMaillard2015}) then gives, for $t>0$,
\[
\Pr\big(|\widehat{w^P}(q)|\ge t\big)\le2\exp\!\Big[-c\min\Big(\frac{t^2}{S_2^2/K_4},\ \frac{t}{S_2/K_\infty}\Big)\Big].
\]
Setting $t=\varepsilon S_2$ yields (E.2); (E.3) is the union bound; (E.4) follows by requiring the right side of (E.3) to be at most $\delta$ and absorbing the factor $2(K-1)$ into $C\log(K/\delta)$. \hfill$\blacksquare$

\paragraph{What this does not claim.} (E.2)--(E.4) are \emph{sufficient} high-probability conditions with unspecified universal constants; they are not the deterministic criterion for a realized acquisition --- that remains Walsh-flatness of $w^P$ (Theorem E.2), which can and should be checked numerically for the ordering actually used. For spectral whitening of all windows up to size $W_{\max}$, the simple route substitutes $\varepsilon\to\varepsilon/(W_{\max}-1)$ in (E.4) via Gershgorin; sharper window-size dependence via matrix-Bernstein or Hanson--Wright arguments (Sec.~E.5) is a concentration refinement, again sufficient rather than necessary. Both branches of the min collapse when energy concentrates on few pixels ($K_4,K_\infty$ small): permutation cannot flatten a spike.

\subsection{E.4 Permutation alone: exact carrier variance and the flat-object counterexample}

\begin{namedthm}{Theorem E.4 (permutation-alone variance)}
Under (E1)--(E2), with no signs and a uniformly random permutation $P$, the carrier $Z_g(P)=\sum_j\chi_g(j)T_{Pj}$ of any non-DC row $g\neq0$ satisfies $\mathbb{E}_PZ_g=0$ and
\begin{equation}
\operatorname{Var}_PZ_g\;=\;\frac{KS_2-S_1^2}{K-1}\;=\;S_2\,\frac{K-K_{\mathrm{eff}}}{K-1}. \tag{E.5}
\end{equation}
Consequently, for offset physical patterns $B_g=\mu S_1+\sigma Z_g(P)$ with positivity margin,
\begin{equation}
\frac{\operatorname{Var}_PB_g}{(\mathbb{E}_PB_g)^2}=\Big(\frac{\sigma}{\mu}\Big)^2\frac{K-K_{\mathrm{eff}}}{(K-1)\,K_{\mathrm{eff}}}\;\le\;\Big(\frac{\sigma}{\mu}\Big)^2\frac{1}{K_{\mathrm{eff}}}. \tag{E.6}
\end{equation}
No matching lower bound of order $1/K_{\mathrm{eff}}$ holds uniformly over objects: for flat $T$ ($T_j$ constant), $K_{\mathrm{eff}}=K$ while $Z_g(P)=0$ for \textbf{every} permutation and every $g\neq0$.
\end{namedthm}

\noindent\textbf{Proof.} For $g\neq0$, $\sum_i\chi_g(i)=0$, so $\mathbb{E}_PZ_g=\big(\sum_i\chi_g(i)\big)\mathbb{E}_PT_{P1}=0$. Exchangeability of $(T_{P1},\dots,T_{PK})$ gives $\mathbb{E}_PT_{Pi}^2=S_2/K$ and, for $i\neq j$, $\mathbb{E}_PT_{Pi}T_{Pj}=(S_1^2-S_2)/(K(K-1))$. Hence
\[
\mathbb{E}_PZ_g^2=S_2+\frac{S_1^2-S_2}{K(K-1)}\sum_{i\neq j}\chi_g(i)\chi_g(j),
\]
and $\sum_{i\neq j}\chi_g(i)\chi_g(j)=\big(\sum_i\chi_g(i)\big)^2-K=-K$, giving (E.5); the second form follows from $K_{\mathrm{eff}}=S_1^2/S_2$. For (E.6), $\mathbb{E}_PB_g=\mu S_1$ and $\operatorname{Var}_PB_g=\sigma^2\operatorname{Var}_PZ_g$; divide. The counterexample: if $T_j\equiv t$, every non-DC Walsh sum of a constant vector vanishes identically, while $S_1^2/S_2=K$. Note (E.5) also exhibits the finite-population correction --- the relative variance is $(\sigma/\mu)^2(1/K_{\mathrm{eff}}-1/K)\cdot K/(K-1)$. \hfill$\blacksquare$

\paragraph{What this does not claim.} Theorem E.4 does \emph{not} say permutation-alone whitens: it provides finite-population stationarity (exchangeability of the carrier along the acquisition, the ($\bigstar$) mechanism of Sec.~5) and an $O(1/K_{\mathrm{eff}})$ \emph{upper} variance scale, but it is not equivalent to i.i.d.\ Rademacher or nonnegative random-pattern carrier noise, and it guarantees no $\Theta(1/K_{\mathrm{eff}})$ \emph{excitation} --- a flat object has exactly zero. Two-sided coefficientwise control requires the sign flip \emph{and} the spectral-spread condition (E.4) together (the ``both'' cell of the Fig.~8 ablation). An earlier claim (R1/v3) that permutation-alone ``cannot attain variance $\sim1/K_{\mathrm{eff}}$'' was too strong as an upper-scale statement and is corrected by (E.6); it is correct only in the two-sided reading. Moreover, \textbf{full (joint) exchangeability belongs to acquisition-order randomization, not to a shared pixel permutation}: a uniform reordering of the acquisition sequence makes the frame sequence finite-population exchangeable over its randomization draw (with pair blocks, not frames, as the exchangeable units when complementary masks are kept adjacent), whereas one common pixel permutation $P$ shared by all rows gives every non-DC row the \emph{same one-dimensional} law --- mean zero and variance (E.5) --- but \emph{not} full joint row exchangeability, by the triple-moment counterexample (E.9) below. Once a single permutation draw is conditioned upon, the sequence is deterministic; finite-population/randomization concentration should be invoked, not a claim that every realized order is stationary or mixing.

\begin{namedthm}{Proposition E.4$'$ (Walsh triple-moment counterexample to joint row exchangeability)}
Let $K=2^m\ge8$, let $\chi_g$ be the Walsh characters, and $Z_g(P)=\sum_{j}\chi_g(j)T_{Pj}$ under one uniform pixel permutation $P$ shared by all rows. Let $\bar T=K^{-1}\sum_jT_j$ and $S_{3,c}=\sum_j(T_j-\bar T)^3$. For distinct nonzero rows $g,h,r$,
\begin{equation}
\mathbb E_P[Z_gZ_hZ_r]=
\begin{cases}
\frac{K^2}{(K-1)(K-2)}\,S_{3,c}, & g+h+r=0,\\[4pt]
0, & \text{$g,h,r$ lin.\ indep.}
\end{cases}
\tag{E.9}
\end{equation}
For a generic object with $S_{3,c}\neq0$, relational triples ($g+h+r=0$) and independent triples therefore have different joint laws, disproving full row exchangeability under a common pixel permutation.
\end{namedthm}

\noindent\textbf{Proof.} Replace $T_j$ by $x_j=T_j-\bar T$ (non-DC Walsh coefficients unchanged). For the without-replacement sample $X_i=x_{P(i)}$, standard finite-population moments give $\mathbb EX_i^3=S_{3,c}/K$, $\mathbb EX_i^2X_j=-S_{3,c}/\{K(K-1)\}$ for $i\ne j$, and $\mathbb EX_iX_jX_k=2S_{3,c}/\{K(K-1)(K-2)\}$ for distinct $i,j,k$. If $g+h+r=0$, character orthogonality gives sign sums $K$ over the diagonal class, $-K$ over each exactly-two-equal pairing, and $2K$ over the all-distinct class; substitution gives the first line of (E.9). If $g+h+r\neq0$, all the corresponding character sums vanish. \hfill$\blacksquare$

\subsection{E.5 Hanson--Wright window-energy bound}

Whitening of window \emph{energy} --- the quantity the window estimator of Sec.~6 actually averages --- admits a direct quadratic-form bound. For a window $A$ of $W$ rows let $P_A$ select those rows and set
\[
\begin{gathered}
Q_A=\|P_A\,U D P T\|_2^2=\eta^\top M\eta,\\
M=\mathrm{diag}(PT)\,U_A^\top U_A\,\mathrm{diag}(PT),
\end{gathered}
\]
with $\eta$ the Rademacher sign vector and $U_A$ the selected rows of $U$.

\begin{namedthm}{Proposition E.5 (window-energy concentration)}
Condition on any realized $P$. Then $\mathbb{E}_DQ_A=(W/K)S_2$, and there is a universal $c>0$ with
\begin{equation}
\begin{aligned}
&\Pr_D\!\Big(\big|Q_A-\tfrac{W}{K}S_2\big|\ge\varepsilon\,\tfrac{W}{K}S_2\Big)\\
&\qquad\le\;2\exp\!\Big[-c\min\Big(\varepsilon^2\,\tfrac{WK_\infty}{K},\ \varepsilon\,\tfrac{WK_\infty}{K}\Big)\Big].
\end{aligned}
\tag{E.7}
\end{equation}
Over a specified family of $M$ windows, uniform relative energy whitening at tolerance $\varepsilon\le1$ holds with probability $\ge1-\delta$ once
\begin{equation}
\frac{WK_\infty}{K}\;\ge\;C\varepsilon^{-2}\log(M/\delta). \tag{E.8}
\end{equation}
\end{namedthm}

\noindent\textbf{Proof.} Since $U$ has orthonormal rows, $U_AU_A^\top=I_W$, so $U_A^\top U_A$ is an orthogonal projector and $(U_A^\top U_A)_{jj}=\sum_{g\in A}U_{gj}^2=W/K$ (Hadamard entries are $\pm K^{-1/2}$). Hence $M\succeq0$, and $\mathbb{E}_DQ_A=\operatorname{tr}M=\sum_j(PT)_j^2\,(W/K)=(W/K)S_2$. Because $\eta_j^2=1$, the diagonal of $M$ contributes deterministically and $Q_A-\mathbb{E}Q_A$ is a pure off-diagonal Rademacher chaos, to which the Hanson--Wright inequality applies (Hanson--Wright~\cite{HansonWright1971}; sub-Gaussian form of Rudelson--Vershynin~\cite{RudelsonVershynin2013}):
\[
\Pr(|Q_A-\operatorname{tr}M|\ge t)\le2\exp\Big[-c\min\Big(\frac{t^2}{\|M\|_F^2},\ \frac{t}{\|M\|_{\mathrm{op}}}\Big)\Big].
\]
The proxies: $\|M\|_{\mathrm{op}}\le\|\mathrm{diag}(PT)\|_{\mathrm{op}}^2\,\|U_A^\top U_A\|_{\mathrm{op}}\le\|T\|_\infty^2=S_2/K_\infty$, and for PSD $M$, $\|M\|_F^2\le\|M\|_{\mathrm{op}}\operatorname{tr}M\le(W/K)S_2\|T\|_\infty^2$. Substituting $t=\varepsilon(W/K)S_2$ gives $t^2/\|M\|_F^2\ge\varepsilon^2WK_\infty/K$ and $t/\|M\|_{\mathrm{op}}\ge\varepsilon WK_\infty/K$, i.e.\ (E.7); for $\varepsilon\le1$ the quadratic branch is active. A union bound over $M$ windows gives (E.8). \hfill$\blacksquare$

For dense objects of bounded dynamic range, $K_\infty\asymp K_4\asymp K_{\mathrm{eff}}\asymp K$ and (E.8) reduces to the mild $W\gtrsim\varepsilon^{-2}\log M$ --- the window only needs logarithmic length. For an object supported on $m$ comparable pixels, $K_\infty\asymp m$ and the condition hardens to $Wm/K\gtrsim\varepsilon^{-2}\log M$: this is precisely where $K_{\mathrm{eff}}$ alone is the wrong functional and the three spread parameters separate (Sec.~6).

\paragraph{What this does not claim.} Proposition E.5 controls window \emph{energy} over the sign draw for a pre-specified finite window family; it is a sufficient concentration tool, not the deterministic necessary-and-sufficient condition (which remains the spectral condition of Theorem E.2(3) applied to $w^P$), and it does not cover all $\binom{K}{W}$ windows simultaneously without paying for the larger union. The constants $c,C$ are universal but not optimized, and no claim is made that (E.8) is tight in its $W$- or $K_\infty$-dependence --- Hanson--Wright and matrix-Bernstein refinements can improve window-family dependence for structured $\mathcal{A}$, but they do not change the deterministic criterion. Finally, all of Appendix~E concerns the carrier statistics of the \emph{design}; the translation of whitening into gain-estimation rates ($1/W$ per window) and into image error (leverage $B_L=1$ for full SRHT inversion) is carried out in Appendix~D and Appendix~F respectively.

\subsection{E.6 Two randomization semantics: pixel randomization (\texorpdfstring{$P_{\mathrm{col}}$}{Pcol}, \texorpdfstring{$D$}{D}) vs.\ acquisition-order randomization (\texorpdfstring{$P_{\mathrm{row}}$}{Prow})}

\textbf{This subsection separates two operations that earlier drafts conflated under the single word ``permutation''.} Write $P_{\mathrm{col}}$ for a uniformly random \textbf{pixel (column) permutation} of the object --- the $P$ of assumption (E3) --- write $D$ for the i.i.d.\ Rademacher \textbf{pixel sign} diagonal of (E3), and write $P_{\mathrm{row}}$ for a uniformly random reordering of the \textbf{acquisition (Hadamard row / time) order}. These control different claims:

\begin{itemize}
\item \textbf{Appendix E's theorems govern cross-row covariance.} By Lemma E.1, the sign-draw covariance $\operatorname{Cov}_D(Z_g,Z_h)=\widehat{w^P}(g+h)$ is a function of the permuted squared object $w^P$ alone; it is therefore \textbf{invariant to the acquisition order} --- $P_{\mathrm{row}}$ does not enter Lemma E.1 or Theorems E.2--E.3 at all. The whitening randomization of the theory is $x=UDP_{\mathrm{col}}T$: pixel signs plus pixel permutation, in \emph{natural} acquisition order.
\item \textbf{The manuscript's temporal-stationarity screen is controlled by acquisition-order structure.} The Brown--Forsythe/Levene screen~\cite{BrownForsythe1974,Levene1960} of Sec.~7 probes temporal homogeneity of the acquired time series; it responds to whatever breaks the coarse-to-fine sequency-energy ordering of natural Hadamard rows --- which either $P_{\mathrm{col}}$ \textbf{or} $P_{\mathrm{row}}$ accomplishes.
\end{itemize}

\paragraph{Factorized ablation (empirical attribution; \texttt{results/perm\_ablation\_r1/}).} All arms share the same ordered natural-Hadamard baseline ($K=4096$, $64\times64$, 8192 paired frames, noiseless carrier), the same 10-object panel, and the same Levene protocol (8 chunks, reject at $p<10^{-3}$); 32 draws per randomized arm. Levene rejection rates:

\begin{center}
\small
\setlength{\tabcolsep}{4.5pt}%
\begin{tabular}{lcccc}
\toprule
arm & $P_{\mathrm{col}}$ & $D$ & $P_{\mathrm{row}}$ & rate \\
\midrule
ordered (baseline) & -- & -- & -- & \textbf{0.700} \\
$P_{\mathrm{col}}$ only (theory's $P$) & Y & -- & -- & \textbf{0.056} \\
$P_{\mathrm{row}}$ only & -- & -- & Y & \textbf{0.122} \\
$D$ only & -- & Y & -- & \textbf{0.434} \\
$P_{\mathrm{col}}+D$ (theory) & Y & Y & -- & \textbf{0.119} \\
$P_{\mathrm{row}}+D$ (code's \texttt{srht\_paired}) & -- & Y & Y & \textbf{0.059} \\
$P_{\mathrm{col}}+P_{\mathrm{row}}$ (old Exp.-B arm) & Y & -- & Y & \textbf{0.059} \\
\bottomrule
\end{tabular}
\end{center}

The $P_{\mathrm{col}}+P_{\mathrm{row}}$ arm reproduces the historical headline ($19/320=5.9\%$) exactly, confirming the seed-compatible reconstruction of the old experiment; that historical 5.9\% is attributable to \emph{either} permutation alone ($P_{\mathrm{col}}$: 5.6\%; $P_{\mathrm{row}}$: 12.2\%).

\paragraph{Code caveat (stated plainly).} The implemented basis \texttt{make\_srht\_paired\_basis} (\texttt{src/basis.py}, \texttt{permute\_rows=True}) applies pixel signs $D$ but permutes \textbf{Hadamard rows} (acquisition order), i.e.\ it implements $P_{\mathrm{row}}+D$ --- \textbf{not} the theorem's $P_{\mathrm{col}}+D$. Both combinations are empirically sufficient for the temporal-stationarity screen (5.9\% and 11.9\%, versus 70\% ordered), while \textbf{$D$ alone is not} (43.4\%): sign flips whiten carrier \emph{amplitudes} (chunk-std CV collapses from 0.76 to 0.03) but leave the temporal sequency ordering intact, so the chunk-variance test still rejects. Text and captions must therefore attach each claim to its own arm: the covariance-whitening claims of E.1--E.3 attach to $P_{\mathrm{col}}$ (+$D$); the temporal-screen results attach to the arm actually acquired ($P_{\mathrm{row}}+D$ for the code's SRHT).

\paragraph{What this does not claim.} The table is an empirical attribution on one object panel and one protocol, not a theorem; no claim is made that $P_{\mathrm{row}}$ enjoys the concentration guarantees (E.2)--(E.4), which are proved for $P_{\mathrm{col}}$. The small increase from adding $D$ on top of $P_{\mathrm{col}}$ ($5.6\%\to11.9\%$) is a property of the variance-homogeneity test on near-white sequences, not a whitening failure. Nothing here alters the theorems of E.1--E.5; it fixes which experimental arm may cite them.

\section{The master finite-noise identity and the flip boundary}

This appendix proves Theorem~1 (the master finite-noise relMSE identity) and Prop~3 (the finite-$N$ flip boundary) of Sec.~8, together with the noise plug-ins and the three basis specializations quoted there. Throughout, $T\in\mathbb{R}^K$ is the object, with $S_1$, $S_2$, $K_{\mathrm{eff}}$ as in Appendix~C; $A\in\mathbb{R}^{N\times K}$ is the design, $b=AT$ the ideal bucket vector; the raw bucket is $R_n=a_nB_n$ with gain $a_n=e^{\ell_n}$, $\ell\in S$ (the drift class of Sec.~4, $p=\dim S$). After the pipeline applies a gain estimate $\hat a_n=e^{\hat\ell_n}$, the corrected bucket is
\begin{equation}
\begin{gathered}
z_n \;=\; R_n/\hat a_n \;=\; (1+\delta_n)\,b_n+\xi_n,\\
\delta_n=\frac{a_n}{\hat a_n}-1=e^{\ell_n-\hat\ell_n}-1,
\end{gathered}
\tag{F.1}
\end{equation}
where $\delta_n$ is the \emph{residual} multiplicative gain error and $\xi_n$ is additive bucket-domain noise after correction, with $\mathbb{E}\,\delta=m_\delta$, $\mathrm{Cov}(\delta)=V_\delta$, $\mathbb{E}\,\xi=0$, $\mathrm{Cov}(\xi)=\Sigma_\xi$, and $\delta\perp\xi$ (uncorrelated, $C_{\delta\xi}=\operatorname{Cov}(\delta,\xi)=0$ --- a hypothesis, revisited in the cross-term note below). The reconstruction is $\hat T=Lz$ for any fixed linear operator $L\in\mathbb{R}^{K\times N}$ (exact inverse, pseudoinverse, DGI correlator, or a regularized estimator linearized around a fixed design).

\subsection{F.1 The master identity}

\begin{namedthm}{Theorem 1 (exact second-moment bridge)}
Assume model (F.1) with $A,T,L$ fixed (all expectations conditional on them), and assume --- as part of the hypothesis, not as background convention --- that:

(H1) $\mathbb{E}\,\xi=0$;

(H2) $\delta$ and $\xi$ are uncorrelated, $C_{\delta\xi}:=\operatorname{Cov}(\delta,\xi)=0$;

(H3) the gain estimate $\hat a$ is \textbf{exogenous} with respect to the reconstruction record --- computed independently of the noise realization $\xi$ (injected/synthetic residuals, an independent calibration record, or a cross-fitted estimate from held-out frames) --- so that (H1)--(H2) hold conditionally on the record being reconstructed;

(H4) the second moments $m_\delta,V_\delta,\Sigma_\xi$ are finite.

Then, exactly and with no Gaussian or small-noise approximation,
\begin{align*}
\frac{\mathbb{E}\|\hat T-T\|_2^2}{S_2}
=\;&\underbrace{\frac{\|(LA-I)T+L\,\mathrm{diag}(b)\,m_\delta\|_2^2}{S_2}}_{\text{bias}}\\
&+\underbrace{\frac{\mathrm{tr}\!\big(L\,\mathrm{diag}(b)\,V_\delta\,\mathrm{diag}(b)\,L^\top\big)}{S_2}}_{\text{residual gain}}\\
&+\underbrace{\frac{\mathrm{tr}\!\big(L\,\Sigma_\xi L^\top\big)}{S_2}}_{\text{additive noise}}. \tag{F.2}
\end{align*}

(a) If moreover $\mathbb{E}\,\delta_n=0$ and the $\delta_n$ are uncorrelated with common variance $v$, the residual-gain term equals $v\,B_L(A,T)$ with the \textbf{leverage}
\begin{equation}
B_L(A,T)=\frac{1}{S_2}\sum_{n=1}^N b_n^2\,\|Le_n\|_2^2. \tag{F.3}
\end{equation}

(b) If instead the residual gains are stationary with $\mathrm{Cov}(\delta_n,\delta_m)=v\,r_\delta(n-m)$, the residual-gain term equals
\begin{equation}
R_{\mathrm{gain}}=\frac{v}{S_2}\sum_{n,m} r_\delta(n-m)\,b_n b_m\,\langle Le_n,Le_m\rangle. \tag{F.4}
\end{equation}
\end{namedthm}

\noindent\textbf{Proof.} Write $z=b+\mathrm{diag}(b)\,\delta+\xi$, so $\hat T-T=(LA-I)T+L\,\mathrm{diag}(b)\,\delta+L\xi$. Decompose $\delta=m_\delta+(\delta-m_\delta)$. The mean of $\hat T-T$ is $(LA-I)T+L\,\mathrm{diag}(b)\,m_\delta$ (using (H1)); its squared norm is the bias term. By the bias--variance decomposition $\mathbb{E}\|X\|^2=\|\mathbb{E}X\|^2+\mathrm{tr}\,\mathrm{Cov}(X)$ and (H2), the covariance of $\hat T$ is $L\,\mathrm{diag}(b)\,V_\delta\,\mathrm{diag}(b)\,L^\top+L\,\Sigma_\xi L^\top$; taking traces and dividing by $S_2$ gives (F.2). For (a), $V_\delta=vI$ and $\mathrm{tr}(L\,\mathrm{diag}(b)\,\mathrm{diag}(b)\,L^\top)=\sum_n b_n^2\|Le_n\|^2$, giving (F.3). Case (b) is the coordinate expansion of the same trace with $(V_\delta)_{nm}=v\,r_\delta(n-m)$. \hfill$\blacksquare$

\paragraph{Remark F.1.1 (same-record gain estimates violate the hypotheses; scope of the exact identity).} Hypothesis (H3) is not decorative. When $\hat a$ is estimated from the \emph{same} noisy record that carries $\xi$ --- the generic situation for a blind pipeline --- the conditional mean $\mathbb{E}[\xi\mid\text{record}]$ is generically nonzero and $\delta$ is correlated with $\xi$, so (H1)--(H2) fail: the mean of $\hat T-T$ acquires an additional bias term $L\,m_\xi$ (with $m_\xi=\mathbb{E}\,\xi\neq0$) inside the bias norm, and the covariance acquires the cross-term $(2/S_2)\,\operatorname{tr}\!\big[L\,\mathrm{diag}(b)\,C_{\delta\xi}\,L^\top\big]$. Neither term is included in (F.2). The exact identity therefore covers the \textbf{exogenous / injected-residual case} --- which is precisely what the injected-residual protocol of Sec.~9 (experiment E4, Fig.~10) validates --- and cross-fitted estimates that restore (H1)--(H2) by construction; for same-record blind pipelines, (F.2) is exact only up to the two extra terms above.

\paragraph{What this does not claim.} With (H1)--(H3) now stated \emph{in} the theorem, the caveat is Remark F.1.1: (F.2) is not claimed for same-record blind gain estimation, where the $L\,m_\xi$ bias and the $C_{\delta\xi}$ cross-term appear. Beyond that, (F.2) is a second-moment accounting identity for a \emph{linear} reconstructor with the design held fixed; it does not bound nonlinear or data-adaptive $L$, and for random-pattern pipelines the constants below are obtained by additionally averaging over the pattern draw. It says nothing about how small $v$ can be made --- that is the estimation-rate question of Sec.~6. Finally, the scalar collapse $v\,B_L$ \emph{requires} uncorrelated residuals: for smooth (coherent) residual gain --- the generic output of a windowed blind estimator --- only the matrix form (F.4) is correct. In particular, the simple summary ``orthogonal inversion gives $v$'' holds only for independent coefficient-wise residuals under exact inversion.

\subsection{F.2 Noise plug-ins}

\begin{namedthm}{Proposition F.2 (read and Poisson noise, exact)}
\emph{(a) Gaussian read.} If detector-domain noise $\eta_n\sim\mathcal N(0,\sigma_{\mathrm{read},n}^2)$ (independent across frames) precedes gain correction, then $\xi_n=\eta_n/\hat a_n$ and, exactly, $\mathrm{Var}(\xi_n)=\sigma_{\mathrm{read},n}^2/\hat a_n^2=\big(\sigma_{\mathrm{read},n}^2/a_n^2\big)\big(a_n/\hat a_n\big)^2$; for calibrated gains $\hat a_n\approx a_n$ the factor $(a_n/\hat a_n)^2\to1$ and $\mathrm{Var}(\xi_n)=\tau_{G,n}^2:=\sigma_{\mathrm{read},n}^2/a_n^2$. \emph{(b) Poisson.} If counts obey $C_n\mid T,a\sim\mathrm{Pois}(\Phi_n)$ with $\Phi_n=\gamma_n a_n b_n$ and $z_n=C_n/(\gamma_n\hat a_n)$, then exactly
\begin{equation}
\mathbb{E}[z_n\mid T,a]=\frac{a_n}{\hat a_n}b_n,\qquad \mathrm{Var}(z_n\mid T,a)=\frac{b_n^2}{\Phi_n}\Big(\frac{a_n}{\hat a_n}\Big)^2. \tag{F.5}
\end{equation}
\end{namedthm}

\noindent\textbf{Proof.} (a) is scaling of a Gaussian. (b): a Poisson variable has mean $=$ variance $=\Phi_n$; dividing by $\gamma_n\hat a_n$ scales the mean by $(\gamma_n\hat a_n)^{-1}$ and the variance by $(\gamma_n\hat a_n)^{-2}$, and $\Phi_n/(\gamma_n\hat a_n)= (a_n/\hat a_n) b_n$, $\Phi_n/(\gamma_n\hat a_n)^2=(b_n^2/\Phi_n)(a_n/\hat a_n)^2$. \hfill$\blacksquare$

Hence --- and only after the post-calibration simplification $\hat a_n\to a_n$, which sets the exact $(a_n/\hat a_n)^2$ factors of (a) and (F.5) to one --- $\Sigma_\xi=\mathrm{diag}(\tau_{G,n}^2+b_n^2/\Phi_n)$ plugs into (F.2), and the Poisson term is \textbf{exact at all photon counts}, including $\Phi_n\ll1$; no log or Gaussian approximation enters the reconstruction layer. Under a total budget $\Phi_n=\omega_n\Phi_{\mathrm{tot}}$, $R_{\mathrm{Pois}}=(S_2\Phi_{\mathrm{tot}})^{-1}\sum_n(b_n^2/\omega_n)\|Le_n\|^2$.

\paragraph{What this does not claim.} Exactness holds \emph{conditionally on the gains}; the difficulty of low photons lives entirely in estimating $\hat a_n$ (the calibrated soft-log analysis; see Appendix~D), not in propagating noise through $L$. The optimal allocation $\omega_n\propto|b_n|\,\|Le_n\|$ requires knowing $T$ and is not claimed to be achievable.

\subsection{F.3 Specializations}

\paragraph{F.3.1 Orthogonal / full-SRHT inversion ($B_L=1$).} If $A=U$ is orthonormal ($N=K$, signed coefficients available) and $L=U^\top$, then $LA=I$ and $\|Le_n\|=1$, so $B_L=(1/S_2)\sum_n b_n^2=1$ by Parseval. With independent residuals,
\begin{equation}
\mathrm{relMSE}_{\mathrm{orth}}=v+\frac{1}{S_2}\sum_n\big(\tau_{G,n}^2+b_n^2/\Phi_n\big). \tag{F.6}
\end{equation}
Exact orthogonal inversion transmits independent coefficient-wise residual gain at unit leverage. A full SRHT acquisition $A=P_{\mathrm{row}}UDP_{\mathrm{col}}$, $L=A^\top$ has identically the same propagation; SRHT's role is \emph{not} a better inverse but temporal stationarization of the carrier so that the blind estimator can shrink $v$ (Sec.~5/7).

\paragraph{F.3.2 Pairwise Hadamard: exact perturbation law.} Let $c_k=h_k^\top T$, $h_k\in\{\pm1\}^K$, with physical masks $B_k^\pm=(S_1\pm c_k)/2$ and reconstruction $T=K^{-1}H^\top c$, so coefficient errors $e_k$ give $\mathrm{relMSE}=(KS_2)^{-1}\sum_k\mathbb{E}e_k^2$. If the paired masks carry gains $a_k^+$ and $a_k^-=a_k^+(1+\Delta_k)$, the pairwise-normalized estimator satisfies the \textbf{exact} algebraic law
\begin{equation}
\hat c_k-c_k=-\,\frac{S_1\,\Delta_k\,(1-x_k^2)}{2+\Delta_k(1-x_k)},\qquad x_k=c_k/S_1, \tag{F.7}
\end{equation}
whence, exactly (Parseval),
\begin{equation}
R_{\mathrm{pair,gain}}
=\frac{K_{\mathrm{eff}}}{K}\sum_{k=1}^K
\mathbb{E}\!\left[\frac{\Delta_k^2(1-x_k^2)^2}{\{2+\Delta_k(1-x_k)\}^2}\right].
\tag{F.8a}
\end{equation}
\emph{Second-moment expansion (bounded increments).} If $|\Delta_k|\le\eta<1$ almost surely, then for $g_x(t)=\{2+t(1-x)\}^{-2}$ the mean value theorem with $|x|\le1$, $|t|\le\eta$ gives $|g_x(t)-g_x(0)|\le|t|/\{2(1-\eta)^3\}$, whence
\begin{equation}
\begin{aligned}
&\left|R_{\mathrm{pair,gain}}
-\frac{K_{\mathrm{eff}}}{4K}\sum_{k=1}^K(1-x_k^2)^2\,\mathbb{E}\Delta_k^2\right|\\
&\qquad\le\frac{K_{\mathrm{eff}}}{2(1-\eta)^3K}\sum_{k=1}^K\mathbb{E}|\Delta_k|^3 .
\end{aligned}
\tag{F.8b}
\end{equation}
If the pair increments share a common marginal distribution (no cross-pair independence required),
\begin{equation}
\begin{aligned}
R_{\mathrm{pair,gain}}
&=\frac{K_{\mathrm{eff}}}{4}\,D_H(T)\,\mathbb{E}\Delta^2
+O_\eta\!\big(K_{\mathrm{eff}}\,\mathbb{E}|\Delta|^3\big),\\
D_H(T)&=\frac{1}{K}\sum_k(1-x_k^2)^2 .
\end{aligned}
\tag{F.8}
\end{equation}
\textbf{The second moment is $\mathbb{E}\Delta^2$, not $\operatorname{Var}(\Delta)$}: the two may be interchanged only under the \emph{extra} assumption $\mathbb{E}\Delta=0$. A deterministic mismatch $\Delta\equiv\epsilon$ is the immediate counterexample --- its $\operatorname{Var}(\Delta)=0$, yet (F.8a) gives a strictly positive risk of order $K_{\mathrm{eff}}\epsilon^2$. For most non-DC Hadamard coefficients of a nonnegative object $|x_k|\ll1$, so $D_H\approx1$; $D_H$ is nevertheless a \emph{pipeline constant to be measured}, not assumed.

\emph{OU increment corollary (exact lognormal moments).} For a stationary Gaussian OU/AR(1) log gain with $\operatorname{Var}(\ell_n)=s^2$ and adjacent correlation $e^{-\rho}$, put $r(\rho)=1-e^{-\rho}$ and $u=s^2r(\rho)$. The adjacent log increment $X=\ell_{n+1}-\ell_n$ is $\mathcal N(0,2u)$ and $\Delta=e^X-1$, so
\begin{equation}
\mathbb{E}\Delta=e^{u}-1,\qquad
\mathbb{E}\Delta^2=e^{4u}-2e^{u}+1=2u+7u^2+O(u^3),
\tag{F.8c}
\end{equation}
with $\operatorname{Var}(\Delta)=e^{4u}-e^{2u}$ and $\mathbb{E}\Delta^2-\operatorname{Var}(\Delta)=(e^{u}-1)^2$. The bounded-increment proposition does not apply directly to this lognormal $\Delta$; instead, split (F.8a) on $\{|X|\le c\}$: on the central event Taylor's theorem gives $\Delta^2(1-x^2)^2/\{2+\Delta(1-x)\}^2=\tfrac14(1-x^2)^2X^2+O_c(|X|^3)$ uniformly in $x\in[-1,1]$, while on the complement the exact fraction admits an exponential envelope $C_c(1+e^{2|X|})$ whose Gaussian-tail expectation is $o(u^m)$ for every $m$. Since $\mathbb{E}X^2=2u$ and $\mathbb{E}|X|^3=O(u^{3/2})$,
\begin{equation}
R_{\mathrm{pair,gain}}
=\tfrac12K_{\mathrm{eff}}D_H(T)\,s^2r(\rho)
+O\!\big(K_{\mathrm{eff}}\,u^{3/2}\big):
\tag{F.8d}
\end{equation}
the leading pair law survives, but its proof must begin from $\mathbb{E}\Delta^2$. The $K_{\mathrm{eff}}$ amplification is structural: pairwise normalization converts a differential gain error into a coefficient error of order $S_1$, and Parseval then divides by $S_2$. Noise terms: read noise gives exactly $R_{\mathrm{pair,read}}=2\sigma_{\mathrm{read}}^2/S_2$; for Poisson counts with pair budget $\Phi_{\mathrm{pair}}$, using $B_k^++B_k^-=S_1$, the coefficient shot variance is $\mathrm{Var}(e_k\mid c_k)\approx S_1^2/\Phi_{\mathrm{pair},k}$ --- an approximation, exact only up to the $c_k$-dependent split between the two buckets --- giving $R_{\mathrm{pair,Pois}}=K_{\mathrm{eff}}/\Phi_{\mathrm{pair}}$ (equal-pair budget) or $(K_{\mathrm{eff}}+1)/(2\Phi_{\mathrm{mask}})$ (equal-mask budget); these differ only by allocation convention.

\paragraph{What F.3.2 does not claim.} (F.7)--(F.8) presuppose a stable DC/$S_1$ normalization; at low photons the denominator bucket sum can be noisy or zero, and a robust implementation needs offset/clipping or a likelihood-ratio estimator. The exact law (F.7) is algebra; only the moment step (F.8) invokes small $\Delta$ and index-independence.

\paragraph{F.3.3 Random patterns / DGI: floor, leverage, and population-moment constants.} \textbf{Three distinct constants must carry three distinct names; no generic ``pipeline $C_0$'' label is permitted.} For a \emph{fixed} design $A$, linear reconstruction $L$, object $T$, and $b=AT$, define the \emph{raw fixed-design sampling floor} and the \emph{residual leverage}
\begin{align}
C_0^{\mathrm{floor}}(A,L,T)&:=\frac{N}{S_2}\,\|(LA-I)T\|_2^2, \tag{F.9a}\\
\Lambda_{\mathrm{lev}}(A,L,T)&:=N\,B_L(A,T)=\frac{N}{S_2}\sum_{n=1}^Nb_n^2\|Le_n\|_2^2, \tag{F.9b}\\
C_0^{\mathrm{lev}}(A,L,T)&:=\Lambda_{\mathrm{lev}}-1. \tag{F.9c}
\end{align}
The first is the raw-MSE sampling floor of the fixed operator; the second propagates exogenous multiplicative residuals and equal-count Poisson noise ($C_0^{\mathrm{lev}}$ is merely the shifted notation $\Lambda_{\mathrm{lev}}-1$). Third, for $N$ iid pattern draws $I_n=\mu\mathbf 1+\sigma x_n$ with iid coordinates ($\mathbb{E}x=0$, $\mathbb{E}x^2=1$, $\mathbb{E}x^3=\gamma_3$, $\mathbb{E}x^4=\beta_4$), the ideal known-$(\mu,\sigma)$ correlator $\hat T_{\mathrm{DGI}}=(N\sigma)^{-1}\sum_nx_nz_n$ and one-sample vector $Z=\sigma^{-1}x(\mu S_1+\sigma x^\top T)$ define the \emph{population raw-moment constant}
\begin{equation}
\begin{aligned}
\frac{\mathbb{E}\|\hat T_{\mathrm{DGI}}-T\|^2}{S_2}&=\frac{C_0^{\mathrm{raw,moment}}}{N},\\
C_0^{\mathrm{raw,moment}}&:=\frac{\mathbb{E}\|Z-T\|^2}{S_2},
\end{aligned}
\tag{F.9}
\end{equation}
whose evaluation is a moment calculation: expanding $\mathbb{E}\|Z\|^2$ with $Z=(\mu S_1/\sigma)x+xx^\top T$ and subtracting $\|\mathbb{E}Z\|^2=S_2$,
\begin{align*}
C_0^{\mathrm{raw,moment}}&=K+\beta_4-2+K_{\mathrm{eff}}\big[K(\mu/\sigma)^2+2\gamma_3(\mu/\sigma)\big]\\
&\;\xrightarrow[\ \mu=0,\ \gamma_3=0\ ]{}\; K+\beta_4-2. \tag{F.10}
\end{align*}
The $K_{\mathrm{eff}}K(\mu/\sigma)^2$ term is the familiar nonnegative-background penalty. For the ideal correlator, independence and unbiasedness give $\mathbb{E}_AC_0^{\mathrm{floor}}=C_0^{\mathrm{raw,moment}}$ and $\mathbb{E}_A\Lambda_{\mathrm{lev}}=1+C_0^{\mathrm{raw,moment}}$; \textbf{for a mean-subtracted or scale-aligned implementation these are not definition-level identities} --- each constant must be measured for the pipeline that is actually run. On the published grid the mean-subtracted pipeline's \emph{raw} fixed-design floor sits near $10^3$ (\texttt{results/prop3\_raw\_metric\_r1/}: $C_0^{\mathrm{floor,raw,fix}}=980.110$--$1114.529$, consistent with the centered-moment value $K+\beta_4-2=1023.800$ to $0.8\%$), its \emph{scale-aligned} floor near $6.2$--$7.2\times10^2$, and its nonnegative-background drift leverage near $10^6$; conflating any two of these is a factor-$1.5$ to factor-$10^3$ error.

\begin{namedthm}{Proposition (exact pipeline specialization)}
Suppose the residuals are exogenous to detector noise and design, with common mean $m_\delta$, common second moment $q_\delta=\mathbb{E}\delta_n^2$, and $\operatorname{Cov}(\delta)=(q_\delta-m_\delta^2)I$. Under the hypotheses of Theorem~1,
\begin{align}
\frac{\mathbb{E}\|\hat T-T\|_2^2}{S_2}
={}&\frac{\|(LA-I)T+m_\delta LAT\|_2^2}{S_2}\notag\\
&+\frac{\Lambda_{\mathrm{lev}}}{N}(q_\delta-m_\delta^2)+R_\xi,
\tag{F.10a}
\end{align}
with $R_\xi=\operatorname{tr}(L\Sigma_\xi L^\top)/S_2$. If $m_\delta=0$,
\begin{equation}
\frac{\mathbb{E}\|\hat T-T\|_2^2}{S_2}
=\frac{C_0^{\mathrm{floor}}}{N}
+\frac{1+C_0^{\mathrm{lev}}}{N}\,q_\delta+R_\xi ;
\tag{F.10b}
\end{equation}
and if $\Sigma_\xi=\tau_G^2I+\Phi_{\mathrm{frame}}^{-1}\operatorname{diag}(b^2)$,
\begin{equation}
R_\xi=\frac{\tau_G^2\operatorname{tr}(LL^\top)}{S_2}
+\frac{1+C_0^{\mathrm{lev}}}{N\,\Phi_{\mathrm{frame}}}.
\tag{F.10c}
\end{equation}
\end{namedthm}

\noindent\textbf{Proof.} The residual mean vector is $m_\delta\mathbf1$ and $L\operatorname{diag}(b)\mathbf1=Lb=LAT$; substitute this mean and the stated covariance into the exact identity (F.2). The Poisson specialization uses $(S_2\Phi_{\mathrm{frame}})^{-1}\sum_nb_n^2\|Le_n\|_2^2=\Lambda_{\mathrm{lev}}/(N\Phi_{\mathrm{frame}})$. \hfill$\blacksquare$

For the ideal iid correlator with centered residuals the pattern-averaged bookkeeping is the familiar display
\begin{align*}
\mathrm{relMSE}_{\mathrm{DGI}}=\;&\frac{C_0^{\mathrm{floor}}}{N}+\frac{(1+C_0^{\mathrm{lev}})\,v}{N}\\
&+\frac{K\tau_G^2}{N\sigma^2S_2}+\frac{1+C_0^{\mathrm{lev}}}{N\,\Phi_{\mathrm{frame}}}, \tag{F.11}
\end{align*}
where for the ideal correlator $\mathbb{E}_A\Lambda_{\mathrm{lev}}=1+C_0^{\mathrm{raw,moment}}$, so in residual-injection experiments the measurable slope is $NB_L\to\Lambda_{\mathrm{lev}}$. \textbf{Status (labeled).} This pattern-averaging step treats the residuals $\delta$ as independent of the pattern draw --- exact for injected residuals (the Fig.~10 protocol), a heuristic decoupling for blind pipelines, where $\hat a$ is computed from the same record the patterns generated. Correlated residuals multiply $v$ by $\chi_\delta=1+2\sum_{h\ge1}\mathrm{corr}(\delta_0,\delta_h)$ --- a heuristic long-run-variance correction valid only when correlations are weak enough that pattern cross-terms still average; the rigorous object remains (F.4).

\paragraph{Correction (drift-leverage constant of the gain term; measured, no fitted parameters).} As previously written, the gain term of (F.11) evaluated with the \emph{scale-aligned floor} constant in place of the leverage is \textbf{wrong for mean-subtracted DGI pipelines} under time-structured (drift-like) residual gain: on the published grid it underpredicts the measured random-arm drift excess by a median factor $\approx1.4\times10^3$ (1364 / 1527 / 1521 across the none / AGC / SCGI-proxy corrections; white-drift diagnostic, \texttt{results/prop3\_nofreeparam\_r1/}). The correct leverage is the \textbf{nonnegative-background constant} $\Lambda_{\mathrm{lev}}=1+C_0^{\mathrm{lev,raw\ background}}$, obtained from (F.10) evaluated at the measured raw pattern moments (at the grid's operating point $C_0^{\mathrm{lev,raw\ background}}\approx0.7$--$1.2\times10^6$, versus the scale-aligned floor $\approx6.2$--$7.2\times10^2$ and the \emph{raw} fixed-design floor $\approx1.0$--$1.1\times10^3$); with it the measured excess is reproduced to 0.3--4.7\% (median ratios 1.003 / 1.047 / 1.039), still with zero fitted parameters. Mechanism: mean subtraction removes the $K_{\mathrm{eff}}K(\mu/\sigma)^2$ background penalty from the \textbf{sampling floor} --- hence the small floor constant in the $C_0^{\mathrm{floor}}/N$ term, confirmed exactly by the flat oracle arm --- but \textbf{not} from the \textbf{drift leverage}: a time-varying gain modulates the large DC background coherently, and the correlator pays the full background leverage. For mean-subtracted pipelines the gain term of (F.11) must therefore carry $\Lambda_{\mathrm{lev}}=1+C_0^{\mathrm{lev,raw\ background}}$, while the sampling-floor term keeps the (raw or aligned, \emph{consistently chosen}) floor constant. This vindicates --- and extends to the gain term --- the ``no universal $C_0$'' warning below, which earlier versions applied to the floor term only.

\paragraph{What F.3.3 does not claim.} There is \textbf{no universal $C_0$}: (F.10) is exact only for the ideal known-$(\mu,\sigma)$ correlator with iid coordinates. Mean subtraction, reference normalization, and photon allocation all change the constants (not the $C_0/N$ structure), so each of the three constants of (F.9a)--(F.9c)/(F.9) must be measured from the pipeline that is run. The sampling-floor constant and the drift-leverage constant of the \emph{same} pipeline are different measured numbers whenever the pipeline subtracts a background --- quoting one for the other is a $\sim10^3$ error at nonnegative-background operating points --- and the \emph{raw} and \emph{scale-aligned} floors of the same pipeline differ by a further factor $\approx1.54$ (the metric-mismatch audit of F.4 below). Table~I's residual-probe value $\Lambda_{\mathrm{lev}}=NB_L\in[7.38\times10^5,7.68\times10^5]$ is a Monte Carlo mean over the bridge experiment's object-by-seed residual probes, \emph{not} a deterministic statistic and \emph{not} the Prop.~3 floor constant. The flux-limited convention $\Phi_n=\gamma b_n$ replaces the last term of (F.11) by $(N\sigma^2S_2)^{-1}\mathbb{E}[\|x\|^2 b_n/\gamma]$.

\subsection{F.4 Prop 3: the finite-\texorpdfstring{$N$}{N} flip boundary}

\paragraph{Drift model and assumptions.} Let $s^2$ be the variance scale of the log-gain process $\ell$ and let $\rho_{\mathrm{pair}}$ be the \textbf{adjacent-pair decorrelation parameter}: for an Ornstein--Uhlenbeck log-gain the paired-mask gain increment obeys $\mathbb{E}[\Delta_k^2]=2s^2r(\rho_{\mathrm{pair}})+O(s^4r^2)$ with $r(\rho)=1-e^{-\rho}$ (exact moments: (F.8c); $r(\rho)=\rho+O(\rho^2)$ at small $\rho$). Assume the small-mismatch regime of (F.8) (so $R_{\mathrm{pair,gain}}(\rho_{\mathrm{pair}})\approx\tfrac12 K_{\mathrm{eff}}D_H(T)\,s^2\,r(\rho_{\mathrm{pair}})$, per (F.8d)) and the DGI decomposition (F.11) with blind residual variance $v_{\mathrm{blind}}(\rho_{\mathrm{pair}},N)$, treated here as an \textbf{externally measured function of $(\rho_{\mathrm{pair}},N)$} (the runner measures it on the grid). \textbf{No deterministic-Lipschitz $\rho^{2/3}$ rate is attributed to the OU drift}: OU paths are not Lipschitz, and the correct interior OU window analysis (lemma below, (F.14)--(F.16)) gives stochastic drift error $s^2\rho W/6$ to first order and window-optimized risk $\propto\rho^{1/2}$, not $\rho^{2/3}$.

\begin{namedthm}{Prop 3 (raw-MSE flip boundary; fixed-point form)}
\textbf{All risks in this proposition are raw relative MSEs of Theorem~1} (the metric caution below explains why a scale-aligned analogue is a different number). Define $\rho^*=\inf\{\rho_{\mathrm{pair}}\ge0: R_{\mathrm{pair}}(\rho_{\mathrm{pair}})\ge R_{\mathrm{rand}}(\rho_{\mathrm{pair}},N)\}$, where $R_{\mathrm{pair}}=R_{\mathrm{pair,gain}}+R_{\mathrm{pair,read}}+R_{\mathrm{pair,Pois}}$ and, for the random arm,
\begin{equation}
R_{\mathrm{rand}}(\rho,N)
=\frac{C_0^{\mathrm{floor}}}{N}
+\Gamma_{\mathrm{blind}}(\rho,N)
+R_{\mathrm{rand,noise}} . \tag{F.12a}
\end{equation}
For an \emph{oracle-known-gain} arm, $\Gamma_{\mathrm{blind}}=0$ (the validated skeleton below). For an \emph{exogenous, centered, frame-uncorrelated} residual with $q_{\mathrm{blind}}(\rho,N)=\mathbb{E}\delta_n^2$,
\begin{equation}
\Gamma_{\mathrm{blind}}(\rho,N)
=\frac{1+C_0^{\mathrm{lev}}}{N}\,q_{\mathrm{blind}}(\rho,N),
\tag{F.12b}
\end{equation}
with $C_0^{\mathrm{lev}}$ the pipeline's drift-leverage constant of (F.9c) (equal to $C_0^{\mathrm{lev,raw\ background}}\approx10^6$ for the mean-subtracted pipeline --- the F.3.3 correction). \textbf{No scalar representation is asserted for a same-record blind estimator}: there the full covariance, nonzero-mean, and residual--noise cross terms of Theorem~1 are required, and any scalar display must be labeled a conditional proxy. At first order define
\[
A_T=\tfrac12K_{\mathrm{eff}}D_H(T)s^2,
\qquad
Q(\rho)=\frac{\Delta_R(\rho)}{A_T},
\]
\[
\Delta_R(\rho)=\frac{C_0^{\mathrm{floor}}}{N}+\Gamma_{\mathrm{blind}}(\rho,N)
+R_{\mathrm{rand,noise}}-R_{\mathrm{pair,noise}} .
\]
Assume $A_T>0$ (equivalently $s>0$ and $D_H(T)>0$) and that the risks, hence $Q$, are continuous in $\rho$ on the comparison interval. Then the leading boundary is the \textbf{fixed-point equation}
\begin{equation}
1-e^{-\rho^*}=Q(\rho^*), \tag{F.12}
\end{equation}
up to the remainder in (F.8d). If $\Delta_R(0)\le0$ then $\rho^*=0$ (the flip is immediate, not absent). If $\Delta_R(0)>0$, a finite crossing on $[0,\rho_{\max}]$ \emph{exists} when $1-e^{-\rho}\ge Q(\rho)$ somewhere on that interval, and is \emph{unique} only under an additional single-crossing condition (e.g.\ strict increase of $1-e^{-\rho}-Q(\rho)$); if the difference stays negative, no crossing occurs in range. Only if $Q(\rho)\equiv Q_0$ is \emph{constant} may one write the explicit solution: $Q_0\le0\Rightarrow\rho^*=0$; $0<Q_0<1\Rightarrow\rho^*=-\log(1-Q_0)$; $Q_0\ge1\Rightarrow$ no finite crossing in the leading OU model. When $q_{\mathrm{blind}}$ depends on $\rho$, the logarithmic display is a fixed point, \textbf{not an explicit solution}.
\end{namedthm}

\noindent\textbf{Proof.} At the boundary, $R_{\mathrm{pair}}(\rho^*)=R_{\mathrm{rand}}(\rho^*,N)$. Substituting $R_{\mathrm{pair,gain}}=\tfrac12K_{\mathrm{eff}}D_Hs^2\,r(\rho^*)$ (from (F.8d), whose proof starts from $\mathbb{E}\Delta^2$, \emph{not} $\operatorname{Var}\Delta$) and (F.12a)--(F.12b), and solving for $r(\rho^*)=1-e^{-\rho^*}$, gives (F.12). Existence under the stated sign/continuity conditions is the intermediate value theorem applied to $\rho\mapsto1-e^{-\rho}-Q(\rho)$; the infimum is a crossing by continuity; uniqueness is exactly the single-crossing assumption and is not otherwise proven. \hfill$\blacksquare$

\begin{statusblock}\textbf{Empirical status of Prop 3 (two regimes; corrected constant \emph{and} corrected metric; \texttt{results/prop3\_nofreeparam\_r1/}, \texttt{results/prop3\_raw\_metric\_r1/}).} The published equal-frame grid ($N=2048$, $K=1024$, OU drift, 10 objects, all constants measured, zero fitted parameters) splits the proposition's status in two:

\emph{(a) The $v=0$ skeleton is VALIDATED, no-free-parameter --- in the raw metric, at the few-percent level.} For a gain-known random arm ($v=0$: oracle gains, so only the floor competes with the pair term), comparing \textbf{raw prediction with raw observation} --- the raw fixed-design floor $C_0^{\mathrm{floor,raw,fix}}=980.110$--$1114.529$ in $Q_{\mathrm{raw}}=2(C_0^{\mathrm{floor,raw,fix}}/N)/(K_{\mathrm{eff}}D_H\sigma_a^2)$, $\rho^*_{\mathrm{pred}}=-\log(1-Q_{\mathrm{raw}})$, against raw pair-versus-raw oracle crossings on the same OU paths --- gives, for the 40 cells with $\sigma_a\ge0.1$: \textbf{40/40 observed raw crossings, median factor 1.04387, maximum factor 1.19985} (median signed $\log_{10}$ ratio $-0.0172$; all within factor 2). The earlier aligned-metric analysis reported a median factor $1.54$: that factor is a \textbf{raw/aligned metric mismatch} --- the raw/aligned floor ratio on the same objects is $1.5398517$ --- \emph{not} evidence for the (F.8) cubic remainder, and the former ``consistent with the $O(\Delta^3)$ correction'' explanation is deleted. The two $\sigma_a=0.05$ aligned crossings become $Q_{\mathrm{raw}}\ge1$ with no raw crossing; they were already outside the claimed small-mismatch validation range. The empirical $\sigma$-scaling exponent is $-2.11$ ($R^2=0.976$) versus $-2.07$ predicted --- the leading-order $\sigma^{-2}$ law is correct. The residual discrepancy grows toward $\approx20\%$ at $\sigma_a=0.5$; only that residual may plausibly be associated with higher-order small-mismatch terms, and even that reading is offered cautiously. Ingredients check out separately: the pair law (F.8) is essentially exact (median measured/predicted $0.997$ over 639 cells), and the oracle random arm sits flat at its floor across the whole grid.

\emph{(b) The blind-arm flip prediction, as previously stated, is WITHDRAWN --- the corrected theory predicts no flip, and none occurs.} With the erroneous pipeline-$C_0$ leverage, (F.12) predicted flips at $\rho^*=0.007$--$0.33$ for $\sigma_a\ge0.1$; empirically ordered pairwise Hadamard stays strictly better than blind random+DGI in \textbf{all 100} (source $\times$ correction $\times$ $\sigma$ $\times$ object) cells over the entire scanned range $\rho\in[10^{-3},20]$. With the corrected leverage $C_0^{\mathrm{lev}}=C_0^{\mathrm{lev,raw\ background}}$, the bracket satisfies $Q(\rho)>1-e^{-\rho}$ for every $\rho$ at every grid ($\sigma_a$, object) for both blind corrections --- the fixed-point equation (F.12) then has no solution in range, i.e.\ the theory predicts \textbf{no flip in range}, exactly matching the 100/100 empirical no-flip. The failure of the original prediction was localized to one constant (the F.3.3 leverage correction, verified to 0.3--4.7\% at white drift), not to the functional form of (F.12); the boundary curves previously shown in Fig.~4 are SRHT-paired-versus-ordered-Hadamard crossovers --- a comparison Prop 3 explicitly does not rank --- summarized by free two-parameter power-law fits, and are not a test of (F.12).\end{statusblock}

\paragraph{Leading-order heuristic (labeled as such).} If $r(\rho)=\rho$, $D_H=1$, noise differences are negligible, $v_{\mathrm{blind}}$ is negligible, $Q$ is treated as constant, and $\rho$ \emph{means} adjacent-pair decorrelation, (F.12) collapses to the engineering rule
\begin{equation}
\rho^*\approx\frac{2C_0^{\mathrm{floor}}}{N\,K_{\mathrm{eff}}\,s^2}, \tag{F.13}
\end{equation}
with first corrections $\rho^*\approx 2\big(C_0^{\mathrm{floor}}+(1+C_0^{\mathrm{lev}})v_{\mathrm{blind}}\big)/[NK_{\mathrm{eff}}D_Hs^2]+2(R_{\mathrm{DGI,noise}}-R_{\mathrm{pair,noise}})/(K_{\mathrm{eff}}D_Hs^2)$ and the nonlinear replacement $\rho\mapsto r^{-1}(\cdot)$ when adjacent increments are not infinitesimal. Note that for mean-subtracted pipelines ``negligible $v_{\mathrm{blind}}$'' means $(1+C_0^{\mathrm{lev,raw\ background}})v_{\mathrm{blind}}\ll C_0^{\mathrm{floor}}$ --- a far stronger requirement than $v_{\mathrm{blind}}\ll C_0^{\mathrm{floor}}$, and the reason the blind-arm flip vanished on the published grid (empirical-status point (b) above). (F.13) is a \emph{leading-order heuristic}, not a theorem; the theorem is (F.12).

\paragraph{Lemma (OU window risk; not a Lipschitz-path theorem).} Let $\ell_n$ be stationary Gaussian with $\operatorname{Cov}(\ell_n,\ell_{n+h})=s^2e^{-\rho|h|}$, and let $\varepsilon_n$ be an independent centered stationary carrier-noise sequence with autocovariance $\gamma_\varepsilon(h)$. For an interior centered window $W=2m+1$ and $\hat\ell_n=W^{-1}\sum_{j=-m}^m(\ell_{n+j}+\varepsilon_{n+j})$,
\begin{align}
\mathbb E(\hat\ell_n-\ell_n)^2
={}&s^2\Big[1-\tfrac2W\sum_{j=-m}^me^{-\rho|j|}\notag\\
&\qquad+\tfrac1{W^2}\sum_{i,j=-m}^me^{-\rho|i-j|}\Big]\notag\\
&+\tfrac1{W^2}\sum_{i,j=-m}^m\gamma_\varepsilon(i-j).
\tag{F.14}
\end{align}
If $\rho W\to0$, the OU contribution is
\begin{equation}
s^2\rho\,\frac{W^2-1}{6W}+O(s^2\rho^2W^2), \tag{F.15}
\end{equation}
and if $\sum_h|\gamma_\varepsilon(h)|\le\sigma_{\mathrm{abs}}^2$,
\begin{equation}
\mathbb E(\hat\ell_n-\ell_n)^2\le\frac{\sigma_{\mathrm{abs}}^2}{W}+Cs^2\rho W .
\tag{F.16}
\end{equation}
An interior optimizer therefore has $W^*\asymp\sigma_{\mathrm{abs}}/(s\sqrt\rho)$ and $\mathrm{MSE}^*\asymp\sigma_{\mathrm{abs}}\,s\sqrt\rho$ (for iid carrier noise, $W^*\sim\sqrt6\,\sigma_\varepsilon/(s\sqrt\rho)$ and $\mathrm{MSE}^*\sim\sqrt{2/3}\,\sigma_\varepsilon s\sqrt\rho$). \textbf{The deterministic Lipschitz $\rho^{2/3}$ rate is therefore not an OU output.} Prop~3 may treat $q_{\mathrm{blind}}(\rho,N)$ as an externally measured function --- as the empirical analysis does --- but any OU \emph{estimator theorem} must use this stochastic $\rho^{1/2}$ analysis instead.

\paragraph{Metric and implementation cautions (raw versus scale-aligned).} The Prop.-3 runner's scale-aligned gain error $\min_c\|c\hat a-a\|^2/\|a\|^2$ is not $q_\delta=N^{-1}\sum_n(a_n/\hat a_n-1)^2$ and is not $\operatorname{Var}(\delta)$; they agree only to first order after a fixed small-error gauge. Likewise, a scale-aligned reconstruction floor is not the raw $C_0^{\mathrm{floor}}$ required by Prop.~3. If $\hat T=T+e$, put $a=\langle e,T\rangle/S_2$ and $b=\|e\|^2/S_2$; the implementation's scale-aligned error, minimizing over $c\in\mathbb R$, is
\begin{equation}
R_{\mathrm{align}}=\frac{b-a^2}{1+2a+b}, \tag{F.17}
\end{equation}
which is not generally proportional to the raw risk $b$ (even for $e\perp T$ it is $b/(1+b)$, not $b$). Therefore a raw-MSE theorem and an aligned-MSE empirical boundary must not be called an exact no-free-parameter validation of one another; the raw re-export of \texttt{results/prop3\_raw\_metric\_r1/} exists precisely to compare raw with raw.

\paragraph{$\rho$-convention warning.} Throughout, $\rho_{\mathrm{pair}}$ (adjacent-pair decorrelation) is the variable of (F.12)--(F.13); it is \textbf{not} the normalized bandwidth $\rho_{\mathrm{bw}}$ of the tall-design thresholds ($p\approx\rho_{\mathrm{bw}}N$, Sec.~4). If a bandwidth parameter with adjacent-increment variance $\sim2s^2\rho_{\mathrm{bw}}/N$ is substituted, the explicit $1/N$ in (F.13) cancels and the boundary reads $\rho^*_{\mathrm{bw}}\approx2C_0/(K_{\mathrm{eff}}s^2)$ --- same physics, different symbol. Every phase-diagram axis must declare its convention.

\paragraph{What Prop 3 does not claim.} The boundary compares two \emph{specific} pipelines (pairwise-normalized Hadamard vs.\ blind random+DGI with the ideal correlator); it does not rank SRHT-paired pipelines, externally calibrated systems, or prior-regularized reconstructions, each of which shifts $\Delta_R$. It inherits every assumption of its ingredients: small mismatch (F.8b)--(F.8d), measured constants $C_0^{\mathrm{floor}}$, $C_0^{\mathrm{lev}}$, and $D_H$ (with the floor/leverage distinction of the F.3.3 correction \emph{and} the raw/aligned metric distinction of (F.17)), a stationary drift with a declared $r(\rho)$, and a $q_{\mathrm{blind}}$ that is an externally measured function, not an OU theorem (the $\rho^{2/3}$ attribution is withdrawn; the OU lemma gives $\rho^{1/2}$). Every prediction and every observation must be stated in the \emph{same} raw or aligned metric --- the strict no-free-parameter claim attaches only to the raw-vs-raw comparison of \texttt{results/prop3\_raw\_metric\_r1/}; a purely aligned comparison is a scale-aligned empirical analogue. It is a mean-square crossing, not a guarantee about any single realization; the single-crossing property itself is assumed, not derived; and at $\rho_{\mathrm{pair}}$ beyond the small-drift expansion only the defining infimum --- not (F.12) --- is meaningful. Finally, the validated regime is the $v=0$ skeleton (empirical-status point (a)); no flip prediction is currently claimed for the blind arm --- the corrected constant predicts, and the data confirm, no flip on the published grid.
